% =====================================================================
%  Harald Høffding, BEGREBET ANALOGI
%  Det Kgl. Danske Videnskabernes Selskab. Filosofiske Meddelelser, I, 4.
%  København: Hovedkommissionær Andr. Fred. Høst & Søn, Kgl. Hof-Boghandel.
%  Bianco Lunos Bogtrykkeri, 1923.  138 pp.
%  DIPLOMATIC TRANSCRIPTION -- GENERATED; DO NOT EDIT BY HAND.
%  Made by ebook-method/begrebet-analogi/build.py --emit from this template, the
%  decisions (decisions.tsv) and the patches (patches.tsv) there.
%
%  WITNESSES.  (1) The SAGA e-book (Begrebet_analogi.epub, ISBN 9788726363067),
%  an OCR of the 1923 print: base text, paragraphs, italics, small caps, notes.
%  (2) The Royal Danish Academy's scan of the print, read by tesseract
%  (ebook-method/collate/ocr_layer.py): every letter, the page turns, the
%  footnote placement.  The Academy's first PDF was damaged and lacked pp.
%  135-138; an intact copy (2026-10-03) supplies them (see pagemap.py; the
%  old repair, repair_scan.py, is now only historical).  Every disagreement
%  no rule settles was read at the image (decisions.tsv); the Greek, which
%  the e-book garbles, was read at the image and set by patches.tsv.  Print
%  governs; misprints are kept and marked PRINTED AS IS.
%
%  The INDHOLD of p. 138, which the e-book lacks, was set from the scan and
%  read at the image.  Letterspacing (the print's emphasis, set \emph) was
%  found by measuring the letter gaps in the scan and confirmed at the image;
%  short or line-broken spaced words may have been missed.
%
%  CONVENTIONS.  "% --- p. N ---" + \opage{N} at the first word of printed page N
%  (a word broken over the turn is split with %).  Footnotes as printed, numbered
%  per page (\footnote[k]).  Small capitals \textsc; italics \textit;
%  letterspacing \emph.  Running heads, folios and signature marks are not
%  transcribed.
% =====================================================================
\documentclass[12pt]{article}
\usepackage[utf8]{inputenc}
\usepackage[T1]{fontenc}
\usepackage{textalpha}
\usepackage{libertinus}
\usepackage{libertinust1math}
\usepackage[danish]{babel}
\usepackage[protrusion=true,expansion=true]{microtype}
\usepackage{setspace}
\usepackage[margin=1.3in]{geometry}
\usepackage{fancyhdr}
\usepackage{hyperref}
\hypersetup{pdftitle={Begrebet Analogi}, pdfauthor={Harald Høffding}, hidelinks}
\pagestyle{fancy}
\fancyhf{}
\fancyhead[LE,RO]{\thepage}
\fancyhead[C]{\textit{Harald Høffding: Begrebet Analogi (1923)}}
\setstretch{1.3}
\usepackage{marginnote}
\usepackage{graphicx}
\DeclareUnicodeCharacter{0254}{\reflectbox{c}}   % ɔ (id est)
\DeclareUnicodeCharacter{2192}{\ensuremath{\rightarrow}}
\DeclareUnicodeCharacter{221A}{\ensuremath{\surd}}
\newcommand{\tocline}[3]{\noindent\makebox[2.2em][r]{#1}\hspace{0.6em}#2\dotfill\ #3\par}
\newcommand{\opage}[1]{\marginnote{\footnotesize\textit{[#1]}}}

\begin{document}

% --- front matter: title (PDF 3) ---
\begin{center}
Det Kgl. Danske Videnskabernes Selskab.\\
Filosofiske Meddelelser, I, 4.\\[5em]
{\LARGE BEGREBET ANALOGI}\\[1.5em]
AF\\[1.5em]
{\large HARALD HØFFDING}\\[12em]
KØBENHAVN\\[0.5em]
{\small HOVEDKOMMISSIONÆR: ANDR. FRED. HØST \& SØN, KGL. HOF-BOGHANDEL}\\
{\small BIANCO LUNOS BOGTRYKKERI}\\[0.5em]
1923
\end{center}
\clearpage

% ===== TEXT (pp. 3--137) =====
% --- p. 3 ---
\setcounter{footnote}{0}%
\opage{3}

\begin{center}\large Indledning.\end{center}

1. I denne Afhandling har jeg sat mig til Opgave at fortsætte den nærmere
Undersøgelse af nogle af de Grundbegreber, jeg i min Bog »Den menneskelige
Tanke« (1910) forsøgte at paavise som de mest karakteristiske for menneskelig
Erkendelse, saaledes som vi kender denne efter dens hidtidige Udvikling. Jeg
har altsaa villet give et Sidestykke til Afhandlingen om »Totalitet som
Kategori« (Videnskabernes Selskabs Skrifter. 1917) og til Afhandlingen om
»Relation som Kategori« (Videnskabernes Selskabs filosofiske Meddelelser.
1920), og det er Begrebet Analogi, jeg her nærmere vil undersøge.

Foreløbig kan Analogi defineres som Forholdslighed mellem to Emner, altsaa en
Lighed, der ikke gælder enkelte Egenskaber ved disse Emner eller enkelte Dele
af dem, men Forholdet mellem Egenskaber indbyrdes eller Dele indbyrdes. Der kan
være meget stor Forskel mellem to Emners Egenskaber eller Dele, hver for sig,
og dog kan der være en Lighed i Forholdet mellem Egenskaber eller Dele. Analogi
er altsaa en særegen Art af Lighed. Den staar i en vis Forstand fjernere fra
Identitet (absolut Lighed) end Kvalitetslighed, sammensat Lighed og
Dækningslighed, som kun angaar enkelte Egenskaber eller Dele, og som lettere
danner Overgangsformer til Identitet, (smlgn. »Den menneskelige Tanke« § 21).
Dog maa bemærkes, at der ikke blot gives en Emneidentitet, men ogsaa en Forholds%
% --- p. 4 ---
\setcounter{footnote}{0}%
\opage{4}identitet, og allerede Aristoteles skelner mellem kvantitativ Analogi
(ἰσότης λόγων), Proportionalitet i streng Betydning, og kvalitativ Analogi, der
f. Eks. kan gælde mellem geometriske Figurer eller mellem Organer hos
forskellige levende Væsener. Og der gives desuden en Forskelsidentitet, som
ligger til Grund ved Dannelsen af en identisk varierende Forskelsrække (Den
menneskelige Tanke § 65) og særlig ved Dannelsen af de for exakt Videnskab
grundlæggende identiske Kvalitetsrækker (Tal, Tid, Grad, Sted). (Den
menneskelige Tanke § 76---82).

Baade Totalitet og Relation har jeg givet Plads mellem Kategorierne, d. v. s.
mellem de Grundbegreber, de Tankeformer, der paa Videnskabens nuværende Stadium
spiller en væsentlig Rolle i Tankearbejdet, bestemmer dets Opgaver og dets hele
Retning (Den menneskelige Tanke § 51---52; 57). Hverken Aristoteles eller Kant
har givet Analogien Plads blandt Kategorierne, og i »Den menneskelige Tanke«
har jeg fulgt dem heri. Dog førtes jeg i det nævnte Skrift til at lægge megen
Vægt paa, at i Modsætning til rent formal Tænken (Logik og Aritmetik) spiller
ved al Tænken over reale Emner Analogi en væsentlig Rolle. Modsætningen mellem
Identitet (Emneidentitet i Logiken, Forskels- og Forholdsidentitet i
Aritmetiken) og Analogi er derfor af stor Betydning for den rette Opfattelse af
vor Erkendelses Natur og Gyldighed. Men naar Analogi virkelig spiller saa stor
en Rolle, medens Identitet egentlig er et uopnaaeligt Ideal udenfor
Abstraktionernes formale Verden, kunde det være berettiget at give den Plads i
Kategoriernes Række, og i det Skema over Kategorier, der findes paa sidste Side
i Afhandlingen »Totalitet som Kategori«, er dette ogsaa sket. Den opføres der
blandt de formale Kategorier, lige efter Begrebet Identitet, fordi den ligger til
% --- p. 5 ---
\setcounter{footnote}{0}%
\opage{5}Grund ved den Reduktion af Kvalitetsforhold, der --- ved Hjælp af
Identitetssynspunkter --- finder Sted ved den strengere Udvikling af Begreberne
Tal, Tid, Grad og Sted. Saadanne Rækker danner Mellemled mellem absolute
Identitetsrækker og kvalitativt forskellige Emnerækker. De dannes ved
analogiserende Abstraktion fra disse sidste. Det er her, Analogien især viser
sin Betydning. Det skulde nu være nærværende Afhandlings Opgave at følge
Begrebet Analogi videre i dets speciellere Betydning for menneskelig Erkendelse.

Jeg dvæler først ved de uvilkaarlige Analogier som gør sig gældende hos
primitive Mennesker og hos Børn, men ogsaa, paa Aandslivets Højdepunkter, hos
Geniet. De fremtræder ikke som Analogier for vedkommende Bevidsthed selv; dette
er først muligt for den kritiske Bevidsthed, der har lært at skelne mellem
Analogi og Identitet og at indse, at enhver begrundende og bevisende
Tankeovergang sker i Kraft af Identitet. Derefter gaar jeg over til at
undersøge Analogi fra et formalt logisk Synspunkt, og i de følgende Afsnit
søger jeg at paavise Analogier mellem Erkendelsens forskellige Funktioner og
mellem dens forskellige Omraader.

Foreløbig bemærkes, at Analogiens Betydning for Erkendelsen viser sig i tre
forskellige Henseender.

Størst Betydning har den sikkert ved at virke spejdende, foregribende og
ansporende. Kategorier lægger sig overhovedet for Dagen i de Spørgsmaal, der
stilles (Den menneskelige Tanke § 52). Analogi kan være den faktiske Vej til
sand Erkendelse, selv om Analogiens Gyldighed bagefter forkastes. Man kan jo
uddrage en rigtig Slutning af falske Forudsætninger. Naturligvis kan man ikke
bevise, at Slutningen er rigtig, uden ved at finde de rette Forud%
% --- p. 6 ---
\setcounter{footnote}{0}%
\opage{6}sætninger, men det er jo ofte, at en Antagelse staar for os som sand,
og at den netop derved stiller os den Opgave at udfinde, hvorfor den i Grunden
er sand. Forudsætningerne (de rette Forudsætninger) er ofte det sidste, der
findes; men man vilde ikke søge efter dem, naar en Analogi ikke først havde
ført os til en Antagelse og derved stillet os en Opgave. Analytisk Metode beror
paa, at der er stillet os en Opgave, og at man saa søger Betingelserne for at
løse denne Opgave. Findes saa de rette Forudsætninger, forkaster man maaske den
Analogi, der ledte paa Vej, ikke blot hvis den er falsk, men ogsaa hvis den er
unødvendig. Grunden til, at en Analogi opstilles, vil ligge i, at der
foreligger Brud paa en kontinuerlig Sammenhæng, som uvilkaarlig har dannet sig
for os, og som en uvilkaarlig Trang fører til at fastholde saa længe som
muligt. Analogien kan her føre til et intellektuelt Vovestykke, der undertiden
krones med Held. Man kan anvende paa Analogien, hvad Pascal siger om Fantasien
idethele, at den bedrager saa meget des lettere, som den ikke altid bedrager
(d'autant plus fourbe qu'elle ne l'est pas toujours). Dersom den ikke bedrager,
eller dersom den viser sig at kunne undværes, afløses den af en
Forholdsbestemmelse, der tilsidst begrundes ved Identitet.

Men Analogi er ikke blot forberedende. Den kan være det sidste Foreningsbaand
mellem forskellige Emner eller Emnerækker, der ikke kan reduceres saaledes, at
et enkelt Emne eller en enkelt Emnerække ses at ligge til Grund, eller
saaledes, at de alle ses at have deres Grund i dybere liggende Emner eller
Emnerækker. En vis Afrunding og Overskuen bliver overfor saadanne Emner kun
mulig ved Analogiens Hjælp. En Forklaring af, hvorfor der gives forskellige
Emner eller Emnerækker, faar vi ikke derved, men
% --- p. 7 ---
\setcounter{footnote}{0}%
\opage{7}vi konstaterer det faktiske Forhold, nemlig det, at vi paa visse
Punkter ikke naar videre end til en Forbindelse, en Syntese, der kun bygger paa
den fjerneste Grad af Lighed, som Analogi er. Det vil vise sig, at vi baade
overfor de forskellige Erkendelsesfunktioner (se Kap. III) og overfor de
forskellige Erkendelsesomraader (se Kap. IV) tilsidst er stillede paa denne
Maade. Vi ender da i aabne Spørgsmaal, idet det staar hen, om en videre
Reduktion, en dybere liggende Kontinuitet skulde kunne findes, eller om der her
i Tingenes Natur skulde være Forskelligheder, som i hvert Tilfælde for os er
irreduktible. Spørgsmaal kan altsaa ikke blot være Visdommens Begyndelse, men
ogsaa dens Ende.

Det vil have sin Interesse tilsidst at dvæle noget ved Analogier, hvis
Betydning ligger i at være Udtryk for poetiske og religiøse Stemninger, Udtryk,
der snart kan minde om det uvilkaarlige Sjælelivs Analogier, snart om den
videnskabelige Tænkens Søgen og Finden. Især vil Anelser eller Forhaabninger om
en Enhed og en Kontinuitet, der ligger dybere, end Erfaring kan paavise, finde
Udtryk i denne Art Analogier.

\begin{center}\large I. Uvilkaarlige Analogier.\end{center}

2. Som allerede bemærket forudsætter Analogi ikke altid, at der udtrykkelig er
skelnet mellem forskellige Emner, henholdsvis deres Egenskaber eller Dele, og
at der saa bagefter, i en vis Modsætning til Forskellighederne bringes en
Sammenhæng tilveje mellem Emnerne paa Grundlag af en Forholdslighed. I
uvilkaarlige sjælelige Fænomener, som Sanseanskuelse, Erindring og Fantasi, kan
det ikke altid paavises, at Elementerne først er tilstede som sondrede og først
bagefter indgaar Forbindelse. Disse uvilkaarlige sjæle%
% --- p. 8 ---
\setcounter{footnote}{0}%
\opage{8}lige Processers Betydning ligger netop i, at de tilvejebringer
Helheder, Kombinationer, der først bagefter kan analyseres. Eftertanken behøver
stadig et Stof, men dette Stof behøver ikke, som man ofte har troet, at bestaa
i en Mængde selvstændige Elementer; det kan være Helheder, som det saa er
Eftertankens Opgave at opløse i Elementer, der ikke uden videre tør betragtes
som identiske med de Elementer, de uvilkaarlig dannede Helheder skyldes. Nye
opdukkende Elementer kan nemlig straks og uvilkaarlig sammenvæves med
Bevidsthedens givne Indhold, saa at der maaske slet ikke mærkes nogen
Forandring. Alligevel er der en Tilbøjelighed hos Eftertanken til at betragte
de Elementer, dens Analyse finder, som Helhedens oprindelige, elementære
Grundlag. Der er her en Analogi tilstede, der idethøjeste kan søge Berettigelse
i den Kontinuitet, hvormed Overgangen kan ske fra et uvilkaarlig dannet
sjæleligt Kontinuum til den Differentiering, hvorom allerede selve Eftertankens
Opstaaen vidner. Egentlig kommer, naar vi taler om de mest primitive og
uvilkaarlige sjælelige Processer, alle Betegnelser tilkort, og det er intet
Under, da Sprogets psykologiske Betegnelser er hentede ud fra det klart
differentierende, kritisk reflekterende Sjæleliv. Selve Udtrykket Analogi
bruger vi egentlig her paa analogiserende Vis. Det gaar her, ligesom naar man
taler om »ubevidste Slutninger«. Disse er ikke Slutninger i egentlig Forstand,
men Processer, Overgange, der kan fortolkes i Analogi med de egentlige
Slutninger. Det Samme gælder fremdeles (som jeg allerede har vist i min
Psykologi og senere i Den menneskelige Tanke § 5), naar vi opstiller Begrebet
Syntese som Udtryk for alt Bevidsthedslivs Natur. I Erkendelsesteorien, hvor vi
undersøger Mulighederne for videnskabelig Erkendelse, maa vi ogsaa drøfte den
Mulighed, at det Givne kunde bestaa af
% --- p. 9 ---
\setcounter{footnote}{0}%
\opage{9}en Hob ligestore og indbyrdes irreduktible Forskelligheder. Kun en
Tilnærmelse kan paavises til en saadan kaotisk Forskelsrække (ligesom til dens
Modstykke, en absolut Identitetsrække); der vilde i hvert Tilfælde behøves en
indgaaende og anstrengt Eftertanke til at godtgøre, at en saadan Række virkelig
forelaa i et bestemt givet Tilfælde.

I den Maade, paa hvilken en enkelt Oplevelse lige fra sin Undfangelsesstund
forbindes med andre samtidige eller fortidige Oplevelser eller med
Overleveringer, i hvilke Individet har levet sig ind, saa at de ogsaa kan
betragtes som Oplevelser, foreligger der allerede en Tydning. Jo mere vi kommer
tilbage til, hvad vi kalder det primitive Menneskes Stade, des nærmere rykker
Oplevelse og Tydning hinanden, og der er tilsidst ingen Modsætning mellem dem
at finde. Naar vi dog forsøger at holde dem ude fra hinanden, gør vi Brug af en
Analogi med, hvad der først gør sig gældende for den prøvende Eftertanke. Og,
som sagt, selve Begrebet Analogi bruges analogisk, naar vi siger, at primitive
Mennesker tyder enhver ny Oplevelse i Analogi med, hvad de kender fra de
Overleveringer, de har levet sig ind i. --- I min Afhandling »Oplevelse og
Tydning« har jeg søgt at beskrive forskellige Maader, paa hvilke Forholdet
mellem Oplevelse og Tydning fremtræder historisk, hvad ekstatiske Oplevelser
angaar. ---

I sit betydningsfulde Værk »Les fonctions mentales dans les sociétés
inférieures« (1910) har \textsc{Lévy-Bruhl} kritiseret den tidligere
almindelige Opfattelse, at det primitive Menneske opfatter alt i Lighed med sig
selv, antropomorfiserer eller analogiserer. Det er efter hans Opfattelse ikke
Sammenligning eller Eftertanke, der her virker, men en aldeles umiddelbar
Identificeren. Hos det primitive Menneske ligger der stedse Overleveringer til
Grund, med hvilke
% --- p. 10 ---
\setcounter{footnote}{0}%
\opage{10}det Nye, der opleves, uvilkaarlig smelter sammen. De fra Stammen
overleverede Forestillinger gør sig gældende med en saadan Følelsesstyrke, at
der ingen Trang til Eftertanke opstaar. Det, der i høj Grad afviger fra det
Overleverede, lægges der foreløbig ikke Mærke til; først senere udtænkes der
maaske en med Overleveringen stemmende Tydning. Naar den Vilde ser et Billede
af en Genstand, han kender, bringer Billedet som Billede ham intet Nyt, thi
Billedet er selve Genstanden, hvorledes han saa end nærmere lægger sig dette
Forhold tilrette. Og med al den logiske Konsekvens, der allerede udmærker dette
Stade, som Lévy-Bruhl med Urette har kaldt det »prælogiske«, frygter den Vilde
derfor for at blive afhængig af den, der har et Billede af ham, og han frygter
for at komme til at lide Mangel paa Vildt, naar en Europæer har tegnet eller
fotograferet mange vilde Dyr paa Egnen og er rejst bort med Billederne. Og som
det gaar med Billeder, gaar det ogsaa med Navne og med Ceremonier: de er, hvad
de betyder. Billeder, Navne og Ceremonier er vel kun en Del af et Hele, der har
andre Egenskaber foruden dem, der afbildes eller benævnes; men denne Del
\emph{er} det Hele. Lévy-Bruhl finder i denne Identificeren af Del og Helhed,
--- hvad han kalder »loi de participation«%
% footnote 1 on p. 10
\footnote[1]{Til denne svarer, hvad \textsc{Grønbech} \textit{(Primitiv
Religion,} p. 8) kalder »Helhedsrealisme«, og som er karakteristisk for det
primitive Sjæleliv.}, --- den mest karakteristiske Ejendommelighed ved det
primitive Sjæleliv. Det er sikkert, at den Tydning, det primitive Menneske
anstiller under Indflydelse af en Overlevering, som omfattes med inderlig
Tilslutning og Ærefrygt, ikke kan følge den formelle Logiks Principer. Disse
gør sig kun gældende, naar Tvivl er vakt; men han er vant til at betragte
Overleveringen som et selvfølgeligt \emph{Grundlag} for sin Opfattelse af alt,
der sker. Tydnings%
% --- p. 11 ---
\setcounter{footnote}{0}%
\opage{11}grundlaget er ikke Værk af Eftertanke. Men det kan lede til
Misforstaaelse, naar Lévy-Bruhl udtaler (p. 425), at de kollektive
Forestillinger, der beherskes af »Participationens Lov«, er »ligegyldige
overfor Modsigelsens logiske Lov«. Der maa skelnes mellem det givne
Tydningsgrundlag (her: den i Ærefrygt fastholdte Overlevering, der nu engang
præges af »Participationsloven«) og selve den Tydning af enkelte Fænomener, der
sker ud fra dette Grundlag. Eftertanke virker til alle Tider og paa alle
Kulturtrin ud fra bestemte Forudsætninger, som ikke selv straks kan gøres til
Genstand for Eftertanke. Ud fra det Grundlag, det primitive Menneske nu engang
har, er det aldeles konsekvent, at han ikke forstaar, hvorfor Europæeren
skelner mellem Billedet og Tingen. Naar Lévy-Bruhl kalder »den primitive
Mentalitet« »prælogisk«, passer det paa Tydningsgrundlaget, ikke paa Tydningen
ud fra dette Grundlag.%
% footnote 1 on p. 11
\footnote[1]{Dr. \textsc{Albert Schweitzer,} der virker som Læge i
Ækvatorialafrikas Urskove, skriver, efter flere Aars Ophold blandt Negerne:
»Naturmennesket tænker dybere, end man i Almindelighed tror. Selv om han
hverken kan læse eller skrive, anstiller han dog Betragtninger i langt videre
Udstrækning, end vi aner.« (\textit{Mellem Floder og Urskov.} Dansk Overs.
1922. p. 134.)} Som der senere vil blive Lejlighed til at omtale, har
Lévy-Bruhl selv i en anden Sammenhæng hævdet, og med Rette, at menneskelig
Bevidsthed, hvad det sidste Grundlag for Domme og Slutninger angaar, egentlig
aldrig helt mister sin »prælogiske« Karakter. Der vil stedse være
Forudsætninger tilstede, som vi ikke drager og ikke kan drage frem, selv i vor
mest logiske Eftertanke. --- Hvad særlig den nøje Forbindelse mellem Navn og
Ting angaar, kan man endnu den Dag idag finde Exempler paa den. Der er
Tilbøjelighed til at anse den Terminologi, man selv bruger, for den eneste
»rigtige«, og det Sprog, man selv taler, synes at
% --- p. 12 ---
\setcounter{footnote}{0}%
\opage{12}være Et med Tingen, medens fremmede Sprogs Benævnelser ofte synes
fremmede for Tingene. Det fortælles, at en Tysker, der besøgte Paris, berettede
hjem, at man dér kaldte Brød for »pain«, og føjede til, »Wir nennen es Brot,
und es \emph{ist} doch auch Brot!«

Overleveringen træder for den Vilde frem som en sluttet Helhed, indenfor
hvilken han ingen Grund har til at foretage nogen Skelnen. Der raader foreløbig
en umiddelbar Tro paa Gyldigheden af overleverede Forestillinger, ligesom paa
de umiddelbart optrædende Sanseiagttagelser. Og Udtrykket »Tro paa« er endda
for reflekteret, thi der skelnes slet ikke mellem Troen og det, der tros paa.
Selve Begrebet Virkelighed eller Gyldighed dannes først, naar Tvivl har gjort
sig gældende. Den Grundsætning, der foreløbig tænkes efter (skønt den
naturligvis ikke formuleres), kan udtrykkes med Spørgsmaalet »Hvorfor ikke?«
Hvorfor ikke stole paa alt, hvad der med levende Magt fremtræder for os som
Sansning eller som Overlevering? Først senere lærer man at spørge »Hvorfor?«
--- og Eftertanken kan da endog naa til at rette dette Spørgsmaal mod selve
det, der hidtil har dannet det ubestridte Grundlag. Dertil kræves, at nye
Oplevelser optræder saa talrige og med en saadan fra alt Tilvant afvigende
Ejendommelighed, at de kan faa Magt over det, man allerede har indlevet sig i.
\textsc{Albert Schweitzer} udtaler som Resultat af flere Aars Samliv med Negere
i Centralafrika: »Man maa ikke tro, at man er kommet tilbunds i Negerens
Tankeverden, hvis man ikke er naaet længere end til at se paa hans overtroiske
Forestillinger og overleverede Retsregler. Inderst i ham lever en dunkel Anelse
om, at Opfattelsen af, hvad der er godt, maa komme ved Eftertanke« (Mellem
Floder og Urskove p. 136). Den opofrende Læge, der gennem historisk Kritik
% --- p. 13 ---
\setcounter{footnote}{0}%
\opage{13}havde lært at skelne mellem de kristelige Dogmer og den etiske
Idealisme, der kommer frem i Bjergprædikenen og i Parablerne, erfarede især
dette, naar han søgte at forklare Negerne sin religiøse Opfattelse. Dog
erfarede han ogsaa, at der ved Siden af den »Hjertets Moral«, der kunde
fremkaldes hos dem, længe holdt sig en Lyst til at lyve og stjæle og en Trang
til Hævn, der efter deres egne Overleveringer var deres Ret og maaske endog
deres Pligt.

Jeg maa her indflette den Bemærkning, at naar jeg i mine »Psykologiske
Undersøgelser« (p. 24 f.) siger, at i Forholdet mellem nye Oplevelser og
indøvede, tilvante Forestillinger ligger Hovedvægten paa det Nye, der »sætter
den hele Proces igang«, saa forudsætter dette et andet Udviklingstrin end den
uforstyrrede primitive Mentalitet, hvor det Nye saa let preller af overfor det
Tilvante og kun formaar at skaffe sig Magt i Sindet, naar det fremtræder paa
særlig indtrængende Maade. ---

Spørger man om Oprindelsen til det System af kollektive Forestillinger, der gør
sig gældende hos den Vilde og bestemmer hans Tydning af, hvad der møder ham,
svarer Lévy-Bruhl, at den »taber sig i Tidernes Nat«. Og da enhver
Menneskegruppe, man hidtil har truffet paa, møder med visse kollektive
Forestillinger, kan vi ikke paavise nogen absolut første Begyndelse. Vi maa
her, ligesom ved Spørgsmaalet om Sprogets Oprindelse%
% footnote 1 on p. 13
\footnote[1]{Smlgn. \textsc{Otto Jespersen:} \textit{Language, its Nature,
Development and Origin.}(1922). p. 416 f. om Metoden ved Drøftelsen af
Spørgsmaalet om Sprogets Oprindelse. Man maa (som allerede paa deres Vis Madvig
og Whitney havde hævdet) gaa ud fra den Maade, hvorpaa Sproget faktisk udvikler
sig, altsaa fra dets Historie, og derfra slutte tilbage til et mere primitivt
Sprog end noget af dem, vi kender.}, nøjes med Analogier, der hentes fra
Kulturtrin, som ligger os nærmest, og disse Analogier maa naturligvis bruges
med stor For%
% --- p. 14 ---
\setcounter{footnote}{0}%
\opage{14}sigtighed. Vi ser, at paa de Kulturtrin, vi kender, undergaar de
Overleveringer, der ligger til Grund ved hvert nyt Slægtleds aandelige
Udvikling, stadige Ændringer, Tilføjelser og Forskydninger, som for en
væsentlig Del netop skyldes selve det nye Slægtleds Erfaringer og Eftertanke.
Det Kollektive ændres stadig under Indflydelse af det Individuelle. Vi kan
ikke, som man ofte forsøgte i ældre Tider, tænke os en Tid, hvor der intet
Kollektivt var, eller hvor det Kollektive, der var, udelukkende skyldtes
enkelte Individers Erfaring og Eftertanke. Men ligesaa lidt kan vi tænke os en
Tid, hvor Sindene udelukkende beherskedes af fælles Overleveringer, uden at
individuelle Erfaringer og Tanker øvede nogensomhelst Indflydelse. Det maa
antages, at der paa alle Udviklingstrin er foregaaet en Vexelvirkning mellem
det Fælles og det Individuelle, og Trinenes Forskellighed kan kun bestaa i, at
paa tidligere Trin har det Fælles, paa senere Trin det Individuelle størst
Indflydelse, dog saaledes, at der atter og atter svinges frem og tilbage mellem
de forskellige Faktorers Overmagt. Selv paa de mest primitive Trin, vi kender,
kan individuelle Forskelligheder gøre sig gældende og Afvigelser fra Stammens
Tro fremtræde.%
% footnote 1 on p. 14
\footnote[1]{Se min Afhandling \textit{Sociologi og Religionsfilosofi} (i Bogen
om Oplevelse og Tydning. 1918) p. 131---138.} Og Exemplerne paa Sligt vilde
utvivlsomt være endnu talrigere, naar vi havde mere Kundskab om Livet indenfor
primitive Samfund, end vi har. Vi vilde have fundet mærkelige Afvigelser og
Initiativer, der nu drukner for os i de herskende Overleveringer.

De primitive Livsforhold frembyder stor Ensformighed og dermed stort Fællesskab
i de Erfaringer, de enkelte Individer kan gøre. Og paa Grund af det stadige
Samliv i Gruppe og Familie vil den Enkeltes særlige Erfaringer og
% --- p. 15 ---
\setcounter{footnote}{0}%
\opage{15}Opfindelser meget hurtig blive Fællesejendom, og hans Navn vil
glemmes. Dertil kommer, at det praktiske Krav stadig gør sig gældende; der
kræves Handling, og Livet leves for en stor Del som Refleks- og Instinktliv.
Enhver Forestilling faar sin væsentlige Betydning ved at udløse Bevægelse. Men
dertil er det ikke nødvendigt at gaa ind paa hele den Sammenhæng, til hvilken
Forestillingens Indhold hører. En enkelt Del eller en enkelt Egenskab er nok
til at udløse Handling, og man behøver ikke at skelne mellem Del og Helhed;
Delen staar uvilkaarlig som Tegn paa Helheden. Billedet, Navnet eller Tegnet
virker, som om selve Genstanden var given deri. Instinktet har nok i en
Antydning for at begynde at røre sig. Deri, at saaledes en Del kan have samme
Indflydelse som Helheden, og at det endog kan være gavnligt, at der ikke
anvendes Tid og Kraft paa at undersøge, om Helheden nu ogsaa virkelig er
repræsenteret ved den enkelte Del eller Egenskab, kunde ligge en Adgang til
Forstaaelse af »Participationsloven« eller »Helhedsrealismen«. Og naar en Del
saaledes atter og atter i Slægtens Erfaring viser sig at have samme praktiske
Betydning som Helheden, ser man Muligheden af, at der kan dannes en Kreds af
kollektive Forestillinger, som lyder Participationens Lov, og som det kan være
vanskeligt for nye Oplevelser at bryde.

Rent logisk kan denne Identificeren af Del og Helhed (pars pro toto) formuleres
som en Slutning i anden Figur: Delen udløser en vis Bevægelse, som Helheden
vilde udløse; altsaa er Del og Helhed det Samme! Og dette er netop en af de
Maader, paa hvilke en Analogislutning kan formuleres. Naar en Instinkthandling
kan lede vild, hænger det sammen med, at vedkommende »Slutning« er gal. Kun
indenfor aldeles bestemte Grænser, som den af ubetinget Handle%
% --- p. 16 ---
\setcounter{footnote}{0}%
\opage{16}trang beherskede Bevidsthed ikke giver sig Tid til at fastslaa, kan
det gaa an at identificere Del og Helhed. Hverken Instinkt eller prælogisk
Mentalitet kan derfor holde sig længere, end Analogislutningen praktisk slaar
til. Tvivlen vil vækkes, naar dette ikke længere er Tilfældet. En nærmere
psykologisk Undersøgelse vil kunne belyse dette Forhold nærmere og derved kaste
Lys over de uvilkaarlige Analogier.

3. Ligesom den voksne Europæer er for tilbøjelig til at lægge sit udviklede og
reflekterende Sjæleliv ind i de saakaldte primitive Menneskers Sjæleliv,
saaledes er der idethele en Tilbøjelighed til at lægge Funktioner, der
forudsætter en vis Udvikling af Eftertanken, ind i de mest umiddelbare og
uvilkaarlige sjælelige Processer. Man taler da om Forestillingsassociation og
om Dømmen, hvor det i Virkeligheden er mere elementære Processer, der foregaar.
Det kommer for en Del af, at Sprogets Udtryk for Sjæleliv væsentlig er hentede
fra det udviklede, differentierede Sjæleliv, og man er da ikke klar over, at
det kun ad Analogiens Vej kan være berettiget at bruge saadanne Udtryk om de
mest elementære psykiske Processer. Jeg har for en Del Aar siden (Psykologiske
Undersøgelser. Vid. Selsk. Skr. 1889) havt Lejlighed til at gøre opmærksom
herpaa og fremdrager nu nogle enkelte Punkter fra denne Undersøgelse, især for
at vise, at hvad Lévy-Bruhl har paavist for de primitive Menneskers Sjæleliv,
genfindes i sjælelige Processer, der den Dag idag foregaar i os alle.

a. Man har ikke været tilstrækkelig opmærksom paa, at der gives flere
forskellige Arter eller Former for Genkendelse. I den nævnte Afhandling (p. 35
f.) har jeg ment at maatte skelne mellem mindst ti Arter af Genkendelse.
% --- p. 17 ---
\setcounter{footnote}{0}%
\opage{17}Det fører til Forvirring, naar man, som det i Regelen sker, taler om
Genkendelse i al Almindelighed og saa dog faktisk holder sig til en eneste
Form, og det en Form, der forudsætter en vis differentieret Udvikling af
Forestillingslivet. Der foreligger da en Miskendelse af det elementære,
uvilkaarlige Sjælelivs Ejendommelighed.

Den mest elementære Form for Genkendelse, som kan paavises, er ligesaa
umiddelbar og uvilkaarlig, som en Sansefornemmelse kan være det. Et Emne staar
her med ét Slag for os som bekendt, med Bekendthedskvalitet, ligesom en
Fornemmelse har sin Sansekvalitet, uden at nogen Association, Sammenligning
eller Dom kan paavises. Et Navn, et Ansigtsudtryk, en Farvenuance kan optræde
med Bekendthedskvalitet, uden at vi formaar at erindre, at vi har iagttaget det
samme Emne før, og uden at vi kan henføre det til nogen bestemt Sammenhæng,
altsaa uden at vide, hvem der har det Navn o. s. v. Maaske giver
Bekendthedskvaliteten os Anledning til at danne den Dom, at vi har hørt Navnet
før o. s. v. Men denne Dom kommer saa netop bagefter. Og hvis man vil opfatte
saadan umiddelbar Genkendelse som Association, maa man formulere det som et
Grænsetilfælde og tænke sig Sagen saaledes, at naar en Syns- eller
Lydfornemmelse gentages, kan den i selve sit Fødselsøjeblik smelte sammen med
Eftervirkningen af tidligere Fornemmelser af samme Beskaffenhed (hvorledes man
saa vil tænke sig saadan Eftervirkning). Der sker altsaa en momentan
Reproduktion, og det reproducerede smelter strax sammen med den Fornemmelse,
som er ved at blive til. Bekendthedskvaliteten kan da teoretisk opfattes som
Produkt af to Elementer. Fysiologisk set foreligger der her et Exempel paa
Øvelse. Al Øvelse af en Funktion
% --- p. 18 ---
\setcounter{footnote}{0}%
\opage{18}skyldes en Sammensmeltning af ny Udløsning med Eftervirkninger af
tidligere Udløsninger.

Denne Paavisning af umiddelbar Genkendelse som elementært Fænomen, liggende paa
Grænsen mellem Fornemmelse og egentlig Erindring (jeg har kaldt den bunden
Erindring), blev ikke anerkendt eller paaagtet, da den først fremkom. Den blev
afvist af danske Psykologer, og Forfattere som Wundt, Meinong og Bergson gik
uden videre ud fra, at Association, Sammenligning eller Dom altid er Betingelse
for Genkendelse. Desmere har det interesseret mig at se, at G. E.
\textsc{Müller,} der er en Autoritet, hvad Læren om Forestillingerne angaar, i
sine som Manuskript trykte psykologiske Forelæsninger (1913) § 26, beskriver
det samme Fænomen, jeg kalder umiddelbar Genkendelse, under Navn af ubestemt
Genkendelse. »Genkendelse af en Iagttagelsesgenstand er«, siger han,
»\emph{ubestemt}, naar vi ved Iagttagelsen af Genstanden blot har det Indtryk,
at vi et eller andet Sted tidligere maa have iagttaget den. Naar Genstanden
ikke blot fremtræder for os som bekendt, men desuden vækker Erindringen om sit
Navn, om de Omstændigheder, under hvilke vi tidligere har iagttaget den, o. s.
v., saa har vi en \emph{bestemt} Genkendelse.« Jeg har mod denne Beskrivelse
blot det at indvende, at Bekendthedskvaliteten ikke \emph{behøver} at være
forbunden med Bevidsthed om (»Indtrykket af«), at vi har opfattet Genstanden
før. Den \emph{kan} være uden Følge. Det er allerede en Slutning, naar vi
mener, at vi har opfattet Genstanden før, en Slutning, som drages fra
Forskellen mellem Fornemmelser med og uden Bekendthedskvalitet, selvfølgelig en
umiddelbar Slutning. Naar umiddelbar Genkendelse strax udløser en Reflex- eller
en Instinkthandling, drages en saadan Slutning ikke, og »Indtrykket« vil ikke
gøre sig gældende.

% --- p. 19 ---
\setcounter{footnote}{0}%
\opage{19}\textsc{Müller} finder (i en haandskreven Tilføjelse til den trykte
Forelæsning) Forklaringen af ubestemt Genkendelse deri, at Eftervirkninger af
tidligere Iagttagelser og maaske utydelige Reproduktioner af disse udløses ved
den nye Fornemmelse, eller at Emner, vi tidligere har opfattet, ved Gentagelse
opfattes lettere og hurtigere end før. Ganske paa samme Maade havde jeg ogsaa
henvist til Øvelsens Lov, der finder Anvendelse baade paa den psykologiske og
paa den fysiologiske Side af Fænomenet. Det kan naturligvis ofte være
vanskeligt at afgøre, om og i hvilken Grad en utydelig Forestillingsassociation
finder Sted. Derfor har jeg betegnet umiddelbar Genkendelse som et
Grænsetilfælde, og jeg har i min Psykologi sammenlignet Forholdet mellem
umiddelbar Genkendelse og Forestillingsassociation med Forholdet mellem
Tangenten og Sekanten: jo nærmere Sekantens to Skæringspunkter med Cirklen
kommer til at ligge ved hinanden, desmere nærmer Sekanten sig til Tangenten.
Gentagen og nøjagtig Iagttagelse har vist mig mange Grader her, og der er al
Anledning til nærmere Undersøgelse.

Allerede før \textsc{Müller's} Forelæsninger tryktes, havde \textsc{Betz}
(Archiv für die gesammte Psychologie. 1910, p. 267) udtalt: »Erst die
vielzitierte Abhandlung von Høffding über Wiedererkennen zeigte, dass die Dinge
bei weitem nicht so einfach liegen, als es den Assoziationstheoretikern
schien«. --- Allerede Nyplatonikeren Olympiodoros%
% footnote 1 on p. 19
\footnote[1]{\textit{Olympiodori in Phædonem Commentaria,} ed. W.
\textsc{Norvin.} 1913. p. 68.} har iøvrigt omtalt Fænomenet. I sin Kommentar
til Platon's Faidon supplerer han sin Mesters Lære om Erindringen ved at gøre
opmærksom paa, at der ikke blot gives en Overgang fra et Element til et andet
paa Grund af Lighed [Lighedsassociation] eller fra et Element til et andet,
derfra forskelligt Element,
% --- p. 20 ---
\setcounter{footnote}{0}%
\opage{20}med hvilket det tidligere er forekommet sammen
[Berøringsassociation], men at ogsaa »det Samme kan forbindes med det Samme«.
Naar En f. Ex. genkender Sokrates, efter en Tid at have glemt ham, saa er det,
siger Olympiodoros, ligesaavel en Sansning, som første Gang, han saa Sokrates;
der foregaar ved saadan Genkendelse ingen Tankeovergang (οὐ γέγονε μετάβασις
γνώσεως). At Platon ikke omtaler denne Art Genkendelse, forklarer Olympiodoros
ved, at han i »Faidon« kun havde Brug for den frie (»bestemte«) Genkendelse,
ved hvilken ifølge ham »Ideerne sprang frem i Sjælen«.

b. Naar der foreligger et sammensat Emne, hvis Elementer (Dele eller
Egenskaber) vi ofte har opfattet i Rækkefølge, behøver vi ikke altid at
gennemgaa alle Elementerne paany, for at Genkendelse kan finde Sted. Successiv
Genkendelse kan erstattes ved partiel Genkendelse. Er Emnets Elementer
\textit{a b c d,} kan det være nok, at \textit{a} genkendes. Vi har her ganske
det samme Fænomen, som \textsc{Lévy-Bruhl} har betegnet som Participation.
Partiel Genkendelse har stor praktisk Betydning, da det ofte vilde tage for
lang Tid, hvis hele Rækken af Elementer skulde genkendes, før vi var sikre paa
at have Emnet for os, saa at vi kunde skride til Handling overfor det. Emnet
indeholder altsaa mere, end vi lægger Mærke til, men vi indskrænker os til det
strengt nødvendige. Vi genkender f. Ex. et Menneske paa hans Gang eller paa en
Særegenhed ved hans Dragt, et kemisk Stof ved dets Lugt o. s. v. Ligesom ved
umiddelbar, ubestemt Genkendelse er det Trangen til Handling, der forkorter
Processen. Nødvendigheden af Indgriben, hvis Genkendelsen slaar til, vejer
praktisk mere end Sikkerheden for, at Opfattelsen er rigtig. Naar spontan eller
instinktiv
% --- p. 21 ---
\setcounter{footnote}{0}%
\opage{21}Handling leder i Ulykke, er det, fordi det enkelte Kendemærke ikke
slaar til.

Partiel Genkendelse kan blive til Sanseillusion, naar det umiddelbart genkendte
Element i forskellige Tilfælde har været forbundet med forskellige andre
Elementer. Naar baade Rækken \textit{a b c d} og Rækken \textit{a f g h} er
forekommet, og jeg, naar \textit{a} kommer igen, strax stoler paa, at ogsaa
\textit{b c d} er der, kan jeg blive overrasket ved, at \textit{f g h} melder
sig istedetfor. Jeg benytter et Exempel fra min Psykologi. Paa en Vandring
langs Jammerbugten paa Jyllands Vestkyst gentog det sig for mig en hel Dag
igennem, at jeg tog Vragstumper, der ragede op af Sandet, for Mennesker. Der
foregik her en uvilkaarlig Supplering af det virkelig sete. Det Billede, der
uvilkaarlig dannede sig for mig, indeholdt andet, end der viste sig at være
givet. Suppleringen sker her ligesaa uvilkaarlig som i den bekendte Erfaring
med den blinde Plet, eller som naar man, efterat en Del af Underlæben er bleven
lokalt bedøvet, mærker et Hul eller en Sænkning i Randen af det Glas, man
drikker af. En saadan Supplering ligger, fra vort moderne Synspunkt set, til
Grund for det primitive Menneskes Participation. Det »Prælogiske«, der kun ved
en betænkelig Analogi kan betegnes som opstaaet ved Association eller ved
Dømmen, kender vi altsaa fra daglig Erfaring. Rent logisk kan Sanseillusion
formuleres som en Fejlslutning; men derfor er det alligevel uberettiget at
lægge en saadan Slutning som virkelig foregaaet psykisk Proces ind i de
elementære Fænomener, vi her har med at gøre. Vi har nu engang ikke Ret til at
identificere de Elementer, den analyserende Eftertanke finder, med dem, der
medvirker i den oprindelige Oplevelse. Her glipper Analogien mellem psykologisk
Analyse og kemisk Analyse.

% --- p. 22 ---
\setcounter{footnote}{0}%
\opage{22}Naar man formulerer Sanseillusionen, eller den partielle Genkendelse
idethele, som en Slutning, bliver det anden eller tredie Figur, vi faar.
Egenskaben \textit{B} findes baade i Emnet A og i Emnet \textit{C}; deraf
sluttes partiel eller absolut Identitet af A og \textit{C.} (Anden Figur).
Eller: Emnet \textit{B} har baade Egenskaben A og Egenskaben \textit{C;} deraf
sluttes, at der er partiel eller absolut Identitet mellem A og \textit{C.}
(Tredie Figur). Som vi skal se, har vi her to af de Maader, hvorpaa
Analogislutning logisk kan formuleres. Det betænkelige ved enhver
Analogislutning fremtræder allerede her klart. ---

I al Participation, altsaa i umiddelbar, partiel og illusorisk Genkendelse gør
der sig en uvilkaarlig Trang til Helhedsdannelse gældende. Selv om kun et
enkelt Element faktisk er givet, gør man det dog strax til et Hele. Derfor har
\textsc{Grønbech} med Rette betegnet det, Lévy-Bruhl kalder Participation, som
Helhedsrealisme. En saadan uvilkaarlig Helhedstrang er særlig karakteristisk
for Bevidsthedslivets første Udviklingstrin, der bærer Naivitetens Præg; den
forsvinder --- heldigvis --- heller ikke senere, skønt den bliver Genstand for
Analyse og Kritik. Den ligger til Grund for, hvad jeg har kaldt konkret og
praktisk Intuition. Ved konkret Intuition forstaar jeg Sanseanskuelse,
Erindring eller Fantasi. I alle disse tre Former, der ikke fra først af staar
som principielt forskellige fra hverandre, opbygges et Hele paa Foranledning af
et enkelt eller nogle faa Elementer. Da det ofte vil være en herskende Følelse,
der bestemmer det Midtpunkt, om hvilket Helhedsdannelsen finder Sted, er
konkret Intuition mere eller mindre beslægtet med praktisk Intuition, ved
hvilken alle Sindets Kræfter kan virke med, uden at der dog behøver at være
nogen Bevidsthed om Kilderne til Helhedsbilledet eller om den Energi, hvormed
det dannes og fastholdes. Al praktisk Tro,
% --- p. 23 ---
\setcounter{footnote}{0}%
\opage{23}overhovedet al personlig Koncentration om en Tanke eller en Opgave er
af denne Art. Ogsaa her mindes vi om den »prælogiske« Bevidsthed, som ifølge
Lévy-Bruhl har en rent emotionel og mystisk Karakter, da det er Stammens
hellige Overleveringer, der raader for, hvorledes den Enkelte tyder, eller
rettere integrerer, hvad han oplever%
% footnote 1 on p. 23
\footnote[1]{Smlgn. min Afhandling om \textit{Begrebet Intuition} (Vid. Selsk.
Oversigt 1914) p. 75---77 og mit Skrift \textit{Henri Bergson's Filosofi} p. 22
f. --- I Modsætning til de to nævnte Arter af Intuition sætter jeg analytisk og
syntetisk Intuition; hin udtrykker Eftertankens prøvende, denne dens samlende
Funktion.}.

c. Ved al Genkendelse spiller Lighed med det tidligere oplevede en Rolle, og
forsaavidt kan Genkendelse betragtes som et Grænsetilfælde af
Lighedsassociation. Hvad Lighedsassociation overhovedet angaar, har der været
Strid om, hvorvidt en saadan overhovedet gives. Min Stilling til dette
Spørgsmaal har, ligefra den første Udgave af min Psykologi, været den, at jeg
ikke antager nogen ren og ublandet Lighedsassociation, da der i alle
paaviselige Exempler paa Lighed stedse vil være visse Forskelligheder tilstede.
Det har ingen været mere klar over end Platon --- det er jo derfor han gør
fuldkommen Lighed til en »Idé«. Og Forskelselementerne vil kunne udløse
Berøringsassociation, som spiller en større eller mindre Rolle ved Overgangen
til den følgende Forestilling. Enhver faktisk Association vil derfor være en
Blanding af Ligheds- og Berøringsassociation, og Diskussionen maa i hvert
enkelt Tilfælde dreje sig om, i hvilken Grad Lighedsmomenter virker med. Min
Formel for Lighedsassociation er derfor \textit{a}\textsubscript{1}
+\textit{a}\textsubscript{2}, og Indices betegner de Forskelligheder, der
stedse gør sig gældende, selv hvor Ligheden er overvejende. Da det er
ejendommeligt for saadanne Overgange i Forestillingslivet, der væsentlig
skyldes Lighed, at de sker uvilkaarlig, og især naar Sindet
% --- p. 24 ---
\setcounter{footnote}{0}%
\opage{24}er stærkt optaget og bevæget, vil det være vanskeligt at paavise
saadanne Overgange experimentalt; det er ved psykologiske Forsøg meget
vanskeligt at undgaa Forberedelse og Forventning, hvorved strax
Berøringsassociation begunstiges%
% footnote 1 on p. 24
\footnote[1]{Se foruden min Psykologi: \textit{Psykologiske Undersøgelser}
(1889), p. 43---59.}.

Medens mange Psykologer (især danske) stillede sig ganske afvisende overfor
Antagelsen af en Lighedsassociation, selv med de anførte Begrænsninger,
udtrykte G. E. \textsc{Müller,} der ganske særlig har undersøgt
Forestillingsforbindelser, sig mere forsigtigt. Han fandt det (i sine ovenfor
omtalte Forelæsninger) ikke sandsynligt, at der gaves ren Lighedsassociation;
naar et Portræt fremkalder Forestillingen om selve Personen, kan det ske ved,
at det først fremkalder Navnet og dette derefter Personen. Der er jo ingen
Tvivl om, at det kan gaa til paa denne Maade. Jeg har dog ikke kunnet
overbevise mig om, at det passer paa alle Tilfælde; jeg kan jo godt genkende et
Portræt, selv om jeg har glemt Navnet, som jeg maaske ogsaa efter Genkendelsen
har Møje med at genkalde mig.

I samme Aar, som Müller lod sine Forelæsninger trykke som Manuskript, deltog
han som Forsøgsperson i nogle Forsøg, som lededes af hans Elev Louise Schlüter,
og som skulde prøve Forholdet mellem to Metoder ved Undervisning i et fremmed
Sprog, af hvilke den ene bygger paa direkte Forbindelse mellem en Genstand og
det fremmede Ord (Anskuelsesmetoden), den anden først nævner det hjemlige Navn
paa Genstanden og lader det genkalde det fremmede Ord. Det viste sig her, at
naar Dagen efter, at en Genstand var forevist og det fremmede Ord lært, en
lignende Genstand forevistes, dukkede det fremmede Ord
% --- p. 25 ---
\setcounter{footnote}{0}%
\opage{25}op uden nogen Støtte af det hjemlige Ord. Dette forudsætter, at den,
den foregaaende Dag, sete Genstand reproduceredes ud fra den nu sete Genstand,
skønt den ikke var ganske den samme. Dette betragter Louise Schlüter som Bevis
for, at der kan gives virkelig Lighedsassociation%
% footnote 1 on p. 25
\footnote[1]{\textit{Beiträge zur Prüfung der Anschauungs- und der
Übersetzungsmethoden} (Zeitschrift für Psychologie. 1913) p. 73; 103.}, og
Müller har erklæret sig enig heri%
% footnote 2 on p. 25
\footnote[2]{I haandskrevne Tilføjelser til et Exemplar af Forelæsningerne, som
Professor Axel Herrlin velvillig har stillet til min Disposition.}.

Jeg skal hertil kun bemærke, at der ved Iagttagelser og Forsøg over dette
Forhold maa skelnes mellem de forskellige Arter af Lighed. I Louise Schlüter's
Exempel var det, saavidt jeg kan se, Kvalitetslighed, det drejede sig om: den
Genstand, der forevistes den anden Dag, var ikke den samme som den, der saas
den første Dag, men »ein ähnliches Exemplar derselben Art«. Det var altsaa
ikke, hvad jeg kalder Dækningslighed, der her gjorde sig gældende. Men da
Forholdet mellem Exempler paa samme Art for det meste er Forholdslighed, idet
det ikke behøver at være Lighed med Hensyn til en enkelt Del eller Egenskab,
der er afgørende for Artsbegrebet, vilde det være værd at undersøge, hvorledes
i Tilfælde som det Schlüter'ske Kvalitetslighed og Forholdslighed forholder sig
til hinanden.

I Grunden er det Schlüter'ske Tilfælde ikke forskelligt fra de Tilfælde, paa
hvilke Müller i sine Forelæsninger (§ 20) har bygget den Associationsform, han
kalder aktiv Substitution, og som er en Lighedsassociation baseret paa en
Berøringsassociation. Naar \textit{a} er associeret med den derfra forskellige
Forestilling \textit{b,} saa kan en Forestilling \textit{α}, der ligner
\textit{a}, ogsaa fremkalde \textit{b,} hvad enten Forholdet mellem \textit{a}
og \textit{α} er Kvalitetslighed eller Analogi.

Saadan aktiv Substitution er af stor Interesse for vor
% --- p. 26 ---
\setcounter{footnote}{0}%
\opage{26}Undersøgelse over Analogien i erkendelsesteoretisk Henseende. Ti naar
det viser sig, at \textit{a} og \textit{α} praktisk set spiller samme Rolle,
hvilken Lighedsart det saa er, ser vi her en psykologisk Mulighed for en
Forvexling af Identitet og Analogi, der viser sig saa skæbnesvanger i
Tænkningens Historie. Der vil fra først af ikke være nogen Tvivl om Gyldigheden
af Overgangen fra \textit{α} til \textit{b.} Praktisk set er der en
Tilbøjelighed til at behandle al Lighed som Identitet. Det ses allerede af det
meget diskuterede Fænomen, at brændt Barn skyr Ilden. Efter Stuart Mill
foreligger her blot en Berøringsassociation. Men Barnet skyr foreløbig alle
Genstande, der ligner den, hvorpaa det engang har brændt sig.

4. I Barnets Leg og Sprog, i kunstnerisk Anskuen og i Geniets Virken gør der
sig uvilkaarlige Analogier gældende, der ligesaalidt som de i det Foregaaende
omtalte skyldes udtrykkelig Sammenligning eller Slutten, og heller ikke skyldes
en Association, ved hvilken ethvert enkelt Led opfattes for sig selv og først
derefter forbindes med andre Led.

a. Barnets Leg skyldes ofte Efterligning af de Voxnes Færd, men Barnet kan
fordybe sig i den med en saa fuldstændig Indleven, at Forskellen mellem hvad
den »er« og hvad den skal »betyde« slet ikke gør sig gældende. Afbrydes Legen
udefra, kan det vække Smerte og Harme. Et Barn legede Husmoder og havde dækket
op paa sit Legebord med forskellige indbildte Retter. Da det nu blev Husets
Spisetid, og man vilde have Barnet til at sætte sig ved det rigtige Bord med
den øvrige Familie, blev det heftig opbragt og gav sig til at græde, »for en
Mor løber da ikke saadan fra det Bord, hun skal sørge for«. Forargelsen viser,
hvor dyb Indlevelsen i Legen havde været. Paa lignende Maade kan det gaa
Tilskueren ved et Drama,
% --- p. 27 ---
\setcounter{footnote}{0}%
\opage{27}naar han pludselig forstyrres i sin fulde Opgaaen i, hvad han ser og
hører. Med Urette har man villet forklare Tiltrækningen ved et tragisk Drama
derved, at Læseren eller Tilskueren skulde være sig bevidst, at han staar
udenfor, hvad der fremstilles. En saadan Viden vilde gøre den fulde Optagethed
umulig. Optagetheden forudsætter, at vi for en Stund identificerer vor Skæbne
med den, som fremstilles. Det er en Illusion, en »Participation«, som maaske
hæves, naar Tæppet falder, eller Bogen slutter, men hvis Genklang eller
»Perseveration« kan gøre sig gældende længe, idet vi kan have en Tendens til
atter og atter at leve i samme Stemning, som da vi var umiddelbart under
Dramaets Indflydelse.

Ogsaa udenfor Legen viser Børn Evne til analogiske Oplevelser. En Tilstand
eller Stilling, der har medført Tilfredsstillelse for dem, vil de have fortsat,
enten saaledes, at de selv fremdeles er med deri, eller, mere uinteresseret, at
Tilstanden blot fortsættes efter samme Skema som før. Det er i begge Tilfælde
Ækvivalenter, der søges. \textsc{AVENARIUS} har i sit interessante Hovedværk
»Kritik der reinen Erfahrung« søgt at vise, at der her --- i Trangen til
Fortsættelse af en Tilstand, selv om man derved kommer ud over de gamle Former
--- gør sig en almenmenneskelig Tendens gældende. Han anfører følgende morsomme
barnepsykologiske Exempel (der afgjort ikke hører til de uinteresserede). En
lille Dreng sidder nok saa bekvemt i en Lænestol og lader sin Gæst, en lille
Pige, staa op. Hans Moder bebrejder ham det, og han rømmer Pladsen, hvorpaa
Pigen indretter sig i Stolen med sit Legetøj paa Skødet. Drengen ser lidt paa
hende og foreslaar saa, at hun skal lægge Legetøjet paa Gulvet og lade ham
sidde paa sit Skød%
% footnote 1 on p. 27
\footnote[1]{\textit{Kritik der reinen Erfahrung}\textsuperscript{2} II p. 493,
kfr. I, p. 136. --- \textsc{Avenarius} finder den samme Tendens i saadanne
Tilfælde, hvor Mennesker ved pludselige Skæbneslag eller ved Alderdomssvækkelse
kun kan fortsætte deres aandelige Liv ved Hjælp af Trosformer, de i deres
kraftige Alder havde forladt. \textit{Kritik der reinen
Erfahrung}\textsuperscript{2} II, p. 295.}. --- Til den uinter%
% --- p. 28 ---
\setcounter{footnote}{0}%
\opage{28}esserede Fortsættelse af en Situation hører derimod følgende Exempel.
En lille Dreng red paa sin Bedstefaders Knæ, men maatte afbryde Situationen, da
han skulde i Seng. Lidt efter kom han med stor Iver styrtende tilbage med sin
Legetøjsbjørn, som han anbragte paa Bedstefaderens Knæ. Saaledes blev der sat
et nyt Led ind i Forholdet, der ellers maatte ophøre. Maaske har der ogsaa
været Sympati med den gamle Mand, hvis Skød nu blev tomt. ---

Ved Barnets Sprogtilegnelse spiller ikke blot Erindring om og Efterligning af
de Voxnes Tale, men ogsaa uvilkaarlig Analogiseren en Rolle. Barnet har f. Ex.
hørt Udtrykket »Jakobs Hat« sammen med det tilsvarende Synsbillede og danner nu
paa egen Haand Udtrykket »Peters Hat« ved given Lejlighed, uden at
Almenforestillingen »Hat« er et nødvendigt Mellemled. Børn danner desuden ofte
paa egen Haand regelmæssige Verbalformer, naar de ikke har lært de
uregelmæssige. Bag efter kan de saa øve sig i at bruge de uregelmæssige Former,
Sprogbrugen kræver, istedetfor de regelmæssige, de selv i uvilkaarlig
Analogiseren havde dannet%
% footnote 1 on p. 28
\footnote[1]{\textsc{Otto Jespersen:} \textit{Language,} p. 130 f.}.

b. Den Kendsgerning, at man forholdsvis sent anerkendte genial Skaben som Andet
og Mere end Produkt af mekanisk Sammenføjning af Enkeltheder i en bevidst Hensigt%
% footnote 2 on p. 28
\footnote[2]{Smlgn. f. Ex. den Maade, hvorpaa \textsc{Holberg} selv forklarer
Tilblivelsen af sine forskelligartede Værker. Se mine \textit{Mindre Arbejder}
II, p. 150 ---154.}, viser, hvor vanskeligt man har havt ved at fatte
uvilkaarlig psykisk Helhedsdannelse. Man betragtede alle sjælelige Processer
fra den analyserende Eftertankes Stand%
% --- p. 29 ---
\setcounter{footnote}{0}%
\opage{29}punkt. Efterat Rousseau (især i Artiklen »Génie« i sit Dictionnaire
de Musique) havde beskrevet Geni som Evne til uvilkaarlig Skaben og bønfaldt
»l'homme vulgaire« om ikke at misbruge dette sublime Ord, var det
\textsc{Kant,} der fastslog Begrebet i den Form, i hvilken det nu almindelig
bruges, idet han opfattede Geni som den sjeldne Originalitet i Naturbegavelse,
der gør det muligt at frembringe Værker af forbilledlig Karakter, Værker, der
ikke kan efterlignes, men som kan vække Originaliteten hos andre
Naturbegavelser (Kritik der Urtheilskraft § 40). Geni er, ligesom Instinkt, en
umiddelbar Forening af Trang og Evne. Det Uvilkaarlige og det Forbilledlige
udgør i Forening Geniets Væsen, \textsc{Goethe} glædede sig (Dichtung und
Wahrheit. 19. Buch) over, at det bedste Ord for uvilkaarlig og mønsterværdig
Skaben saaledes blev reddet. --- En lignende Skæbne som Ordet Geni truer Ordet
Fantasi, naar man bruger det om al fri Forestilling, selv om den blot er
Erindring, og ikke begrænser det til at betyde Evne til Nydannelse af konkrete
Individualforestillinger. \textsc{Meumann} har sikkert Ret, naar han mener%
% footnote 1 on p. 29
\footnote[1]{\textit{Intelligenz und Wille}\textsuperscript{2}, p. 136.}, at
den snevrere Betydning af Ordet Fantasi først opstod under Indflydelse af den
nye Tids mere indtrængende Studium af kunstnerisk Frembringen. Kunstnerisk
Fantasi fører til Værker, der for en Betragter kan synes at være Frugt af en
Eftertanke, der med fuld Bevidsthed lægger de enkelte Træk ind paa en saadan
Maade, at de samvirker med andre Træk til Dannelse af et Hele. I Anledning af
et Træk, som Solger havde fremhævet i sin Analyse af Karaktererne i
»Wahlverwandtschaften«, sagde \textsc{Goethe:} »Ich habe selber nicht daran
gedacht, aber es liegt darin«. (Gespräche mit Eckermann. 21. Januar 1827).

% --- p. 30 ---
\setcounter{footnote}{0}%
\opage{30}Kunstneren kommer saaledes uvilkaarlig til at sige mere, end han er
sig bevidst. At han analogiserer, og derved kommer til at forbinde Elementer,
der ellers var adskilte, mærker han ikke, før maaske senere, hvis den prøvende
Eftertanke vaagner. Ved en Indleven, eller, som Tyskerne siger, Indfølen
(Einfühlung) i Emnet, der uvilkaarlig leder ham til at søge og finde Udtryk for
dets Værdi, kan han vække, om ikke en selvstændig Skaben, saa dog en
selvstændig Evne til en lignende Indleven som hans egen. I denne sidste
Indleven øver ikke blot Emnets Indhold, men ogsaa den Form, under hvilken det
fremtræder i Værket, en afgørende Virkning, --- ligesom Kunstnerens Indleven i
Emnet fra først af fører ham til en formende Virksomhed, der kunde lade Emnet
fremtræde i sin hele Fylde. I den Definition, \textsc{Julius Lange} (Om
Kunstværdi. 1876) gav af Kunstværdien, som den Værdi, Emnet har for Kunstneren,
og som den gennem hans Værk faar for Beskueren, kom (som jeg i en Anmeldelse af
Lange's smukke Bog i »Nær og Fjern« 1876 søgte at paavise) Fremstillingsformen
ikke til sin Ret; den betragtede Lange som noget rent ydre og mekanisk. Hos det
kunstneriske Geni er fra først af Trang (Følelse af Værdi) og Evne (formende
Kraft) forbundne --- og det paa den inderligste Maade, idet der er Trang til at
bruge Evnen, og ikke blot til at føle Værdien, og idet Følelsen af Værdi
forudsætter en Evne til at aabne sig for, hvad der kan fremkalde denne Følelse.
\textsc{Friedrich Vischer} har med Rette hævdet, at Indfølingen baade maa
bestemme Værkets Form og dets Indhold. Lange har ved at betone Indlevetheden i
Emnet peget paa det Uvilkaarlige i Fantasien og hos Geniet, men dertil hører
ogsaa den uvilkaarlig formende Virken. Det kunde efter Lange's Defini%
% --- p. 31 ---
\setcounter{footnote}{0}%
\opage{31}tion se ud, som om Kunstværket blot var et ydre, uvæsentligt Middel,
for at Andre end Kunstneren selv kunde erfare Emnets Værdi. Paa saa udvortes
Maade er Trang og Evne ikke knyttede sammen, selv om Gennemførelsen af den
formende Virken paa mange Punkter kan muliggøre, ja kræve bevidst Overlæg og Valg.

Begrebet »Einfühlung« kendte hverken Lange eller jeg, da vi førte vor
Diskussion. Det var nogle Aar iforvejen bleven opstillet af den tyske Æstetiker
\textsc{Robert Vischer,} men blev først mere almindelig bekendt gennem
\textsc{Volkelt's} og \textsc{Friedrich Vischer's} Skrifter%
% footnote 1 on p. 31
\footnote[1]{\textsc{Johannes Volkelt:} \textit{Der Symbolbegriff in der
neueren Ästhetik.} 1876, p. 55. --- \textsc{Friedrich Vischer:} \textit{Das
Symbol} (Festskrift til Zeller. 1887) p. 170---187. --- \textsc{Rousseau} har
kendt Indfølelse paa første Haand. Hans Naturglæde udspringer deraf. Smlgn.
\textit{Confessions} (éd. Bever) I, p. 263 f.: Je dispose en maître de la
nature entière; mon coeur, errant d'objet en objet, \textit{s'unit,
s'identifie} à ceux qui le flattent. --- Men hos Rousseau var Trangen større
end Evnen. Et Kunstværk kunde han ikke forme (uden det Kunstværk, hans egne
Confessions er). Men hans inderlige Følelse gav Stødet til moderne Poesi.}.

Medens Billedet for den primitive Bevidsthed stod som Bærer af en magisk Kraft,
der gjorde dets Besidder til Herre over Emnet, er det i Kunsten Bærer af en
Værdi, som fødtes i Kunstnerens Sjæl og gennem hans Værk fødes igen i
Beskuerens Sjæl. Begge Steder er der en Analogi, der virker paa uvilkaarlig
Maade, og som først af den prøvende Eftertanke ses som Analogi. At man ikke har
forstaaet det uvilkaarlige Sjælelivs Maade at virke paa, førte til Miskendelse
baade af Geniets og af det primitive Menneskes Natur. Men den Forstaaelse, der
kan vindes af det uvilkaarlige Sjæleliv, sker selv ved Hjælp af Analogi, da
Forstaaelsen gaar ud fra det differentierede Bevidsthedsliv, som kun kan
tydeliggøre sig de primitive Synteser ved at opløse dem.

Svarende til Billedets magiske Betydning staar i Kun%
% --- p. 32 ---
\setcounter{footnote}{0}%
\opage{32}sten dets symbolske Betydning, men i Religionens Historie brydes helt
op i den nyeste Tid Magi og Symbolik.

5. Symbolet beror paa Participation, idet en Del staar som Udtryk for det Hele
og virker som saadant. Det Smaa kan blive Symbol for det Store, naar det i sig
frembyder Forhold, der har Lighed med Forhold indenfor det Store. For den store
Humor kan alt enkelt, hvor begrænset det end er, blive Symbol ved at ses som
Del af et stort Hele og som frembydende samme Type som Helheden. »Det Ringeste
forsvarer sin Plads i Livets Krans«. Men Livets Krans er aldrig helt bunden, og
den Mulighed er stedse tilstede, at nye Blomster skal indflettes. Dette vil den
store Humor ikke være blind for; den ser humoristisk paa sine egne Kranse%
% footnote 1 on p. 32
\footnote[1]{\textit{Den store Humor.} Kap. 6 (Forstaaelse og Humor).}, og den
er enig med \textsc{Goethe} (Gespräche mit Eckermann. 5. Juli 1827) i, at der i
Billeder, der bruges som Udtryk for Livet, stedse bliver noget for Forstanden
inkommensurabelt tilbage. Billedets Værdi beror netop fra en væsentlig Side
paa, at det giver som et Hele, hvad der for prøvende Eftertanke staar som en
uafsluttet Sammenhæng af Elementer, maaske stedse skiftende Elementer. Fra
denne Side set giver Billedet mindre end den altid fremskridende Erfaring og
Eftertanke. Som Helhedsbillede giver det mere, idet det som saadant har
Elementer, der ikke findes i det Emne, hvorfra det er hentet. Naar Gud kaldes
Fader, kommer Elementet Almagt med ind, der desværre mangler hos enhver i
Erfaringen given Fader, som ofte maa kæmpe haardt baade for sig og for sine.
Eller (for at bruge \textsc{FR. Vischer's} Exempel) naar Løven er Symbol paa
Højsind, har den en Egenskab, den virkelige Løve mangler. Pseudohallucinationer
egner sig derfor ofte bedre som Grundlag for Symboler end egentlige Hallucina%
% --- p. 33 ---
\setcounter{footnote}{0}%
\opage{33}tioner, der i Regelen tager alle Virkelighedens Elementer med.

Den religiøse Symbolik er især karakteriseret ved, at dens Følelsesgrundlag har
dannet sig under Kampen for det, der for vedkommende Menneskegruppe er det
højeste Værdifulde. De Billeder, i hvilke der findes Udtryk for Erfaringerne
under Kampen for Tilværelsen, er hentede fra Naturens og Livets store
Modsætninger: Lys og Mørke, Liv og Død, Glæde og Sorg, Magt og Afmagt, Godt og
Ondt; og der gør sig, alt som de mest elementære Naturreligioner viger Pladsen
for højere Former for Religion, en Trang gældende til at se disse Modsætninger
repræsenterede og virkende gennem mere eller mindre individualiserede Væsener.

Men ligesom Begrebet Analogi idethele er ogsaa Begrebet Symbol først dannet paa
den prøvende Eftertankes Standpunkt. Forskellen mellem Del og Helhed, mellem
Tegnet og det Betegnede gør sig ikke gældende fra først af. Hvad der for
Kunstneren og den kunstneriske Beskuer kun gælder i Betagethedens Øjeblikke, er
for den positivt-religiøse Bevidsthed, saalænge den bestaar i sin fulde
Ejendommelighed, den stadige Tilstand. Enhver Skelnen mellem Del og Helhed,
Billede og Emne er kættersk. Med Urette mente Hegel, at der ved al symbolsk
Opfattelse maatte føles en Utilfredshed paa Grund af det stedse tilstedeværende
Misforhold mellem Form og Indhold. Idet Hegel forudsatte en saadan Utilfredshed
allerede paa primitive Trin og idethele indenfor den positivt-religiøse
Bevidsthed, blev det ham forholdsvis let at vise, at det var en Befrielse for
den religiøse Bevidsthed at gaa over til Former, der kunde bestaa for prøvende
Eftertanke, og Forskellen mellem Religion og Filosofi blev for ham tilsidst
% --- p. 34 ---
\setcounter{footnote}{0}%
\opage{34}kun en Formforskel, medens det aandelige Indhold skulde være det samme%
% footnote 1 on p. 34
\footnote[1]{\textit{Den nyere Filosofis Historie.}\textsuperscript{2} III, p.
182---186.}. Hegel lægger den Spaltning, som Eftertanken kan føre til, ind i
selve den religiøse Bevidsthed, for hvem Forholdet mellem Del og Helhed, Tegn
og Emne tværtimod er et Identitetsforhold. Det benægtes f. Ex. fra den positive
Religions Side, at Udtrykket Fader brugt om Gud er billedligt. Langt fra at det
skulde være Erfaringer om det menneskelige Faderforholds Betydning, der skulde
ligge til Grund, naar Gud kaldes Fader, forholder det sig netop omvendt: det er
Gud, der er Grund og Forbillede for ethvert Faderforhold i Verden (Brevet til
Efeserne III, 15). Formedelst denne Identitet er det muligt, at Handlinger, der
for Eftertanken staar som symbolske, som Bøn, Offer og Sakrament, kan antages
at øve Indflydelse paa Guddommen. Dette er konsekvent, thi formedelst
Identitetsforholdet er der en real Sammenhæng mellem Tegn og Genstand. Dertil
kommer, at den positivt-religiøse Symbolik ikke blot er Resultat af den
Enkeltes, men for en hel Menneskegruppes Livserfaringer og derved for nye
Slægter optræder som givet og selvfølgeligt Udtryk for, hvad der kan erfares om
Menneskelivet og dets Kaar.

Begreber som Analogi og Symbol er Korrelatbegreber, udtrykker et gensidigt
Forhold. Som vi allerede har set ved Fadersymbolet, vender derfor Forholdet sig
ganske naturlig om. Fra den Erfaringens Verden, i hvilken Menneskelivet føres,
hentedes der Symboler for en højere Verden; men tilsidst bliver selve den
Verden, hvorfra Symbolerne oprindelig hentedes, til et Symbol. Hos
Middelalderens Mystikere træder denne Vending frem paa karakteristisk Maade.
For Hugo a St. Victore er hele Legemverdenen et Symbol paa den aandelige
Verden. Og idet Goethe benytter
% --- p. 35 ---
\setcounter{footnote}{0}%
\opage{35}middelalderlige Udtryk for at betegne Opsvinget over alt Endeligt,
udtaler han

Alles Vergängliche ist nur ein Gleichniss.

Den middelalderlige Tænken bevægede sig, som man træffende har sagt, ligesaa
hyppig fra Symbol til Symbol, som fra Erfaring til Symbol. Der var en Gang for
alle vakt en Grundstemning, som stedse søgte nye Udtryk, og i Forhold til
hvilken forskellige Livserfaringer kunde betyde det samme. Her virker det
emotionelle Grundlag, som idethele er bestemmende for al »Participation«, al
»prælogisk Mentalitet«. Først kritisk Eftertanke stiller den Fordring til
religiøse Forestillinger, som Schleiermacher hævdede, at hvert enkelt Symbol
skal dannes for sig paa Grundlag af særlig Oplevelse. I Troslæren maa, siger
han, den ene Sætning ikke udledes af den anden, men skal opstilles umiddelbart
i Henhold til en enkelt bestemt Oplevelse, og saa først bagefter sammenlignes
med andre, forud opstillede Sætninger. (Der christliche Glaube § 17). Saalænge
den én Gang udviklede Grundstemning holder sig, vil derimod nye Symboler atter
og atter udspringe at det, det ene udviklet i Analogi med det andet. Saaledes
har \textsc{Yrjo Hirn} (Det heliga Skrinet. Studier i den katolska Kyrkans
Poesi och Konst. 1909) paa en smuk Maade vist, hvorledes den inderlige Dyrkelse
af Madonna fører til en Række af Billeder, der udtrykker Grundstemningen fra
forskellige Sider: Hostiens Hjemsted i et Skrin paa Alteret, det Sted, hvor det
Guddommelige fremtræder i legemlig Form, blev Symbol for Madonna, i hvis Skød
Gud blev Menneske. Men Madonna er tillige Morgenstjernen, der bebuder Solens
Komme, og hun er Vaaren, hvor Lys og Varme afløser Vinterens Mørke og Kulde.
Den mægtige Følelsesexpansion skyder stadig nye Skud.

% --- p. 36 ---
\setcounter{footnote}{0}%
\opage{36}Den uvilkaarlige Billeddannelse er Forudsætning for al Dogmatik. Ikke
alle Symboler bliver til Dogmer. Der bliver først Tale om Dogmer, naar visse
Symboler staar som Normer, de Enkelte bør følge under deres Forsøg paa at tyde
deres Livserfaringer. Der foretages da et Valg, og hvad der vælges staar fast i
Kraft af det religiøse Samfunds Autoritet overfor sine Medlemmer. Dogmet hviler
paa en Blanding af Oplevelse, Eftertanke og Autoritet. Ud fra det dogmatisk
Fæstnede føres saa en Kamp mod frit udviklede Symboler, saavel som mod
Forsøgene paa at opfatte selve Dogmerne som blotte Symboler. Dogmatikens Kamp
mod Symboliken er en Kamp mellem »det er« og »det betyder«.

En fri Symboldannelse indenfor religiøs Ramme frembyder Jesu Parabler, i hvilke
han gennem Fortællinger eller Træk, der er grebne ud af Livet, anskueliggør det
forjættede Gudsriges Forhold og Fordringer. En etisk-religiøs Livsvisdom er her
udtalt i poetisk Form. Det billedlige er valgt med velberaad Hu og gaar ofte
(især i Johannesevangeliets Parabler) over fra Symbolik til Allegori.

Symbolet bliver til Allegori, jo mere der udtrykkelig skelnes mellem de to Led
i Symboliseringen. I Allegorien er Forholdet mellem de to Led mere udvortes end
i Symbolet. Den poetiske Symbolik udspringer, ligesom den religiøse, af den
Betagethed, et Emne har vakt; men den udelukker ikke, som den religiøse
Billeddannelse i Religionernes klassiske Tider, at Betagetheden rytmisk kan
skifte med klar Bevidsthed om den Billeddannelse, der har fundet Sted.

Hos intuitive Naturer, hvor Anskuelse og Fantasi har Overvægt over Eftertanke,
vil der være en Tilbøjelighed til at leve i Symboler saa længe som muligt.
Endnu stærkere
% --- p. 37 ---
\setcounter{footnote}{0}%
\opage{37}vil denne Tilbøjelighed være hos kirkelige Naturer, for hvem det er
en Trang at fastholde overleverede Symboler, ikke blot paa Grund af deres
Anskuelighed, men fordi de er det overleverede Udtryk for Fædrenes Tro, altsaa
paa Grund af en Solidaritetsfølelse overfor Fortiden. De hører en Aarhundreder
gammel Tone i de kirkelige Hymner og genfinder Slægtens dyrekøbte Erfaringer i
de kirkelige Dogmer. Dog maa det mærkes, at det ikke blot kan være intellektuel
Kritik, der sprænger de overleverede Former; det kan ogsaa være en ny og stærk
Følelsesbevægelse, der gør det umuligt at finde sine Livserfaringer udtrykte i
disse Former%
% footnote 1 on p. 37
\footnote[1]{Smlgn. min \textit{Religionsfilosofi} § 95---96. ---
\textit{Oplevelse og Tydning} p. 102---107.}.

\begin{center}\large II. Analogi og Logik.\end{center}

6. I de uvilkaarlige Analogier er Ligheder og Forskelligheder fra først af
saaledes indvævede i hverandre, at det kan være vanskeligt at finde den Grænse,
hvor Ligheden ender og Forskelligheden begynder. Det bestemte Forsøg paa at
finde en saadan Grænse gennem skarp Bestemmelse af de forskellige Grader af
Lighed (og Forskel) betegner den videnskabelige Tænknings Begyndelse. Platon
staar i Tænkningens Historie som den Første, der har gjort opmærksom paa, at
der er mange Trin op fra de for umiddelbar Opfattelse fremtrædende Ligheder til
den rene Identitet, den »Lighedens Ide«, som ene gør streng Videnskab mulig.
Hvad Parmenides under Navn af »det Ene« havde proklameret som eneste Betingelse
for sand Viden, bestemmes af Platon nøjagtigere som Identitet. --- I min
Afhandling om den platoniske Dialog »Parmenides« (Vid. Selsk. filosofiske
Meddelelser Nr. 3) udtaler jeg mig dog maa%
% --- p. 38 ---
\setcounter{footnote}{0}%
\opage{38}ske for stærkt, naar jeg siger (p. 28), at Platon »udtrykkelig
erklærer« Enhed og Identitet for det Samme; forsigtigere hedder det p. 30, at
det »røbes«, at Enhed og Identitet er det samme. Sagen er, at Platon i den
underlige anden Del af Dialogen tilsyneladende nægter dette, men da han kun
begrunder denne Nægtelse ved, at \textit{flere} Ting (τὰ πολλά) kan være Et (ɔ:
gaa ind under samme Begreb) uden at falde sammen, men ikke fører noget Bevis
for »det Enes« Vedkommende, om hvilket der netop er Tale, finder jeg her en
uvilkaarlig Indrømmelse. Saa meget er vist, at Parmenides' »det Ene« er dunkelt
og mystisk, naar det ikke betyder Identitet, altsaa udtrykker al Tænkens
Grundlov og tillige al Tænkens højeste Ideal.

For Platon var nu Lighedens Ide ikke blot Grundlov og Ideal, men Udtryk for den
sande Væren, og Lighedsgraderne var ligesaa mange lavere og højere Trin for
Existens. Alligevel var der ved Opstilling af denne Ide givet en Anvisning til
Opløsning af de uvilkaarlige Analogier. Naar denne Opløsning foretages med
haard Haand, træder absolut Lighed og absolut Forskel i Modsætning til
hinanden, Identitet i Modsætning til Kaos, og det bliver Erkendelsens Opgave at
finde en Vej mellem disse Modsætninger. Det er karakteristisk for den skarpe
græske Tanke, at ikke blot Ideen om absolut Identitet, men ogsaa Ideen om
absolut Forskellighed fremdrages, om end i polemisk Hensigt (af Zenon). Og mod
Slutningen af sin Løbebane indsaa Platon, at Identitetsprincipet tilsidst stod
som afmægtigt overfor de mangfoldige Forskelligheder, som Tilværelsen
frembyder, hvorfor han henviser til Analogiseren som en Slags »uægte Tænken«
(λογισμὸς νόθος. Tim. p. 52 B). Analogi staar da her som en sidste Udvej, hvor
ingen højere Lighedsgrader kan paavises. Medens den i den
% --- p. 39 ---
\setcounter{footnote}{0}%
\opage{39}mytologiske Tænken var alt, kommer den nu til at betegne det laveste
Erkendelsestrin. Programmet for europæisk Tænken var hermed givet. Og det
bliver da Opgaven at paavise de forskellige Mellemformer mellem kaotisk
Forskellighed og absolut Identitet, thi det er i disse Mellemformer, at
menneskelig Tænkning i Virkeligheden bevæger sig. Rangfølgen af disse
Mellemformer (af hvilke jeg i »Den menneskelige Tanke« har beskrevet sex) vil
bero paa, hvorvidt de foreliggende Emner kan ordnes i et saadant indbyrdes
Forhold, at Slutning bliver mulig. Det vil da vise sig, at jo mere Slutning
bliver mulig, des mere nærmer vi os til den rene Identitet. Det viser sig
tillige, at Identitet ikke længere staar som den højeste Existensform, men som
den Tankeform, der gør begrundet Erkendelse mulig uden hvert Øjeblik at benytte
Iagttagelsens Hjælp. Tydeligst dæmrer Identitetsprincipet i saadanne Rækker,
hvor hvert Led forholder sig til sin Forgænger som dette til sin Forgænger. Er
f. Ex. \textit{A} Subjekt til \textit{B} og \textit{B} til \textit{C,} eller
\textit{A} mindre end \textit{B} og \textit{B} mindre end C, saa kan \textit{B}
skydes ud, og vi faar en ny Dom. Forudsætningen er \textit{B's} Identitet med
sig selv, hvad enten det sættes i Relation til \textit{A} eller til C. Men en
saadan Forudsætning gælder for hvert Led i enhver Række, vi opstiller.
Forsaavidt ligger Identitetsprincipet til Grund ved enhver Rækkedannelse; men
Erkendelse bestaar ikke blot i Dannelse af enkelte Begreber, men i Paavisning
af forstaaelige Forhold mellem Begreber, i rationelle Rækkedannelser.

Udtrykket »Participation«, der saa heldig er valgt til at betegne uvilkaarlig
Analogi, kan godt beholdes, efter at Logiken er begyndt at tale et Ord med
angaaende Forudsætningerne for Tankegangen. Det betegner en partiel Lighed, der
i den uvilkaarlige Analogiseren blev gjort til en absolut
% --- p. 40 ---
\setcounter{footnote}{0}%
\opage{40}Lighed trods de Vanskeligheder, der er ved at begrunde Retten til ud
fra Lighed paa et enkelt Punkt at antage en Overensstemmelse idethele. Det er
netop disse Vanskeligheder ved »Participation«, der stiller videnskabelig
Tænken dens Opgaver.

Søger man en Forudsætning, der trods alt ligger til Grund baade for den
førvidenskabelige Tænkens uvilkaarlige Analogiseren og for den videnskabelige
Søgen efter begrundede Overgange, og som ligger til Grund for alle
Rækkedannelser ligefra den kaotiske Forskelsrække til den absolute
Identitetsrække, saa ligger en saadan Forudsætning i Syntesens Lov, der gælder
baade i empirisk Psykologi og i rationel Erkendelsesteori. Det ligger i
Begrebet Række --- hvad enten Leddene holdes sammen ved blot Hukommelse eller
ved begrundet Forstaaelse af Overgangene --- at forskellige Led sammenfattes
med hverandre, enten som forskellige eller som lignende. Hvis Tanken stedse
tabte de foregaaende Led, naar den gik over til de følgende, vilde ingen
Rækkedannelse ske. Jo mere det bliver muligt at begrunde Overgangen fra Led til
Led, des mindre en Rolle behøver den blotte Hukommelse, der er enevældig ved
den kaotiske Forskelsrække, at spille; Overgangene sker da i Kraft af en Lov.

Syntesens Grundlov raader paa alle Trin af Forestillingsforbindelse, og derved
bliver det muligt at sammenligne disse forskellige Trin og paavise Forholdet
mellem dem. Hvis de forskellige Videnskaber, der befinder sig paa forskellige
Trin fra Kaos til Identitet, ikke havde en saadan fælles Grundform, kunde vi
ikke opstille selve Begrebet Videnskab som fælles for dem alle. Vi kan nu ikke
blot belyse den ene Videnskab gennem Analogi med andre
% --- p. 41 ---
\setcounter{footnote}{0}%
\opage{41}Videnskaber, og en Del af en Videnskab gennem Analogi med andre Dele,
men ogsaa den førvidenskabelige Tænken gennem Analogi med den videnskabelige,
og vi kan indenfor Psykologien belyse de elementære sjælelige Funktioner ved
Analogi med de mere udviklede. Saadanne Analogier er ikke blotte Fiktioner, som
\textsc{Vaihinger} mente i sin »Philosophie des Als Ob«%
% footnote 1 on p. 41
\footnote[1]{En interessant Fremstilling og Kritik af »Philosophie des Als Ob«
har \textsc{Jørgen Fr. Jørgensen} givet i Tidsskriftet »Vor Tid«. 1914---15. p.
227-249.}. I Kraft af den Kontinuitet, der gennem Syntesens Grundform gør sig
gældende i alt Aandsliv, og særlig indenfor det intellektuelle Liv, bliver
Analogier som de nævnte mulige; naturligvis er deres logiske Værdi højst
forskellig. Det er f. Ex. ingen Fiktion, naar en Tangent kaldes en Sekant, hvis
to Skæringspunkter med en Cirkel i højeste Grad nærmer sig til hinanden; thi
Syntesen kan nærme sine Elementer til hverandre i saare mange Grader. Ifølge en
moderne Retning i Geometrien berører Tangenten endog altid Cirkelen i mere end
ét Punkt, saa at Forskellen mellem Sekant og Tangent tilsidst falder bort.

7. Endog mellem de to yderste Grænser for vor Erkendelse, den kaotiske
Forskelsrække og den absolute Identitetsrække, er der Analogi tilstede. Den
eneste Tankeoperation, de gør mulig, er --- foruden unyttige Substitutioner ---
Tælling. Man kan tælle, hvor mange absolute Forskelligheder eller absolute
Identiteter, man kender. Som vi senere vil se, tager Aritmetiken netop fat der,
hvor den rene Logik ender. En anden Analogi mellem dem ligger deri, at ingen af
de to Rækker kan være umiddelbart given. Kun gennem streng kritisk Prøvelse,
der ikke vil kunne godtgøres at være udtømmende, vilde det kunne vises, at der
slet ingen Lighedspunkter eller slet ingen For%
% --- p. 42 ---
\setcounter{footnote}{0}%
\opage{42}skelligheder var indenfor en Række af Emner. De to Rækker betegner
Grænser for vor Erkendelse, der ikke kan bevises at være tilstede i
nogetsomhelst enkelt Tilfælde; kun gennem vilkaarlige Paastande vil deres
Gyldighed kunne hævdes. Men disse Rækker tjener til indirekte at karakterisere
den menneskelige Erkendelse gennem Paavisning af, hvad der vilde gøre den
umulig fra først af eller være en Afslutning uden mulig Fortsættelse.

Man kunde sige, at Erkendelsesteorien kun har til Opgave at drøfte de
Principer, der gør Erkendelse mulig, og at den derfor allerførst maa opstille
det Princip, der overhovedet ligger til Grund for al Dannelse af Begreber og
Domme og Slutninger, nemlig Identitetsprincipet. Erkendelse bestaar af Domme,
kunde man sige, og enhver Dom udsiger et Identitetsforhold (partielt eller
absolut). Begrebet Identitet synes derfor at burde aabne Kategoriernes Række.
Hertil maa svares, at Erkendelsesteorien bør undersøge alle Muligheder. Enhver
virkelig Erkendelse er stedse kun en Tilnærmelse til Identitet, og alle Grader
af Tilnærmelse maa undersøges. Erkendelsesteoriens egentlige Emne er
Brydningsgraderne mellem Identitet og Forskellighed. Selve Identiteten har kun
den Betydning, at den stiller fremskridende Reduktion af Forskellighederne som
den store Opgave. Selve Identiteten er jo en Relation. At hige efter at finde
Identitet er derfor at søge at stille en Relation istedetfor en anden. Altsaa
forudsætter Begrebet Identitet Begrebet Relation og det dermed korrelate Begreb
Syntese, og disse bliver de første Kategorier. Desuden kan Identitet kun
defineres som et Grænsetilfælde af Lighed eller af Kontinuitet, og disse to
Begreber gaar derfor --- naar vi gaar analytisk frem%
% footnote 1 on p. 42
\footnote[1]{Smlgn. \textit{Den menneskelige Tanke.} p. 177 f.} --- forud for
Identitetsbegrebet.

% --- p. 43 ---
\setcounter{footnote}{0}%
\opage{43}Naar det siges, at Erkendelse bestaar af Domme, maa det ikke glemmes,
at Dommens Subjekt stedse er taget ud af en hel Sammenhæng, og at denne Udtagen
kun bliver mulig ved den Forskel, der er mellem det og andre Emner, der er Led
i samme Sammenhæng. Kun ved en saadan Skelnen bliver et Emne til for os med en
saadan Selvstændighed, at der kan søges og findes nærmere Bestemmelser ved det.
Som man træffende har sagt, har ethvert Begreb Frynser eller Skaar, der viser,
at det egentlig er et Fragment. Først naar denne Udsondring er foretagen, kan
en Dom blive til. Hvis ethvert Emne druknede i eller var helt sammensmeltet med
det Medium, i hvilket det fra først af existerer, vilde der aldrig fældes nogen
Dom. En Forskelsrelation er baade logisk og psykologisk Forudsætningen for en
Doms Tilbliven.

Naturligvis forudsættes Identitetsprincipet ved al logisk Tænken, ogsaa ved
Opstilling af en Kategorilære og særlig ved Opstilling af Kategorierne Syntese
og Relation. Erkendelsesteorien maa naturligvis bygge paa de Principer, al
Videnskab bygger paa, og da det netop er disse Principer, Erkendelsesteorien
vil fremdrage og belyse, kan den forsaavidt siges at bygge paa sig selv eller
at bevæge sig i en Kreds. Man har undertiden taget Anledning heraf til at
godtgøre Umuligheden af en Erkendelsesteori. Saaledes Schelling, Lotze og i
nyeste Tid \textsc{Leonard Nelson} (dels i hans Skrift Ȇber die sogenannte
Erkenntnistheorie« 1908, hvis første Afsnit hedder »Die Unmöglichkeit der
Erkenntnistheorie«, dels i et Foredrag paa Filosofkongressen i Bologna 1911).
Nelson's Tankegang er denne: Kriteriet for Erkendelsens Sandhed maa enten være
Erkendelse eller ikke være Erkendelse. Hvis det er Erkendelse, maa denne
Erkendelse selv behøve et Kriterium. Og hvis det ikke er
% --- p. 44 ---
\setcounter{footnote}{0}%
\opage{44}Erkendelse, maa det dog være noget, der kan være Genstand for
Erkendelse, og saa behøver vi et Kriterium for, at denne Erkendelse er sand.
--- Denne Tankegang kender vi allerede fra Oldtiden. Den blev fra Skeptikernes
Side fremført mod Akademikernes Forsøg paa at opstille et
Virkelighedskriterium, idet man spurgte, ved hvilket Kriterium vi kan afgøre,
at det saaledes opstillede Virkelighedskriterium er det rette!%
% footnote 1 on p. 44
\footnote[1]{Se nærmere herom \textit{Relation som Kategori.}} Til Grund for
denne Skepsis saavel som for de moderne Indvendinger mod Erkendelsesteoriens
Mulighed ligger en Forvexling af Begreberne Kredsgang og Uafsluttelighed.
Erkendelsesteoriens Opgave maa stedse tages op paany. Det maa stedse paany
prøves, om de Principer, Erkendelsesteorien opstiller som betingende
Erkendelsens Gyldighed, falder sammen med de Principer, der ligger til Grund
ved selve Forsøget paa at opstille Erkendelsesteorien. Er der en Forskel her,
maa Undersøgelsen begynde forfra. En Afslutning af disse Undersøgelser kan ikke
tænkes. Det kan meget vel være, at en ny Undersøgelse af det menneskelige
Tankearbejde kan føre til Nødvendigheden af en ny Erkendelsesteori. Ogsaa før
Karneades, og i nyere Tid før Hume og Kant havde man en Erkendelsesteori; men
det viste sig, at den ikke var gennemført, og at dens fundamentale Betydning
ikke var godtgjort.

Erkendelsesteori bliver til ved Undersøgelse af selve Tankearbejdet, d. v. s.
af den Forandring, Anvendelsen af visse Principer frembringer i Henseende til
Klarhed og Sammenhæng i vore Forestillinger. Finder vi fremskridende Klarhed og
Sammenhæng, siger vi, at der erkendelsesteoretisk er sket et Fremskridt, selv
om vi maaske kommer til at stryge visse Antagelser, vi forhen lagde Vægt paa.
Jo flere
% --- p. 45 ---
\setcounter{footnote}{0}%
\opage{45}og jo dybere Forskelligheder mellem Emnerne, der kan fremdrages, og
jo mere det alligevel lykkes os at finde begrundet Sammenhæng mellem dem ved
Anvendelse af Identitetsprincipet, des mere bliver vi overbeviste om dette
Princips Gyldighed. Som vi har set, kan Identitetsprincipet allerede siges at
dæmre i de førvidenskabelige Analogier saavel som i de Rækkedannelser, der
endnu ikke danner Grundlag for Slutninger. Naar, som i en regelmæssig
varierende Forskelsrække (f. Ex. Farverækken), hvert Led ligger fast, saa at
ingen Ombytning kan finde Sted, er Identitetsprincipet kommet mere til
Gyldighed, end hvor Ordenen enten er aldeles vilkaarlig (som i den kaotiske
Forskelsrække), eller hvor Ordenen skyldes et helt udvortes Synspunkt, f. Ex.
et Hukommelsesskema (som i den ubestemt varierende Forskelsrække). Og
Videnskabernes Historie viser os mange Trin i denne Henseende, før fuld Klarhed
og Sikkerhed opnaas.

8. I de Analogier, der spiller saa stor en Rolle i det førvidenskabelige
Tankeliv, rører sig allerede den Energi, der senere kan føre til videnskabelig
Tænkning. De begrunder den Forventning, hvormed Menneskene møder overfor de
Emner, der fremstiller sig for dem, og som gør, at de ikke staar passive
overfor, hvad Iagttagelsen frembyder. Med en sangvinsk Tillid til alle
opdukkende Forestillinger, hvorfra de saa kommer, begynder Tankelivet sin
Udviklingsgang, ja, allerede forud for frie Forestillingers Optræden ytrer der
sig en Analogi til denne Tillid i det uvilkaarlige Hang til øjeblikkelig
motorisk Reaktion, hvorom Reflexer og Instinkthandlinger vidner. Livet begynder
med et Overskud af Energi og har en Tendens til at brede og udfolde sig. Før
der kan danne sig Erindring i Ordets snævrere Betydning (med Klarhed over, at
det, der fore%
% --- p. 46 ---
\setcounter{footnote}{0}%
\opage{46}stilles, hører hjemme i Fortiden), kan der røre sig Forventning, og
Forestillinger, der ikke er forbundne med nogen Forventning, men hvis Indhold
bestemt staar som forbigangent, fastholdes kun, naar en højere Udvikling er
naaet. Hume lagde med Rette Vægt paa Trangen til Expansion, til at løbe Linen
ud, fortsætte en Tankegang, man én Gang er slaaet ind paa. Men naar han af
denne Tendens vilde udlede Troen paa rationale og kausale Principer, oversaa
han, at disse Principer træder tydeligst frem, naar Forventninger skuffes eller
afvises, og naar vi saa nødes til at søge en Orden og Sammenhæng af anden Art
end den, vore uvilkaarlige Analogier og Forventninger bebudede. Rationalitet og
Kausalitet er Normer, der begrænser og inddæmmer vort uvilkaarlige Tankeløb,
der i sig selv har en Tendens til at forfølge visse Synsmaader ud over alle
Grænser. De Hypoteser, Videnskaben paa sin Vej har udformet, men har ladet
falde igen, vidner om, at det for en stor Del er gennem Skuffelser, vi arbejder
os frem til Sandheden. Dersom der ikke var kommet Skuffelser, var vi paa mange
Punkter ikke komne ud over det førvidenskabelige Udviklingstrin; men
Skuffelserne vilde ikke være komne, hvis den overskydende Energi ikke havde
fostret Forventninger, der uvilkaarlig prøvedes ved, at vi handlede efter dem.
Børns Trang til at gribe og bearbejde alle Ting fører dem til at gøre vigtige
Erfaringer. De experimenterer uvilkaarlig. Bevægelsestrangen træder
efterhaanden i Nysgerrighedens Tjeneste. I sin Selvbiografi fortæller den tyske
Kirurg Schleich, at han fem Aar gammel allerede havde faaet Ur, men knuste det
ved at banke paa det med en Sten, og da han saa fik Skænd derfor, sagde han:
»Jungens müssen doch wissen, wa da'inn is!« Denne Trang til at vide, hvad der
er »indeni«, rører sig lige fra
% --- p. 47 ---
\setcounter{footnote}{0}%
\opage{47}den barnlige Experimenteren til de højeste Spekulationer, uden at der
altid rører sig en Anelse om, hvor fin og hvor komplet den Tingenes Sammenhæng
er, som vi søger at trænge ind i. Hin Trang inddæmes efterhaanden, som det
erfares, at der raader bestemte Love for Tingene, og at det er gennem
Forstaaelse af dem, vi maaske kan trænge ind i, hvad vi kalder Naturens Indre.

I et aandfuldt Skrift (De la contingence des lois de la nature. 1875) udviklede
\textsc{Émile Boutroux} den Opfattelse, at Naturlovene væsentlig betyder Former
eller Baner for Udfoldelsen af frie Kræfter, Flodlejet for en Strøm, der
arbejder sig fremad. Et Skridt videre gik \textsc{Henri Bergson,} idet han
fordrede af Filosofien, at den skulde trænge ind bag Naturlovenes Mekanisme og
gøre sig til Et med den Kraft, den Livsudfoldelse, der stadig hæmmes og
begrænses af Naturordenen, og han fandt et Grundlag herfor i en umiddelbar
Intuition af vor egen indre Kraftudfoldelse (élan vital), i Analogi med hvilken
vi saa maa opfatte alle andre Existensformer end vor egen (Introduction à la
Métaphysique. 1903). --- Saadanne Tankeforsøg skyldes i Virkeligheden en
Omsætning i metafysisk Form af hvad Psykologi og Erkendelsesteori lærer. De er
gode Exempler paa, hvad vi i en senere Sammenhæng vil komme lidt nærmere ind
paa, at al Metafysik (Kosmologi) bygger paa Analogi, her paa Analogi med den
Maade, paa hvilken vi i Videnskabens Historie ser menneskelig Expansionstrang
og Utaalmodighed brydes med den faste Tingenes Orden, Videnskaben arbejder paa
at finde. Af særlig Interesse her er et Kapitel i \textsc{Ernst Mach's} Værk
»Erkenntnis und Irrtum« (1905). Kapitlet har til Overskrift »Sinn und Wert der
Naturgesetze« og søger at vise, at Naturlovene bliver til for os som
Begrænsninger, vi under Ledelse af Erfaringen
% --- p. 48 ---
\setcounter{footnote}{0}%
\opage{48}sætter for vore Forventninger, og at der gennem saadan fremskridende
Begrænsning foregaar en stedse nøjere Tillempelse af Tanken efter
Kendsgerninger. Allerede daglig Erfaring viser os en stadig Inddæmning af det
uvilkaarlige Tankeliv.

9. Ligesom Loverkendelsen er en Begrænsning af Friheden til at danne
Forestillingsforbindelser, saaledes er Begrebsdannelsen en Begrænsning af
Friheden til at forme de enkelte Forestillinger. Vi danner Begreber for at være
sikre paa, at en Forestilling, der dukker op paany i vor Tankegang, har samme
Indhold, som da den forhen fremtraadte; kun derved kan den indtræde som fast
Led i en Tankerække og staa i klare og bestemte Forhold til de andre Led i
Rækken. Gælder det blot om at fastholde, hvad der er ejendommeligt for et Emne
i en aldeles bestemt Situation, opnaas dette ved, hvad man kunde kalde et
konkret Individualbegreb. Saaledes er »Napoleon ved Waterloo« et konkret
Individualbegreb, naar vi har mærket os Napoleons hele Optræden ved nævnte
Lejlighed. Skal Begrebet gælde for alle Sider i et Emnes Væsen og for dets
Fremtræden under alle forskellige Situationer, f. Ex. for Napoleon efter hele
hans Karakter og Levnedsløb, faar vi et typisk Individualbegreb. Hver enkelt
Situation eller hver enkelt Handling tjener da kun som Exempel paa dette
Begreb. Og paa lignende Maade forholder det sig med Almenbegreber, der skal
gælde for en hel Gruppe af Emner. De »fælles« Egenskaber vil her ikke, ligesaa
lidt som ved de typiske Individualbegreber, være tilstrækkelige til at være
Indhold af et Begreb for sig. Det »Fælles« kan ikke tænkes for sig alene. Mere
eller mindre tydelig staar derimod et enkelt Emne som Exempel paa hele Rækken,
eller man
% --- p. 49 ---
\setcounter{footnote}{0}%
\opage{49}benytter et Skema eller et Ord, der, hvis det behøves, kan erstattes
ved et Exempel.

Mellem de forskellige Exempler indbyrdes vil der ikke være fuldstændig
Identitet i alle Henseender. Det er et eller andet Lighedspunkt (maaske
Identitetspunkt), der fører til at stille Emner sammen som Exempler paa samme
Begreb. En Indianer, der første Gang saa en Bog aabne sig, kaldte den for en
Østers. Det var allerede et Forsøg paa en Begrebsdannelse. Naar man tidlig
talte om Fødder baade hos Dyr og Mennesker, og senere udstrakte Begrebet til
Borde, Stole, Bjerge o. s. v., var det ogsaa Forsøg paa Begrebsdannelse. Men
heller ikke ved videnskabelig Begrebsdannelse dækker de enkelte Exempler
fuldstændig hverandre. Det er ved Analogiens Hjælp, at man indenfor et Begrebs
Omfang gaar fra Exempel til Exempel. Under Begrebet Keglesnitslinie hører
Cirkel, Ellipse, Parabel, Hyperbel, og Forholdet mellem disse Emner er ikke
Identitet, men Analogi. Kun fra et bestemt Synspunkt, som Definitionen af
Keglesnit angiver i Ord, kan de betragtes som identiske, forsaavidt som ethvert
af dem principielt kan bruges som Exempel ligesaa godt som de andre. Selve
Definitionen af Keglesnit (som en Figur, der fremkommer ved et plant Snit
gennem en cirkulær Kegle) kræver en Konstruktion, der fører til et eller andet
af de anførte Exempler. Det er Forholdet mellem Egenskaber eller Dele ved en
Række Emner, der gør disse Emner til Exempler paa samme Begreb. Ved Begrebet
Farve er der foreløbig kun Kvalitetslighed mellem de enkelte Exempler, og et
egentligt Begreb faas kun ved (ligesom ved Keglesnittene) at holde sig til,
hvorledes de opstaar. Hver enkelt Farvekvalitet forholder sig til en vis Grad
af Lysbrydning, som de andre Farver forholder sig til visse andre Grader af
% --- p. 50 ---
\setcounter{footnote}{0}%
\opage{50}Lysbrydning. For Kortheds Skyld (naar man ikke har Tid til at
fremdrage Exempler) kan Ord eller Tegn (Skemaer) udtrykke Begrebet. Saaledes er
Bogstaver i Aritmetiken Skema for Begrebet Størrelse, hvis Exempler haves i
Talrækkens Led. »Matematikere«, siger \textsc{Pascal} (De l'esprit
géométrique), »og alle, som gaar metodisk frem, giver kun Tingene Navne for at
forkorte Tankeløbet«. Ordet eller Skemaet skal give Anvisning til Konstruktion
af et Exempel. Pascal kræver, at man for Sikkerheds Skyld skal prøve at sætte
det, der defineres, istedetfor Definitionen. Men hertil maa bemærkes, at da
»det, som defineres«, fremtræder i forskellige Skikkelser, kan denne
Substitution aldrig blive absolut. Intet Exempel indeholder blot de Træk,
Definitionen angiver. Hvert enkelt Exempel kræver andre Exempler, og det er
Analogien imellem dem, der er Begrebets egentlige Indhold. Begrebets Betydning
bestaar i, at det kan sætte en Tankerække igang. \textsc{Goblot} siger derfor,
at et Begreb ikke er andet end Mulighed for en Række af Domme (une virtualité,
une possibilité indéfinie de jugements)%
% footnote 1 on p. 50
\footnote[1]{\textit{Traité de Logique.} 1918. p. 87.}. Han begynder derfor sin
Fremstilling af Logiken med Dommene og gaar derfra over til Begreber og
Slutninger. Allerede tidligere havde \textsc{Sigwart} hævdet Dommenes Prioritet
i Forhold til Begrebet. Men Sagen er vistnok mere indviklet, end de to
fortræffelige Logikere fremstiller den. Der er en stadig Vexelvirkning mellem
Begreb og Dom. Selve Begrebsdannelsen foregaar gennem Domme, men Dommene er kun
tydelige, naar de Begreber, der forbindes i dem, hvert for sig er bestemte.
Tænkningens Historie viser os, hvorledes Udgangspunktet snart tages i de
enkelte Begreber, snart i Dommen. I Erkendelsesteorien
% --- p. 51 ---
\setcounter{footnote}{0}%
\opage{51}fremtræder dette i Forholdet mellem Grundbegreber (Kategorier) og
Grundsætninger (Principer)%
% footnote 1 on p. 51
\footnote[1]{Smlgn. \textit{Den menneskelige Tanke.} § 51; 57; 106.}.

Der er noget i vor Tænken (og i vor Bevidsthed idethele), som aldrig fremtræder
for umiddelbar Iagttagelse, men som maa forudsættes eller, om man vil,
konstrueres, fordi det, som umiddelbar Iagttagelse opviser, ellers er
uforstaeligt. Saaledes er Syntese, Bevidsthedslivets Grundform og al
Erkendelses første Grundbegreb, ikke selv et iagttageligt Element. Heller ikke
vor Villen som saadan kan umiddelbart iagttages%
% footnote 2 on p. 51
\footnote[2]{\textit{Psykologi}\textsuperscript{6} p. 177 f. (Forkortet Udgave
p. 33 f.; 84 f.)}. Og saaledes gaar det med al indre og ydre Virksomhed; vi
iagttager dens Resultater og slutter til den Energi, Arbejdet skyldes. Selve
Tankevirksomheden iagttages ikke. Enhver enkelt Tanke er Led i et Tankeløb, og
vi slutter os til Tankens Energi af de Forandringer i Tankeløbet, der er
uafhængige af øjeblikkelige Iagttagelser (uden forsaavidt Forandringen bestaar
i, at saadanne Iagttagelser bringes i bestemt Sammenhæng med det forhen givne
Bevidsthedsindhold). Overalt iagttager vi kun Succession, ikke den Virksomhed,
der strækker sig gennem alle Successionens Led, og som vi ofte møjsommelig maa
efterspore. Her staar Hume's Kritik stadig ved Magt, paa Aandsvidenskabens, som
paa Naturvidenskabens Omraade. Hans voldsomme Isolation af de enkelte Led i
Tankeløbet ledede indirekte til den rette Opfattelse af Begrebet Virksomhed paa
alle Omraader. Men det er Leibniz, der direkte har angivet den rette Anvendelse
af Begreber som Kraft, Energi og Virken, idet han hævdede, at før der kan tales
om Kraft, maa der være paavist en lovmæssig Sammenhæng mellem skiftende Tilstande%
% footnote 3 on p. 51
\footnote[3]{\textit{Den nyere Filosofis Historie.}\textsuperscript{8} Il p. 136 f.}.

% --- p. 52 ---
\setcounter{footnote}{0}%
\opage{52}Den egentlige Tænken adskiller sig kun fra det sædvanlige
Forestillingsløb ved at være rettet mod Betingelserne for at løse en Opgave,
mod Midler og Veje til at naa et Maal. Den bestaar i en Analyse af den
Sammenhæng, indenfor hvilken Opgaven eller Maalet antages at have sin Plads, og
søger saa at finde en Lov mellem de Elementer, Analysen fører til. Et simpelt
Exempel frembyder Orientering, naar man er faret vild. De forskellige Retninger
prøves og sammenlignes. Hvis man derimod vandrer uden bestemt Maal, slaar man
ind paa den første den bedste Vej. Allerede \textsc{Hobbes,} der saa ofte maa
være Syndebuk ved de Kritiker af »Associationspsykologien«, man ikke synes at
kunne blive ked af at gentage, har indset dette, naar han udtaler, at det er
gennem stadig Fastholden af et Maal (frequens ad finem respectio. De Corpore
XXV, 9), vi lærer at finde Orden og Sammenhæng.

10. Begreberne Analogi og Analogislutning er først med Klarhed fremdragen af
Aristoteles. Medens det for Platon, selv efterat han havde haft sin sidste
Holmgang med Eleaterne, stod som det store Maal at føre al Erkendelse tilbage
til én Idé, er Aristoteles overbevist om, at vi ender i en Flerhed af Begreber
(Kategorier), mellem hvilke der ikke gælder Identitet, men kun Analogi. Mellem
de forskellige Omraader, paa hvilke Tænkningen arbejder, kan der maaske ogsaa
paavises Analogier; saaledes mener Aristoteles, at Forholdet mellem Stof
(Mulighed) og Form (Virkelighed) kommer igen overalt, --- men selve dette
Forhold faar en forskellig Karakter paa de forskellige Omraader (d. v. s.
indenfor de forskellige Videnskaber).

Hvad Bevisførelsen angaar, skelner Aristoteles mellem Deduktion og Induktion.
Ved deduktiv Slutning forstaar han syllogistisk Henføren under et Slægtsbegreb
ved Hjælp
% --- p. 53 ---
\setcounter{footnote}{0}%
\opage{53}af et Artsbegreb som Mellemled. Hans Slutningslære er altsaa bygget
paa en forudsat Klassifikation. Ved Induktion forstaar han en fuldstændig
Opregning af alle Emner, der hører ind under et Begreb. (ἡ ἐπαγωγὴ διὰ πάντων.
ed. Bekker p. 68 b). Men saa gives der ogsaa, tilføjer han (p. 69 a), en
Slutning ud fra et enkelt Exempel, eller rettere fra Exempel til Exempel. Denne
Slutning gaar altsaa hverken som Deduktionen fra et Hele til en Del, eller som
Induktionen fra alle Dele til en Helhed, men den gaar fra en Del til en anden
Del (ὡς μέρος πρὸς μέρος). Den adskiller sig fra Induktionen ved, at ikke alle
Dele, men kun en enkelt Del foreligger. Aristoteles' Benævnelse for
Analogislutning er Paradigma (παράδειγμα), der betyder Exempel, men ogsaa kan
betyde Forbillede. Det enkelte Emne kan ved Analogislutning siges at staa som
et Forbillede, en Type, i hvis Lys en hel Kreds af Emner ses, idet de frembyder
lignende Forhold mellem Egenskaber eller Dele indbyrdes. --- Selve Ordet
ἀναλογία bruger Aristoteles mest i Betydning af matematisk Proposition,
Forholdsidentitet (ἰσότης λόγων), medens han om kvalitativ Forholdslighed
hyppig bruger det ubestemte Udtryk τὸ ἀνάλογον.

I nyere Tid er man mere og mere klar over, at der ingen principiel Forskel er
paa Induktion og Analogi, da man paa Erfaringens Grund aldrig kan være sikker
paa at have alle mulige Emner af en vis Art givne%
% footnote 1 on p. 53
\footnote[1]{Undtagen naturligvis i saadanne Tilfælde, hvor det i selve
Begrebets Definition ligger, at der maa være et vist Antal Emner. Smlgn.
\textit{Den menneskelige Tanke.} p. 172 f.}. Den ufuldstændige Induktion
adskiller sig da kun fra Analogislutningen ved at have flere Exempler at gaa ud
fra. Naturligvis er det af Betydning at have saa mange Exempler som muligt at
gaa ud fra. Men for den moderne Metode%
% --- p. 54 ---
\setcounter{footnote}{0}%
\opage{54}lære er Hovedsagen den, om der ved Analyse af de foreliggende Emner
(et eller flere) kan findes en Sætning, dér foreløbig kan opstilles for saa at
prøves ved, at de Slutninger, der drages fra den, finder Bekræftelse i nye
Erfaringer. Denne Metode, der først i stor Stil blev anvendt af Galilei og
Newton, er karakteristisk for moderne Naturvidenskab. Analogi og Induktion
giver, som \textsc{Kant} udtrykker sig i sin Logik (ed. 1801, p. 210), blot
»logische Präsumtionen«. Han føjer underlig nok ikke til, at det dog ofte kan
have Betydning at prøve disse Formodninger.

\textsc{Kant} selv har klart set, at der gives en Art Slutten, som hverken er
Deduktion, Induktion eller Analogislutning. Han henviser her til det Princip,
ved hvilket matematisk Tænkning er forskellig fra formel logisk Tænkning. I
Kraft af dette Princip kan der konstrueres en identisk varierende
Forskelsrække, d. v. s. en Række, indenfor hvilken Forskellen mellem de paa
hinanden følgende Led stedse er den samme, --- hvor der altsaa er
Forskelsidentitet. Det simpleste Exempel herpaa er Talrækken, som opbygges ved,
hvad Kant kalder »Synthesis des Gleichartigen«. Foruden i Talrækken optræder
det samme Princip, vistnok under Indflydelse af Talrækken som Forbillede
(altsaa i Analogi med den), i Rækkerne af Tider, Grader og Steder, naar
(ligesom ved Tallet) alle specielle kvalitative Egenskaber udrenses fra disse
Begreber. Domme, der bygges paa saadanne Rækkedannelser (der kan kaldes
identiske Kvalitetsrækker), kalder Kant syntetiske Domme a priori. I Analogi
med disse paa Forskelsidentitet byggede Rækker har Kant nu konstrueret et
Begreb om »mulig Erfaring« eller »Natur« som et Skema, indenfor hvilket ethvert
Emne, vi skal kunne tillægge Virkelighed, maa finde sin Plads.
% --- p. 55 ---
\setcounter{footnote}{0}%
\opage{55}Vi vil paa et senere Punkt faa Lejlighed til at drøfte denne
Konstruktions Betydning.

Medens vi er ved Omtale af Analogibegrebets Forhistorie, maa der henvises til,
at i filosofiske Forsøg paa at naa til en Verdensanskuelse, en afsluttende
Erkendelse staar Striden i nyere Tid mellem dem, der vil naa til Afslutning
gennem Slutninger, altsaa ved Identitetsprincipets Hjælp, og dem, der hævder,
at Analogi er det sidste Tankemiddel, naar det gælder Afslutning. Det er
Modsætningen mellem Platon og Aristoteles, der gentager sig i nye Former. Paa
Platon's Side staar Spinoza og Hegel, paa Aristoteles' Side Leibniz,
Schopenhauer, Lotze og i nyeste Tid Bergson. Af \textsc{Søren Kierkegaard's} Papirer%
% footnote 1 on p. 55
\footnote[1]{\textit{Papirer.} Udg. af \textit{P. A. Heiberg} og \textit{V}.
\textit{Kuhr.} V. p. 29.} ser man, at han har været opmærksom paa Modsætningen
mellem Identitet og Analogi paa de afgørende Punkter, i god Overensstemmelse
med den kritiske Retning af hans Erkendelsesteori. Han var bleven klar over, at
systematisk Afslutning var umulig ved Hjælp af Identitetsprincipet, og da
Induktion og Analogi ingen streng Slutten tillader, kom han til det Resultat,
at der stedse sluttes af (ligesom der begyndes) med et Spring, og det var ham
sikkert en Tilfredsstillelse at faa Anvendelse her for et af sine
Yndlingsbegreber. Men Problemet er endnu mere indviklet, end han saa, og netop
Springet stiller det paany: Springet foregaar i en Kurve og tilvejebringer
altsaa en ny Kontinuitet, skønt i et andet Plan.

11. Naar der foreligger en Mangfoldighed af Emner for Tænkningen, vil dennes
første Skridt naturlig være at ordne dem i Rækker, hvis Led er saa nøje
forbundne som muligt%
% footnote 2 on p. 55
\footnote[2]{I det følgende gør jeg Brug af det Forsøg paa at opstille og ordne
de vigtigste Former for Rækkedannelser som jeg har gjort i \textit{Den
menneskelige Tanke.} p. 162---173}.
% --- p. 56 ---
\setcounter{footnote}{0}%
\opage{56}Det uheldigste Forhold vil være det, at der mellem alle Emnerne
indbyrdes bestaar lige store Forskelligheder, saa at de kan stilles i
hvilkensomhelst Orden, det skal være. Som allerede ovenfor vist, vilde det være
urigtigt at tro, at saadanne kaotiske Forskelsrækker var ejendommelige for den
primitive Bevidsthed; denne er netop karakteriseret ved en uvilkaarlig Tydning,
der ligger til Grund for den Eftertanke, den formaar at anstille. Først for en
kritisk Undersøgelse, der ikke er tilfreds med nogen foreliggende Tydning, kan
det se ud, som om det givne var et Kaos. Som vi har set, vil Virkeligheden af
en kaotisk Forskelsrække ikke kunne bevises. Men den foreløbige Antagelse af et
Kaos kan indeholde en Opfordring til nærmere Undersøgelse, og her tilbyder
Analogien et Udgangspunkt. Hvad der skulde være fælles for Leddene i en saadan
Række, vilde være netop den Omstændighed, at de er Elementer i et Kaos; det
gælder for \textit{A,} for \textit{B,} for \textit{C,} hele Rækken igennem.
Heri ligger Mulighed for en Slutning i den aristoteliske anden Figur, en af de
Maader, paa hvilke en Analogi kan formuleres. En positiv Slutning kan ikke
udledes af denne Figur, mens der kan være Anledning til at undersøge, om der
dog ikke skulde gives andre Lighedspunkter mellem \textit{A, B} og \textit{C}
end netop det ene, tarvelige at høre til et Kaos. Anden Figur har overhovedet
den Betydning at stille en Opgave. Vi ser da allerede i dette simple Exempel
Analogiens Betydning for Tænkningen, som ansporende, æggende og vækkende.
Analogien spørger, men om den svarer, er et helt andet Spørgsmaal. Og selve den
ansporende Virken forudsætter, at Participationen er opløst, thi ved den gøres
jo netop Lighed paa et enkelt Punkt til Lighed, hvad Helheden angaar.

Meget videre kommer vi ikke, naar Leddene ligger fast,
% --- p. 57 ---
\setcounter{footnote}{0}%
\opage{57}ikke uden videre kan bytte Plads, som naar Rækkefølgen bestemmes rent
kronologisk. Det skulde da være, at der netop rejstes Spørgsmaal om, hvorfor
Emnerne erfaredes eller erindredes netop i denne bestemte Orden. Ogsaa dette
forudsætter en vaagen Eftertanke. Et Barn, der lærer om den mosaiske
Skabelseshistorie, kunde f. Ex. spørge, hvorfor Planter blev skabt før Dyr,
eller hvorledes Lyset kunde blive skabt før Sol og Maane.

Hvis vi har en regelmæssig varierende Forskelsrække, som f. Ex. Farverækken,
ligger ikke blot Leddene fast, men der ytrer sig ogsaa en vis Kontinuitet i
Rækkefølgen, idet Kvalitetsligheden med Nabofarven stedse er større end den med
en fjernere liggende Farve. Og der kan være analoge Forhold indenfor Rækken,
idet f. Ex. Violet kan siges at forholde sig til Blaat, som Orange til Rødt.
Men især kan det Spørgsmaal komme frem, hvorledes denne bestemte kontinuerlige
Orden er at forklare. Goethe mente at kunne finde en Forklaring ved at holde
sig til de kvalitative Analogier. Men trods hans Protest er Videnskaben gaaet
ud over dette Kvalitetsstandpunkt ved at paavise en Analogi mellem
Lysbrydningens Grader og Kvalitetsrækkens Led. Istedetfor Analogien indenfor
Rækken, faar vi Analogi mellem to Rækker, af hvilke den ene er kvantitativ, den
anden kvalitativ. Det vil vise sig i det følgende, at en Analogislutning i
denne Form er af afgørende Betydning for den moderne Tids Opfattelse af Tilværelsen.

Hvis en Række er symmetrisk, d. v. s. kan læses i omvendt Retning, uden at de
enkelte Led byttes om eller noget Led udskydes, er der en ny Mulighed for
Analogi indenfor Rækken. Hvis (for at bruge de Morgan's Exempel) A er i Strid
med \textit{B,} og \textit{B} med \textit{C,} er der Analogi mellem \textit{A}
og \textit{C} med Hensyn til deres Forhold til \textit{B.} Analogislut%
% --- p. 58 ---
\setcounter{footnote}{0}%
\opage{58}ningen kan formuleres i anden Figur: \textit{A} er Fjende af
\textit{B, C} er Fjende af \textit{B.} Ogsaa tredie Figur kan bruges:
\textit{B} er Fjende af \textit{A} og er Fjende af \textit{C.} Men Analogien
fører her til at spørge, hvorledes det nærmere Fjendskab er mellem \textit{A}
og \textit{C.} Det kan være, at det fælles Forhold gør \textit{A} og \textit{C}
til Venner. Men det kan ogsaa være, at det fælles Fjendskab ikke er nok til at
overvinde deres gensidige Ligegyldighed eller det Fjendskab, der af andre
Grunde bestaar imellem dem.

Ved en fremskridende Forskelsrække (d. v. s. en saadan, i hvilken hvert Led
forholder sig til sin Forgænger, som dets Efterfølger forholder sig til det
selv) og ved partielle Identitetsrækker (d. v. s. saadanne, i hvilke hvert Led
er Prædikat til det foregaaende og Subjekt til det følgende, eller Følge af det
foregaaende og Grund til det følgende) er ikke blot lignende Analogislutninger
som de allerede fremdragne mulige, men der kan ogsaa drages Slutninger ved
simpel Udskydning af Mellemled. En fremskridende Forskelsrække dannes f. Ex. af
Tallene, Tiderne, Graderne og Stederne, som vi ovenfor kom til at omtale fra et
andet Synspunkt. En partiel Identitetsrække dannes, naar Prædikater selv bliver
Subjekter for nye Prædikater og Følgerne Grunde til nye Følger. De to Arter af
Rækker er karakteristiske for de formale Videnskaber (Logik og Matematik). Men,
som vi allerede har set, faar saadanne Rækker ogsaa Betydning for rent
kvalitative Forskelsrækker. Kan det vises, at en fremskridende Forskelsrække
eller en partiel Identitetsrække indenfor vore Emners Kreds Punkt for Punkt
forløber parallelt med en fast ordnet Kvalitetsrække, saa at der til et bestemt
Led i den ene Række svarer et bestemt Led i den anden, saa kan der fra en
Overgang indenfor den ene Række sluttes til en Overgang indenfor
% --- p. 59 ---
\setcounter{footnote}{0}%
\opage{59}den anden Række. Al exakt Erfaringsvidenskab beror paa saadan
Sammenstilling af kvantitative og kvalitative Rækker. Hvor saadan
Sammenstilling er mulig, kan man betragte Kvalitetsrækkerne, som om de var rent
logiske eller matematiske Rækker. Man kan f. Ex. betragte Farverækken, som om
Forskellen mellem Farverne blot bestod i Forskellen mellem Lysbrydningens
Grader, og man kan betragte uundgaaelige kvalitative Successioner, som om de
blot var logiske Overgange fra Grund til Følge (i en partiel Identitetsrække)
eller matematiske Overgange fra en Lysbrydningsgrad til en anden. Men det er og
bliver dog Analogier, hvorpaa der her bygges, selv om der her er sket det store
Fremskridt i Forhold til de ovenfor omtalte Analogislutninger, at der ikke blot
gives Anledning til Spørgsmaal, men begrundes bestemte Forventninger og
Formodninger, altsaa dannes Hypoteser, der kan søges bekræftede enten ved, at
der opsøges nye Led i den kvalitative Række, eller ved, at den kvantitative
Række udformes med større Nøjagtighed end før, maaske endog saaledes, at der
konstrueres en ny kvantitativ Række. Medens visse Uregelmæssigheder i Uranus'
Bane kunde forklares ud fra Newton's Gravitationsteori, kan Uregelmæssighederne
ved Merkur's Bane ikke forklares ud fra denne Teori, og \textsc{Einstein} har
da forsøgt en Forklaring ud fra den elektromagnetiske Teori, af hvilken saa
selve Gravitationsteorien skulde være at udlede%
% footnote 1 on p. 59
\footnote[1]{\textsc{Einstein:} \textit{Über die spezielle und die allgemeine
Relativitätstheorie}\textsuperscript{6} \textit{,} p. 70.---
\textsc{Eddington:} \textit{Space, Time and Gravitation.} 1920. p. 123---125.}.

Om de to Rækker vil vedblive at forløbe parallelt, kan ikke vides med fuld
Sikkerhed. Nok saa lang en Fortsættelse af Kvalitetsrækken borger ikke for, at
der jo kan opstaa nye Emner, der ikke ordner sig efter den i den kvan%
% --- p. 60 ---
\setcounter{footnote}{0}%
\opage{60}titative Række angivne Regel. Allerede \textsc{Leibniz} har set, at
en fuldstændig Verifikation er en Umulighed. »En Hypoteses Held«, siger han%
% footnote 1 on p. 60
\footnote[1]{\textit{Nouveaux Essais.} IV, 17, 5. (Op. phil. ed. Erdtmann. p.
397.)}, »beviser ikke dens Sandhed«. Til fuldstændig Stadfæstelse af en
Hypotese vilde der høre, hvad Leibniz kalder en fuldstændig Omvending (un
parfait retour) af Tankegangen, d. v. s. en Godtgørelse af, at man ud fra
Kvalitetsrækken med Nødvendighed kom tilbage til denne bestemte Kvantitetsrække
og ikke til nogen anden Række. Man maatte kunne vende Slutningen om og gøre
dens tidligere Præmisser til Konklusion. Der kan jo drages en objektiv rigtig
Slutning af gale Præmisser, og Videnskabens første (eller sidste)
Forudsætninger er derfor vanskelige at fastslaa.

Et Identitetsforhold er der i hvert Tilfælde ikke tilstede her, skønt en
dogmatisk Tendens atter og atter fører til at antage et saadant for at kunne
hvile ud i et Første (eller Sidste). Kvalitet kan nu engang ikke udledes af
Kvantitet, selv om Studiet af Kvalitetsrækken kan føre til at opstille
forholdsmæssig varierende Kvantitetsrækker. Kvantitet er selv en Kvalitet, der
(som ved Tal, Tid, Grad, Sted) egner sig til Opstilling af fremskridende
Forskelsrækker. Videnskaben gaar overalt, hvor det kan lade sig gøre, ud paa at
danne saadanne Rækker for saa ved deres Hjælp at finde Lovmæssighed indenfor
andre Rækker, hvis kvalitative Karakter ikke paa den Maade kan reduceres. Vi
vil senere komme til at omtale nogle særlige Exempler paa den store Analogi, al
moderne exakt Erfaringsvidenskab beror paa, ofte uden at den selv er sig det
bevidst. Vi vil ogsaa komme til nærmere at betragte Kant's store Fortjeneste af
Klarhed paa dette Punkt, saavel som den Betænkelighed, der endnu bliver tilbage
overfor hans Teori.
% --- p. 61 ---
\setcounter{footnote}{0}%
\opage{61}Foruden Kant er det især den berømte Fysiker Clerk Maxwelt (i
filosofisk Henseende en Discipel af Kant i andet Led), der har hævdet
Analogibegrebets Betydning i denne Sammenhæng. Til ham har Ernst Mach paa
væsentlige Punkter sluttet sig. ---

Til Slutning skal bemærkes, at selv hvor absolute Identitetsrækker kan dannes,
møder vi Analogien som Forudsætning. Det var det eleatiske Ideal at reducere
alt til Identitet. Kunde det naas, vilde alle Tilværelsens fundamentale Emner
kunne fremstilles i Rækken \textit{A = B = C = D} o. s. v. Men (som jeg
udførlig har vist i »Den menneskelige Tanke« p. 177---201) enhver absolut
Identitet mellem reale Emner forudsætter et bestemt Synspunkt, i Forhold til
hvilket de kan betragtes som identiske, altsaa substitueres for hverandre. Fra
det Synspunkt, som Loven for en Gruppe Emner angiver, kan disse Emner betragtes
som identiske, skønt de kan have mange andre Egenskaber end dem, Loven angaar.
Det er det samme Forhold som mellem Exempler paa samme Begreb.

\begin{center}\large III. Analogi mellem Erkendelsesfunktioner.\end{center}

12. Fra et psykologisk Synspunkt bliver det mere og mere umuligt at antage
skarpe Modsætninger mellem Sjælelivets forskellige Trin. Der viser sig en
stadig Kontinuitet mellem de forskellige Trin, og dermed bliver en Analogi
mellem de forskellige Trin, og derved igen mellem de forskellige sjælelige
Funktioner mulig, selv om vi ikke, som der allerede ovenfor var Lejlighed til
at indskærpe, uden videre kan identificere det højere Udviklingstrins
Funktioner med det laveres. Der er her en Sammenhæng, som trodser de
Distinktioner, man har villet hævde. Hvis f. Ex. den saa ofte gjorte Forskel
mellem Stof og Form i Sjæle%
% --- p. 62 ---
\setcounter{footnote}{0}%
\opage{62}livet kunde gennemføres, vilde der være en Grænse baade for Analogi
og for Kontinuitet. Men allerede Sansningen trodser denne Distinktion. Det ses
allerede, naar vi sammenligner Sansefornemmelserne med de tilsvarende
fysiologiske Processer (saavidt vi kender dem); hvad der fysiologisk er en
sammensat Proces, fremtræder i Sansefornemmelsen sammenfattet i et forholdsvis
enkelt Element. Vi sporer allerede her Syntesens Lov. Men selve dette Element
bestemmes fra første Færd dels ved sit Forhold til umiddelbart forudgaaende
Elementer, dels ved sit Forhold til samtidig optrædende Elementer. Og meget
hurtig væves der reproducerede Elementer sammen med det. En absolut enkelt og
selvstændig Fornemmelse, der skulde være et absolut passivt (ɔ: udefra bestemt)
Fænomen, kan ikke paavises. Kun tilnærmelsesvis kan en saadan fuldstændig
Isolation paavises. Allerede \textsc{Salomon Maimon} \textsc{%
% footnote 1 on p. 62
\footnote[1]{Smlgn. \textit{Den nyere Filosofis Historie}.\textsuperscript{8}
III, p. 122.}} indvendte mod Kant's Lære om Fornemmelserne som blot »Stof«, at
en absolut passiv Fornemmelse var at sammenligne med √2, altsaa kun
tilnærmelsesvis kunde paavises. Vi finder saaledes allerede i den simpleste
Sansning en Analogi baade til, hvad der paa højere Trin tydelig træder frem som
Syntese, og til den Skelnen, den gensidige Bestemmelse af Forskelligheder, der
er ejendommelig for Eftertanken.

En Overgang mellem Sansning og Eftertanke danner Erindrings- og
Fantasianskuelse. Ved Anskuen gør ikke blot Forskelsforhold, men ogsaa
Lighedsforhold sig gældende; de betinger den Sammensmeltning eller den
Organisation, enhver Anskuen bærer Præg af%
% footnote 2 on p. 62
\footnote[2]{Forskellen mellem Sammensmeltning og Organisation som to Arter af
sjælelig Helhedsdannelse havde jeg Anledning til at fremhæve i \textit{Den
store Humor.} p. 15---34.}. Erindring og Fan%
% --- p. 63 ---
\setcounter{footnote}{0}%
\opage{63}tasi frembyder, ligesom Sansningen, en Analogi til den Maade, paa
hvilken Eftertanken stræber at bringe Sammenhæng mellem de Emner, der er den
tilgængelige. De uvilkaarlige Forventninger, der er Erindringens og Fantasiens
første Form, danner Analogi til de Hypoteser, Eftertanken ad Analysens Vej
opstiller. --- En Mellemform mellem Anskuen og Eftertanke fremtræder i den
uvilkaarlige Ændring af et Anskuelsesbillede, ved hvilken det kan blive
artikuleret, skifte Retning og Tyngdepunkt, uden at Eftertanke behøver at
spille nogen Rolle med. Der vil især hos nogle Naturer (dem, der ovenfor er
kaldt de intuitive) være en Tendens til at fastholde Anskuen saalænge som
muligt og holde Eftertankens Analyse borte. Det er et Exempel paa den hyppige
Forvexling af Identitet og Analogi, naar mange Psykologer lægger Domme ind i
Sansningens og Anskuelsens uvilkaarlige Funktioner. Dommen bliver først
nødvendig, naar umiddelbar Anskuen ikke, end ikke gennem Artikulationer og
Forskydninger, kan rumme alle de Emner, der tildrager sig Opmærksomhed. Og
Dommen vil altid være et ufuldkomment Surrogat for Anskuelsen, da Eftertanken
neppe kan udskille alle i Anskuelsen indeholdte Elementer og bringe dem i den
nye Forbindelsesform, Dommen er. Derfor vil den enkelte Dom til sin Supplering
kræve nye Domme for derved at faa flere Decimaler til Bestemmelse af det
irrationelle Forhold mellem Anskuen og Eftertanke. --- En og samme Type
behersker, psykologisk set, Erkendelseslivets Udvikling gennem forskellige
Trin. Udførligere end det her var Stedet til, har jeg søgt at godtgøre dette i
\textit{Den menneskelige Tanke.} p. 34---54.

Det kan have sin Betydning at se, at der ogsaa, hvad andre Sider af Sjælelivet
angaar, gør sig Kontinuitet og en ved den betinget Analogi gældende. Svarende
til Modsæt%
% --- p. 64 ---
\setcounter{footnote}{0}%
\opage{64}ningen mellem Sansning paa den ene Side, Erindring, Fantasi og
Eftertanke paa den anden Side fremtræder paa Følelseslivets Omraade
Modsætningen mellem elementære og ideelle Følelser. Elementære Følelser hænger
sammen med Livsfornemmelser, der svarer til Forandringer i Blodomløb, Aandedræt
og Fordøjelse, eller med Fornemmelser, der svarer til Paavirkninger udefra.
Forestillinger spiller ingen Rolle, hverken Erindringer, Fantasier eller
Reflexioner. Allerede ved saadanne Følelser gør sig Kvalitetsforskelligheder
gældende, f. Ex. mellem Kraft- og Mathedsfølelse, mellem Lethedsfølelse og
Beklemthed, mellem Frihedsfølelse og Angstfølelse. Ved ideelle Følelser spiller
Forestillinger (Erindringer, Fantasier, Reflexioner) en større eller mindre
Rolle med. Derved faar Tilstandene en mere bestemt Karakter, en mere eller
mindre udpræget Retning. Den ubestemte Letheds- og Frihedsfølelse bliver til
Glæde ved det, der med Rette eller Urette staar som Aarsag til den; den
ubestemte Angstfølelse bliver til Frygt for noget, der staar som truende, o. s.
v. Men skønt vi her har forskellige Trin i Følelseslivets Udvikling for os, er
der en stadig Overgang mellem dem. Ideelle Følelser er aldrig uafhængige af
vore organiske Tilstande eller af de Sansefornemmelser, der i Øjeblikket
opdukker, men kan netop være prædisponerede ved disse. Fælles for alt, hvad vi
i Psykologien kalder Følelse, er Forskellen mellem Lyst og Ulyst; den gælder
for de ideelle som for de elementære Følelser og er kun Symptomer paa, at der
overhovedet er Følelse tilstede. Ingen Psykolog falder paa at ville aflede de
specielle Kvaliteter, der fremtræder i elementære og ideelle Følelser, af de
abstrakte Begreber Lyst og Ulyst. Derimod har det Interesse i denne Sammenhæng at be%
% --- p. 65 ---
\setcounter{footnote}{0}%
\opage{65}mærke Analogierne mellem elementære og ideelle Følelser, saavel som
den Maade, hvorpaa Overgangen sker fra hine til disse. At opfatte de udprægede
Kvaliteter, som først de ideelle Følelser frembyder, som oprindelige, vilde
stride mod, hvad genetisk Psykologi, særlig Undersøgelsen af Dyrets og det
spæde Barns Følelsesliv lærer.

I sin interessante Bog »Følelseslivets Grundelementer« (1920) har \textsc{Carl
Jørgensen} givet en Række gode Beskrivelser af forskellige Følelser. Man mærker
helt igennem den gode Iagttager. Men han lægger slet ingen Vægt paa de mere
ubestemte (og dog kvalitative) Følelser, jeg kalder elementære, og han
benægter, at Forestillinger spiller nogen Rolle ved Følelseslivets Udvikling
til mere udprægede Former. Og i sin Kritik af min Opfattelse overser han, 1) at
jeg allerede for de elementære Følelsers Vedkommende antager visse
Kvalitetsforskelligheder (der altsaa ikke skyldes Forestillingers Indflydelse),
2) at jeg efter min Terminologi ved Glæde forstaar Glæde ved noget (en »Glæde«,
der ikke var Glæde ved noget, vilde jeg regne under Magt-, Letheds- eller
Frihedsfølelse), 3) at jeg ikke antager, at det skyldes Forestillinger, at
Følelser overhovedet faar Kvalitet, men kun at der under Indflydelse af
Erindringer, Fantasier og Reflexioner opstaar, hvad jeg udtrykkelig har
betegnet som »\textit{nye} Følelseskvaliteter«. (Se den af Forf. benyttede
femte Udgave af min Psykologi p. 304; 305; 340). Jeg kan ikke erkende, at
Forfatteren har anvendt Omhu for at sætte sig ind i min Opfattelse; saa kunde
han ikke have døbt denne »Lyst- og Ulystteorien« og givet det Udseende af, at
jeg antager Lyst- og Ulystfølelse uden nogensomhelst speciellere Kvaliteter. ---

Paa Villieslivets Omraade (naar vi tager Ordet Villie i videste Betydning)
finder vi ogsaa en Række af Trin: Trang, Drift, Overvejelse, Forsæt,
Beslutning. Analogien
% --- p. 66 ---
\setcounter{footnote}{0}%
\opage{66}mellem dem beror paa de ved dem bestemte sjælelige Tilstandes aktive
og koncentrerede Karakter. (Herom henviser jeg til min Afhandling »Begrebet
Villie« i tredie Række af mine »Mindre Arbejder«). ---

Paa alle Sjælelivets Omraader gør der sig mellem de forskellige Trin
kvalitative Forskelle gældende. Det ene Trin kan ikke reduceres til det andet,
men der findes Analogier mellem dem, der vidner om en og samme Types Herredømme.

13. Et er det, at det indenfor Psykologien kan lykkes at finde Kontinuitet
mellem de forskellige Udviklingstrin trods deres kvalitative Forskellighed, og
at der derfor kan paavises Analogier mellem mere elementære og mere udformede
sjælelige Funktioner, et andet er, om der ikke paa dette Punkt er en principiel
Forskel mellem den psykologiske og den erkendelses-teoretiske Betragtning af
Tænkning og Erkendelse. Efter nogles Mening er det helt andre Principer, der
maa anlægges, naar man spørger om vor Erkendelses Gyldighed, end naar man
spørger om dens faktiske Beskaffenhed og Udvikling. Jeg skal nu ikke her komme
tilbage til det almindelige Spørgsmaal om Forholdet mellem Psykologi og
Erkendelsesteori, som jeg har drøftet i »Totalitet som Kategori« (p. 3---20),
men skal blot gøre nogle Bemærkninger om Forholdet mellem gyldig og faktisk
Tænkning. Det er klart, mener jeg, at da det er Mennesker, der tænker de
gyldige Tanker, maa de gyldige Tanker ogsaa være faktiske Tanker, Tanker, der
virkelig tænkes, medens omvendt Tanker, der faktisk tænkes, ikke altid er
gyldige Tanker. Hermed er Forholdet mellem Psykologi og Erkendelsesteori givet.
Saa højt end videnskabelig Tænken hæver sig, saa mange nye Tankemidler og
Tankeformer den end udvikler for at løse sin Opgave,
% --- p. 67 ---
\setcounter{footnote}{0}%
\opage{67}saa maa den dog være psykologisk mulig, og Psykologien maa, med
Benyttelse af Videnskabernes Historie, undersøge dens Mulighed og dens faktiske
Udvikling. En Anekdote, Kant etsteds fortæller i en anden Anledning, udtrykker
Forholdet i et godt Billede. En Indianer saa ved et Besøg hos en Europæer
Proppen springe af en Ølflaske og Skummet bruse ud, og Europæerne morede sig
over hans store Forbavselse derover. Men da sagde han: »Jeg undrer mig ikke
over, at det kommer ud, men over, hvorledes I har faaet det derind!« Saaledes
er netop Psykologiens Stilling overfor Videnskab, saavel som overfor Kunst og
Religion. Psykologien kan ligesaa lidt frembringe Videnskab, som den kan
frembringe Kunst eller Religion. Men den kan og maa stille sig den Opgave at
undersøge, hvorledes disse Former for aandelig Kultur er mulige. Historien
viser ogsaa, at Videnskab opstaar af den uvilkaarlige Trang til Orientering i
Verden. Hvis vi er i Tvivl om, hvorvidt vi staar overfor noget virkeligt eller
ikke, eller om hvilken Vej vi skal slaa ind paa, naar vi mærker, at vi er paa
Vildspor, saa faar vi allerede Brug for formal Videnskab, d. v. s. for
Slutninger, Maalinger og Beregninger, og for real Videnskab, d. v. s. for
Udfindelse af Aarsagsforhold, Udviklingsprocesser og Helhedsdannelser, hvad de
foreliggende Emner angaar. En af de vigtigste Forskelle mellem førvidenskabelig
og videnskabelig Tænken er, som vi allerede har haft Lejlighed til at omtale,
at ved hin spiller uvilkaarlig Expansion en væsentlig Rolle; man spørger (hvis
man overhovedet spørger) ikke om »Hvorfor?« men »Hvorfor ikke?« --- og lader
staa til, saalænge det gaar, medens den videnskabelige Tænken fødes af Tvivlen
og prøver hvert enkelt Skridt. Naturligvis anvender et videnskabeligt Geni
uvilkaarlig den strenge
% --- p. 68 ---
\setcounter{footnote}{0}%
\opage{68}Tænknings Principer; i Geniet er, ligesom i Instinktet, Trang og Evne
uadskillelig forbundne.

Allerede Kant gjorde opmærksom paa, at »selv den jævne Menneskeforstand« tænker
efter visse Principer, og i nyeste Tid er det især af \textsc{Émile Meyerson,}
støttet til Videnskabernes Historie, energisk gjort gældende, at al Videnskab
udvikler sig af den sunde Menneskeforstand (sens commun),%
% footnote 1 on p. 68
\footnote[1]{Det danske Udtryk »sund Menneskeforstand« er bedre end det tyske
»der gemeine Verstand«, det franske »sens commun« og det engelske »common
sense«, idet der i disse sidste Benævnelser ligger noget gængse, fælles og
traditionelt, der ikke altid er »sundt«. \textsc{Léon Brunschwicg} har
træffende skelnet mellem »bon sens« som »fonetion de discernement« og »sens
commun« som »tradition de consentement universel« og finder denne Forskel
allerede gjort hos Montaigne og Descartes.
(\textit{L}\textit{'}\textit{expérience humaine et la physique.} 1922, p. 574
f.).} idet den Skridt for Skridt berigtiger dennes Principer, hvor det behøves.
Der er ikke Identitet, men Analogi mellem Videnskab og sund Menneskeforstand.
Som Meyerson udtrykker det: »Den sunde Menneskeforstands Begreber er blevne
frembragte ved en Proces, der vel er ubevidst, men iøvrigt strengt analog med
den Proces, ved hvilken de videnskabelige Teorier formes« (Identité et
Réalité\textsuperscript{2} p. 393). Den væsentlige Grund til, at Videnskaben
ikke kan blive staaende ved den sunde Menneskeforstand, finder Meyerson deri,
at skønt denne sidste kan føres til at søge og maaske finde Lovmæssighed i
Oplevelserne, fyldestgør den ikke Trangen til at sætte noget varigt, noget
konstant istedetfor de mere eller mindre flygtige Oplevelser, og han forklarer
derved Atomlærens stadige Genopdukken i Videnskabens Historie. (Anf. Skr. p.
420---423). Men det vilde maaske være nok saa rigtigt at sige, at Videnskaben
ikke finder den Lovmæssighed, ved hvilken den sunde Menneskeforstand stanser,
tilstrækkelig speciel, nøjagtig og omfattende, og at Atomteorien netop stedse
dukker op paany,
% --- p. 69 ---
\setcounter{footnote}{0}%
\opage{69}fordi kun denne Teori muliggør en aldeles nøjagtig Forestilling om
den mest specielle Lovmæssighed. --- Til dette Differenspunkt mellem den
udmærkede franske Erkendelsesteoretiker og mig vil jeg siden faa Lejlighed til
at vende tilbage. ---

Ved følgende Exempel kan maaske Analogien mellem Psykologi og Erkendelsesteori
betegnes. Psykologien gør opmærksom paa, at for at en Tidsbevidsthed skal
udvikle sig, maa der i Bevidstheden gøre sig en Modsætning gældende mellem
noget kontinuerlig bestaaende og noget skiftende. Hvis en Tilstand glemtes,
saasnart en anden Tilstand opstod, vilde Tidsbevidsthed ikke være mulig. En
simpel Fornemmelse af Forandring, hvad vi kan kalde en Tidsfornemmelse,
forudsætter en forholdsvis fast Baggrund, i Modsætning til hvilken Forandringen
kan mærkes. I simpleste Tilfælde vil denne konstante Baggrund dannes af
Livsfølelsen (den til den givne organiske Tilstand svarende Følelse), som kan
holde sig, medens forskellige mere periferiske Oplevelser afløser hverandre. En
egentlig Tidsforestilling (som forskellig fra den blotte
Forandringsfornemmelse) forudsætter, at visse Oplevelser kommer igen efter et
Mellemrum, saa at en Rytme kan opfattes. Her vil en Trang eller en Drift kunne
danne den faste Baggrund, og en Rytme af Savn, Stræben, Opnaaelse (eller
Skuffelse) kan erfares. Det er denne praktiske Trehed, der fra først af bærer
Tidsopfattelsen. Praktisk gaar Interessen for Tidsrytme forud for Interessen
for Tidsmaaling, hvad Religionens og Magiens Historie vidner om.%
% footnote 1 on p. 69
\footnote[1]{\textsc{Hubert:} \textit{La représentation du temps dans la
religion et dans la magie} (Hubert et Mauss: Mélanges d'histoire des
religions), p. 195.} En objektiv Tidsbevidsthed udvikler sig ved Iagttagelsen
af regelmæssig skiftende Naturfænomener (Dag og Nat, Aarstiderne o. s. v.)
% --- p. 70 ---
\setcounter{footnote}{0}%
\opage{70}og danner, ved de nøjagtige Maalinger, Astronomien gør mulig,
Overgangen til den videnskabelige Tidsopfattelse. De Forandringer, denne
Tidsopfattelse har gennemgaaet i den nyere Tid, har jeg havt Lejlighed til at
omtale i Afhandlingen »Relation som Kategori« (p. 62---77). Det er især
Spørgsmaalet om Udgangspunktet for Tidsmaalingen, der bliver afgørende. I
Konsekvens af den kopernikanske Teori hævdede Bruno, at hver Stjerne maatte
have sin Tid, og i Konsekvens af den elektromagnetiske Lysteori har Einstein
hævdet, at Tidsmaalingen beror paa, om det Koordinatsystem, der lægges til
Grund, forholdsvis er i Hvile eller i Bevægelse. Selv ved den mest exakte og
specielle gennemførte videnskabelige Teori om Tid og Tidsmaaling falder
Analogien til den primitive Tidsopfattelse ikke bort, ligesom den gennem
Aarhundreders stadige Arbejde har udviklet sig af den. ---

Paa lignende Maade som med Tidsopfattelsen forholder det sig med Begreberne
Tal, Grad og Sted. Ogsaa de har udviklet sig til den Form, i hvilken
Videnskaben nu arbejder med dem, af kvalitative Iagttagelser, og Analogien med
Begyndelsesstadierne er ikke falden bort. Paa alle Stadier virker Principerne
for Orientering i Emnernes Verden.

14. Videnskabernes Historie viser os mærkelige Exempler paa, hvorledes visse
Grundtanker atter og atter dukker op paany, skønt i saa forskellige Former og
med saa forskellig Begrundelse, at man ikke sjælden er tilbøjelig til at
benægte baade den historiske Sammenhæng, som dog faktisk er tilstede, og den
blivende Betydning, de kan have for fortsat Forsken. Der gør sig trods alt en
Kontinuitet gældende i den videnskabelige Udviklingsgang og i Sammenhæng dermed
interessante Analogier, som tyder paa en vis
% --- p. 71 ---
\setcounter{footnote}{0}%
\opage{71}Grundtype for menneskelig Tænkning og for de Krav, den skal fyldestgøre.

a. Atomteoriens Historie viser os, hvorledes der gennem Tiderne gør sig to
tilsyneladende modsatte Tendenser gældende i menneskelig Tænkning, paa den ene
Side Trangen til at finde Noget, som bestaar trods alle Forandringer, og paa
den anden Side Trangen til at forudsætte sondrede Udgangspunkter for de
Forandringer, Erfaringen godtgør. Atomteori opstaar ved Samvirken af disse to
Tendenser. Den skyldes den Erfaring, at man ikke naar til virkelig Forstaaelse
ved at blive staaende ved de kvalitative Forandringer, Sanserne viser os.
Demokritos erklærede derfor, at Sansekvaliteterne kun havde tilsyneladende,
ikke virkelig Existens. Naar den ene Kvalitet afløser den anden, synes der at
opstaa noget og gaa noget tabt; men naar vi anvender »et finere Tankeorgan« end
Sansningen, naaer vi til en »ægte Erkendelse«, som viser os, at det forholder
sig helt anderledes: Intet opstaar eller forgaar, men hvad der sker, er, at
visse Smaadele grupperer sig paa en anden Maade end før. Disse Smaadele er
ensartede, thi ellers vilde de ikke kunne virke paa hverandre. Det er i Tankens
Navn, Demokritos hævder sin Atomlære, og han blev derfor af Sextus Empiricus
stillet sammen med Platon, idet begge hævdede, at kun det, der kan tænkes (τὰ
νοητά), kan være sandt.%
% footnote 1 on p. 71
\footnote[1]{Om en Uklarhed i Demokrits Lære om Sansning og Tænkning se min
Afhandling om Demokrit og Platon. (Mindre Arbejder. III. p. 98).}

Historisk har Demokrits Atomlære øvet en stor Indflydelse. Hans Skrifter er,
paa nogle Fragmenter nær, gaaet tabt, et af de største Tab indenfor
Oltidslitteraturen. Men den blev i den Form, Epikuros havde givet den,
fremstillet paa glimrende Maade i Lucretius' berømte Læredigt;
% --- p. 72 ---
\setcounter{footnote}{0}%
\opage{72}den paavirkede i Renaissancetiden Bruno, og den blev genoptagen af
Basso og Gassendi, fra hvilken sidste Newton optog den (i sin Optik). Fra
Newton optog igen John Dalton den,%
% footnote 1 on p. 72
\footnote[1]{Smlgn. \textsc{Roscoe:} \textit{John Dalton,} p. 127---145.} idet
han, følgende Lavoisier's kvantitative Metode, viste, at Atomet af et Grundstof
har en konstant Vægt, og at sammensatte Stoffer opstaar ved Forbindelser af
forskellige Grundstoffer i bestemte Vægtforhold.

Der var for den nyere Tids Forskere det betænkelige ved Demokrits og Newton's
Atomer, at de var absolut haarde og udelelige. Saa maatte der gaa Kraft tabt
ved deres Sammenstød. Huyghens og Leibniz antog derfor, i Konsekvens af deres
Hypotese om Kraftens Bestaaen, at Atomerne er elastiske. Meget stærkt
fremhævede Leibniz dette overfor Newtonianeren Clarke, der ligesom andre af
Newton's Disciple, f. Ex. Maclaurin, forkastede Kraftens Bestaaen og netop
fandt et afgørende Bevis mod denne i de absolut haarde Atomer. Leibniz antog,
at den Kraft, Atomet som Helhed mister ved Sammenstød, bevares i dets Dele, der
sættes i Vibrering (agitation) formedelst Sammenstødet. Der er altsaa ikke
gaaet Kraft tabt, men den er fordelt paa flere Punkter. Det er, som naar store
Penge vexles i Smaapenge.%
% footnote 2 on p. 72
\footnote[2]{\textit{Lettres entre Leibniz et Clarke.}(Leibnitii opera
philosophica ed. Erdmann, p. 775).} Om denne Udtalelse af Leibniz har
\textsc{Henri Poincaré} bemærket, at man ikke paa nogen bedre Maade kunde
udtrykke den Hypotese, der ligger til Grund for den mekaniske Varmeteori%
% footnote 3 on p. 72
\footnote[3]{Se Citatet af Poincaré hos \textsc{Léon Brunschwieg:}
\textit{L'expérience humaine et la causalité physique,} p. 227.}. Gennem denne
Modifikation af Atomteorien dannes Overgangen til den nyeste Tids Opfattelse,
ifølge hvilken Atomiciteten er relativ, idet ethvert Atom er et helt lille
Verdenssystem, hvis Elementer kaldes Elektroner. Demokrits »finere Tænkning«
var altsaa endnu
% --- p. 73 ---
\setcounter{footnote}{0}%
\opage{73}ikke fin nok, men den nyere Tids Videnskab har udviklet Finheden i
høj Grad --- og hvorfor skulde Tænkningen ikke kunne blive endnu finere?

\textsc{Émile Meyerson} har i sine Skrifter, og især i en interessant
Diskussion i det franske filosofiske Selskab (Bulletin de la Société française
de philosophie 3, mars 1910) med stor Sagkundskab paavist Analogien mellem
antik og moderne Atomistik. Han vil ikke sige, at Demokrit og Lucretius har
»Prioritet« overfor moderne Kemikere og Fysikere; endog det at kalde dem for
Forgængere vilde, mener han, være at miskende Videnskabernes Histories sande
Aand. Og dog er, hævder han, Sammenstillingen af antik og moderne Atomteori
ikke vilkaarlig; der finder en virkelig Analogi Sted, idet de to Teorier
udspringer af samme Tankegang.

b. Som med antik Atomistik i Forhold til moderne Atomistik forholder det sig
idethele med antik Tænkning i Forhold til moderne Tænkning. Modsætningen er
tydelig, og ofte nok fremhævet.%
% footnote 1 on p. 73
\footnote[1]{Se f. Ex. \textit{Den menneskelige Tanke,} p. 123---134.} For
Platon, den mest udpræget græske Tænker, var det Videnskabens Opgave at danne
et System af Begreber (»Ideer«), under hvilke alle foreliggende Emner kunde
henføres. I et saadant System vilde man have Nøglen til al Tilværelse, mener
han. Han søgte paa sin Maade Identiteten bag alt forskelligt og skiftende.
Muligheden af et klassifikatorisk System, indenfor hvilket alle forskellige
Emner fandt deres Plads under stedse højere Lighedssynspunkter, var ham et
Vidnesbyrd om, at Ideernes Verden er den sande Verden. I den nyere Tid er
Begrebet Lov sat istedetfor Begrebet Ide. Man flygter ikke for
Forskellighederne og Forandringerne, men finder i Lovene for Emnernes Existens
og for de Forandringer, der foregaar,
% --- p. 74 ---
\setcounter{footnote}{0}%
\opage{74}den Identitet, som Platon fandt i Ideernes Verden. Bacon og Bruno
fortolker allerede de platoniske Ideer som Emnernes Love, og Galilei's
Platonstudier var af ikke ringe Betydning for hans Forsken. Identitetsprincipet
raader i de antike Tænkeres klassifikatoriske Systemer som i den nyere Tids
forklarende Systemer. I antik Videnskab søges Kvaliteterne reducerede til
Identiteter gennem Arts- og Slægtskabsbegreber. I moderne Videnskab opfattes
Kvaliteter som Identiteter ved Hjælp af Kvantitetsbestemmelser. Ogsaa moderne
Tænkning har sin Ideverden, nemlig i den Kreds af Synspunkter, ud fra hvilke
det bliver muligt, ikke blot at klassificere ved Identitetsprincipets Hjælp,
men ogsaa --- ved Hjælp af det samme Princip --- at naa rationelle
Forklaringer. Ideverdenen har da ligesom sænket sig ned i Erfaringens Verden,
ti det er ingen Selvfølge, at de Emner, Iagttagelsen frembyder, skulde kunne
gøre klassifikatorisk og explicativ Erkendelse mulig. Det kunde jo være, at
Tilværelsen havde været af den Beskaffenhed, at der ingen Mulighed var hverken
for Klassifikation eller Forklaring! Men enhver stor eller lille Sejr,
Videnskaben vinder, er et Vidnesbyrd om det modsatte.

Analogien mellem antik og moderne Videnskab lægger sig ogsaa for Dagen deri, at
de Vanskeligheder, der frembød sig for Platon's Ideer, og som hans Genialitet
og hans Sandhedskærlighed gav ham saa klart et Blik for (se første Del af
Dialogien »Parmenides«), i nye Former gør sig gældende for den moderne
Videnskabs Loverkendelse. For Platon var Vanskeligheden den: naar Ideernes
Verden er afsluttet i sig selv, idet den ene Ide kan udledes af den anden, og
alle Ideer tilsidst af en højeste Ide, hvorledes kan der da bestaa en Emnernes
Verden, hvis hele Natur ikke er udtømt ved de Elementer, der finder Plads i
Ideerne? Paa tilsvarende Maade staar det som et stadigt Pro%
% --- p. 75 ---
\setcounter{footnote}{0}%
\opage{75}blem for moderne Videnskab, hvorledes de enkelte Emner i deres
konkrete, individuelle Ejendommelighed skulde kunne udledes af de almene Love,
dersom disse skal være Videnskabens sidste Ord eller maaske endog staa som
Udtryk for Tilværelsens inderste Væsen. Her er et Problem, der gør sig gældende
baade paa Naturvidenskabens og paa Aandsvidenskabens Omraade. Problemet stiller
sig i nye Former; men det følger stedse Tanken paa denne Vej.

c. Indenfor den nyere Tids Forsken gør sig en betydningsfuld Analogi gældende
mellem det 17. og det 19. Aarhundredes Naturvidenskab. Det 17. Aarhundrede
grundlægger den nyere Tids Tankeliv. De Hovedtanker, som senere fik deres
storartede experimentale Bekræftelse, kommer her for første Gang frem. Man
anvendte dristig Identitetsprincipet paa Forholdet mellem Aarsag og Virkning.
Descartes og Spinoza udtalte, at hvis der i Virkningen er andet og mere, end
der er i Aarsagen, saa er fuld Forklaring ikke naaet. Huyghens og Leibniz
anvendte denne Tanke i deres Hævden af Kraftens Bestaaen: hverken kan der ved
Overgangen fra Hvile til Bevægelse opstaa ny Kraft, ej heller kan der ved
Overgangen fra Bevægelse til Hvile forsvinde Kraft. Overfor denne Tanke paastod
Newtonianeren \textsc{Clarke} og \textsc{Maclaurin,} at den stred mod Erfaring%
% footnote 1 on p. 75
\footnote[1]{\textsc{Maclaurin:} \textit{Exposition des Découvertes de
Newton.}Trad. franç. 1749, p. 88.}. Men Tanken genoptoges --- som
selvindlysende! --- ved det 19. Aarhundredes Midte af de Forskere, der omtrent
samtidig opdagede Sætningen om et Ækvivalensforhold mellem de forskellige
specielle Naturkræfter, ifølge hvilket der ved Overgang fra en Energiform til
en anden hverken opstaar eller forsvinder Energi. Et Studium af disse Forskeres
Arbejder viser, at de lededes af den Tanke, at der maatte
% --- p. 76 ---
\setcounter{footnote}{0}%
\opage{76}være et Ækvivalensforhold mellem Aarsag og Virkning. Det stod for dem
som selvindlysende; men selv det »selvindlysende« har sin Historie, og dette
bestemte selvindlysende har ikke været til før det 17. Aarhundrede. Heldigvis
lod de omtalte Forskere \textsc{(Robert Mayer, Colding, Joule, Helmholtz)} sig
ikke nøje med en Deduktion ud fra en »selvindlysende» Sætning, men søgte
Beviser fra Iagttagelser og Experimenter. En »ledende Tanke« er jo ikke en
Tanke, der uden videre kan bruges som Oversætning i en Syllogisme eller som
Forudsætning for en Deduktion, men er en Tanke, der kan tjene som Forbillede
(»Paradigme«), som Grundlag for en Analogi, der kan føre til at vinde
Erkendelse paa nye Omraader. Tanken om Kraftens Bestaaen har i og for sig
Betydning, nemlig paa den rene Mekaniks Omraade. Men det var jo muligt, at dens
Betydning ikke strakte sig ud over dette Omraade, saa at den ikke kunde hjælpe
til Forstaaelse af de specielle Naturkræfters indbyr-Forhold. Denne Mening
havde Positivismens Grundlægger Auguste Comte, for hvem de forskellige
Videnskabsgrene var principielt adskilte fra hverandre, saa at en Tanke, der
kan passe paa det ene Omraade, ikke havde nogen Betydning for de andre
Omraader. Fysikens forskellige Grene vilde efter hans Mening vedblive at staa
adskilte fra hverandre. Han antog ganske vist, at det i Længden ikke vilde være
muligt for den menneskelige Aand at tænke strengt videnskabelig paa et Omraade
og uvidenskabelig paa andre Omraader; der vil uvilkaarlig stræbes efter Enhed i
Metode og Teori. Men han tillagde ikke den Analogi, der her gjorde sig
gældende, videnskabelig Betydning; han udledte den nærmest af Vanens
Indflydelse. Grundlæggerne af Sætningen om Energiens Bestaaen gik derimod ud
fra den Mulighed, at der objektivt kunde paavises Forhold mellem
% --- p. 77 ---
\setcounter{footnote}{0}%
\opage{77}de specielle Naturkræfter indbyrdes, der svarer til Forholdet mellem
Bevægelse og Hvile i den rene Mekanik, og det var i Tillid til denne Analogi,
at de gik til deres Arbejde. --- Vi vil i det følgende træffe flere Exempler
paa Analogiens Betydning, naar det gælder Forholdet mellem forskellige
Videnskaber; men det er ikke altid, at Analogien er saa frugtbar som i det nys
omtalte Tilfælde.

I min Bog »Ledende Tanker i det 19. Aarhundrede« har jeg ment at kunne
karakterisere dette Aarhundrede som gennemførende og verificerende i Forhold
til de forudgaaende Aarhundreders Tanker, ved hvis Fremkomst den nye Tids
Videnskab overhovedet blev til. I sin smukke og lærerige Fremstilling af
Fysikens og Kemiens Historie i det 19. Aarhundrede mener \textsc{Niels
Bjerrum,} at jeg har været for beskeden paa Naturvidenskabens Vegne. Dertil vil
jeg sige, at det ingenlunde har været min Mening at benægte, at det store
Tankearbejde, der indenfor Naturvidenskaben øvedes i det 19. Aarhundrede, har
selvstændig og fremragende Betydning, og at der er kommet meget nyt og
overraskende frem i Løbet af dette Aarhundrede. Men for en sammenlignende
historisk Betragtning er det netop nyt og overraskende at se, hvorledes Tanker,
der hidtil kun var komne frem i Erkendelsesteoriens og Mekanikens abstrakte
Former, nu gennemførtes i det enkelte paa Erfaringens og Experimentets Grund.
Ved »ledende Tanker« maa først og fremmest forstaas Retninger, i hvilke
Tænkning og Forskning gaar. Et er Retningen, et andet Tankearbejdet. Selve
Indledelsen af en ny Retning kræver dog Tankearbejde, og det kan være et meget
vanskeligt Spørgsmaal at afgøre, om der til Grundlæggelsen af en helt ny
Retning i Forskningen (som da den nye Videnskabs Principer opstilledes i det
17. Aarhundrede) kræves min%
% --- p. 78 ---
\setcounter{footnote}{0}%
\opage{78}dre eller større Tankearbejde end til at gennemføre en allerede
indledet Retning i det enkelte. Som alt Arbejde er ogsaa Tankearbejde
Overvindelse af en Modstand, og den Modstand, Tænkningens første Pionerer
mødte, var maaske nok saa stor som den, det 19. Aarhundredes Forskere, der
gennemførte hines Tanker, mødte paa deres Vej.%
% footnote 1 on p. 78
\footnote[1]{Den literære Skæbne, de fire ovenfor nævnte Forskeres første
Arbejder foreløbig fik, vidner om, at der ogsaa her var Modstand nok.}
Modstanden har været af noget forskellig Art, og en sikker Maalestok for den
har vi ikke. Det var i hvert Tilfælde ikke min Opgave at undersøge dette Punkt.
Jeg finder det 19. Aarhundredes Storhed i den energiske Trofasthed, hvormed
ledende Tanker, der allerede havde set Lyset forud, fastholdtes og gennemførtes.

d. Indenfor Filosofiens Historie har \textsc{Hans Brøchner} i sin Bog »Bidrag
til Opfattelse af Filosofiens historiske Udvikling« (1869) givet interessante
Bidrag til Belysning af Analogiens Betydning for filosofisk Tænken. Han
finderen Sammenhæng mellem de store Perioder i Filosofiens Historie, betragtede
som Helheder, og de enkelte Tankeforsøg indenfor enhver af dem, idet disse
Tankeforsøg snart foregriber, hvad der først fremtræder i fuldt udpræget Form,
naar Perioden har naaet sit Højdepunkt, snart vedbliver at bære Periodens
Særpræg, efterat den egentlig er sluttet. Brøchner antager en hemmelighedsfuld
Indflydelse af en Periode paa de enkelte Tænkere; men hvad der foreligger er
aabenbart Analogier mellem Tankeforsøg, som finder Sted under tilnærmelsesvis
samme historiske Forhold, nærmere eller fjernere fra disse Forholds Kulmination.

De store Perioder, Filosofiens Historie frembyder, indledes ofte ved, at der
fremsættes Anskuelser, der først udføres nærmere længere frem i Perioden.
Saaledes foregribes
% --- p. 79 ---
\setcounter{footnote}{0}%
\opage{79}i den ældste græske Filosofi de Tanker, der først hos Platon kom til
fuld Klarhed. I Begyndelsen af Renaissancetiden gøres der atter og atter Tilløb
til, hvad først Giordano Bruno giver i systematisk Sammenhæng. I mange
Henseender foregriber Kant den romantiske Idealismes Tanker. --- Paa den anden
Side kan Tanker, der er fostrede indenfor en Periode som berettiget Udtryk for
dens særlige Opgaver, virke hæmmende og begrænsende paa Tænkere, hvis
personlige Stræben i og for sig maatte føre dem ud over Periodens Standpunkt.
De Muligheder, der hos Aristoteles og Plotinos frembød sig for et andet
Standpunkt end det klassisk græske, hæmmedes af Forudsætninger, som var fælles
for al græsk Filosoferen. Men saadan hæmmende Indflydelse fra Fortiden kan ofte
fremkalde dybere gaaende Tanker, end der ellers vilde være mulige. Saaledes har
græsk Tænkning øvet en gavnlig Indflydelse i Middelalderen, og Middelalderens
Tænkning i den nyere Tid.

Brøchner betragter Filosofiens Historie som en i sig selv afsluttet Verden.
Mere indviklet bliver Forholdene, naar det tages i Betragtning, at den Maade,
hvorpaa der filosoferes, i hver given Tid er betinget, ikke blot ved
forudgaaende Tiders Filosoferen, men i høj Grad ogsaa ved nye Erfaringer og ved
det Arbejde, der finder Sted indenfor andre Videnskaber. Men selv indenfor sin
bestemte Begrænsning har Brøchner's fine og aandfulde Betragtning sin Interesse
til Belysning af Analogiens Optræden i menneskelig Tænken.

\begin{center}\large IV. Analogi mellem Erkendelsesomraader.\end{center}

15. Til de forskellige Erkendelsesomraader, altsaa til de forskellige Grupper
af Videnskaber, svarer forskellige Grup%
% --- p. 80 ---
\setcounter{footnote}{0}%
\opage{80}per af Kategorier. Fælles for alle Videnskaber er de fundamentale
Kategorier: Syntese og Relation, Kontinuitet og Diskontinuitet, Lighed og
Forskel, og de er tillige fælles for den førvidenskabelige og den
videnskabelige Tænken. Til Logik og Matematik svarer de formale Kategorier:
Identitet, Negation, Rationalitet. Til Natur- og Aandsvidenskaberne svarer de
reale Kategorier: Kausalitet, Totalitet, Udvikling. Til de etiske Videnskaber
svarer de ideale Kategorier: Værdibegreberne.

Analogien mellem den førvidenskabelige og den videnskabelige Bevidsthed har jeg
allerede omtalt. Jeg gaar nu over til at paavise Analogier mellem de formale
Videnskaber indbyrdes, mellem formale og reale Videnskaber, mellem reale
Videnskaber indbyrdes, og endelig mellem formale og reale Videnskaber paa den
ene Side og de etiske Videnskaber paa den anden Side.

\begin{center}16. Logik og Aritmetik.\end{center}

Enhver Videnskab maa selv, gennem Analyse af sine Emner og sine Metoder,
fastsætte sine nødvendige Grundbegreber i saa simpel og saa frugtbar Form som
muligt. Men derefter kan fra filosofisk Side det Spørgsmaal opkastes, om de
saaledes opstillede Grundbegreber ikke igen forudsætter endnu simplere
Begreber, og hvis de gør det, om de da uden videre kan udledes af disse eller
maa opstilles som nye, speciellere Begreber. Et Begreb kan meget vel forudsætte
et andet Begreb uden at kunne udledes af dette. Hvad Forholdet mellem Logik og
Matematik angaar, bliver Spørgsmaalet dette, om de ligger paa samme Linie,
hvilende paa samme Grundlag, saa at Matematik betyder en ligefrem Fortsættelse
af Logik, eller om Matematiken møder med nye Grundbegreber, mellem hvilke og de
rent logiske Grundbegreber der kun kan bestaa Analogi, ikke Identitet.

Den rent logiske Tænken vilde, som paavist i det fore%
% --- p. 81 ---
\setcounter{footnote}{0}%
\opage{81}gaaende, staa afmægtig, dersom alle foreliggende Emner dannede enten
en kaotisk Forskelsrække eller en absolut Identitetsrække. Den er afmægtig,
baade hvor der slet ingen reale Ligheder er, og hvor der slet ingen reale
Forskelligheder er. Selv under disse, logisk set, fortvivlede Forhold kan
matematisk Tænkning gøre et Skridt endnu, og et Skridt, der kan føre mange
Skridt efter sig. Forskellighederne eller Identiteterne kan nemlig tælles, d.
v. s. der kan dannes en Række, mellem hvis Led der er Forskelsidentitet, en
Række, der stedse kan fortsættes efter Forskelsidentitetens Lov, samme Lov,
efter hvilken Rækken begynder, og ligegyldigt, om Forskellighedernes eller
Identiteternes faktisk foreliggende Række er begrænset. Talrækken skyldes en
Konstruktion, der blot bestaar i Gentagelse af en og samme Tankeakt: ethvert
nyt Led findes ved at statuere en Forskel, der er identisk med Forskellen
mellem tidligere Led i Rækken. Talbegrebet bygges paa et Lighedspunkt mellem
reale Forskelligheder (nemlig at de er Genstande for Tanke eller Opmærksomhed)
og et Forskelspunkt mellem reale Ligheder (nemlig at de er Genstand for
forskellige Tankeakter). --- Allerede den førvidenskabelige Bevidsthed forsøger
at danne saadanne Rækker. Barnet »tæller« sine Kegler (»nok en, nok en, nok en«
o. s. v.), saa længe til det bliver ked deraf; det kniber kun med
Opsummeringen. Den Vilde tæller paa sine Fingre og opsummerer ved at knytte
Haanden. Selv om der ikke er flere Emner givne, kan Operationen foretages ved
Hjælp af selvskabte Tegn. Gennem saadan frembringende Gentagelse bliver
Tallæren andet og mere end formal Logik. Af Identitetsprincipet kan Muligheden
af stedse at føje nye Led til de allerede givne ikke udledes.
Identitetsprincipet anvendes logisk paa Kvali%
% --- p. 82 ---
\setcounter{footnote}{0}%
\opage{82}teter som paa Kvantiteter; men gennem frembringende Gentagelse
opstaar ikke nye Kvaliteter; der dannes kun en Kvantitetsrække, eller, som det
simplest kan udtrykkes, det rene Størrelsesbegreb bliver til. Aritmetikens
Grundbegreber er altsaa speciellere end den rene Logiks, og de kan ligesaa lidt
udledes af disse, som overhovedet et Artsbegreb kan udledes af et Slægtsbegreb.
Hvis man vilde sige, at et Tal er et Almenbegreb eller et Klassebegreb, ligesom
de andre Begreber, Logiken opererer med, saa er der dog den vigtige Forskel
mellem Talklasser (f. Ex. Femklasser, Sexklasser o. s. v.) og andre Klasser (f.
Ex. Hestenes Klasse), at medens disse sidstes Definitioner indeholder
kvalitativt forskellige Elementer, indeholder en Talklasses Definition kun
indbyrdes identiske Elementer%
% footnote 1 on p. 82
\footnote[1]{Smlgn. allerede \textsc{Platon:} \textit{Staten} VI, p. 526 A,
hvor det hedder om de aritmetiske Enere: ἴσον τε ἕκαστον πᾶν παντὶ καὶ οὐδὲ
σμικρὸν διαφέρον, μόριον τε ἔχον ἐν ἑαυτῷ οὐδέν.}. Begrebet Hest har til
Indhold: »uparret-taaet Hovdyr med særlig udviklet Midttaa og lang Mellemfod,
med fuldstændigt Tandsæt og med Benring om Øjet«. Begrebet 6 har sit Indhold
givet i Rækken 1 +l + 1 + 1 + 1 + 1, der fremkommer ved Gentagelse af en
Forskelsidentitet, en Gentagelse, vi selv kan foretage uden at afvente
Iagttagelse. For Logiken er Talrækken kun et Exempel paa en særegen Art
Rækkedannelse, en identisk varierende Forskelsrække eller tillige en
fremskridende Forskelsrække (alt eftersom Talen er om Ordenstal eller
Kardinaltal). Det er for Logiken ligegyldigt, hvorledes en saadan Række bliver
til; den tager det som et Faktum, at den kan dannes. Men hvad der for Logiken
er et enkelt Exempel, det er for Aritmetiken af grundlæggende Betydning.
Talrækken grundes paa et Postulat, idet der hævdes Muligheden af en
Konstruktion som den omtalte. Det er gennem dette Postulat, at vi kommer fra
Logiken over i
% --- p. 83 ---
\setcounter{footnote}{0}%
\opage{83}Aritmetiken. Og den samme Proces, der fører til at danne Talrækken,
fører til at fortsætte den gennem negative, brudte, irrationelle og imaginære Tal.

Dette Postulat, der indeholder, hvad man kan kalde »matematisk Induktion« eller
»Gentagelsens Lov« (loi de recurrence), blev i rent matematisk Form anvendt af
Maurolyco (16. Aarh.) og af Pascal (17. Aarh.)%
% footnote 1 on p. 83
\footnote[1]{\textit{Œuvres de Pascal,} éd. Brunschwicg. III, p. 456; VIII, p.
363 (med Pierre Boutroux's Noter).}. \textsc{John Locke} anvender det som
Grundlag for Begreberne Tal, Rum og Tid (Essay II, 12, 8 og 13, 1---2). Som
karakteristisk for Matematik i Modsætning til Logik er det opfattet af Leibniz,
Kant, Boole og Poincaré%
% footnote 2 on p. 83
\footnote[2]{Smlgn. \textit{Totalitet som Kategori,} p. 33---39.}. Platon's
Lære om Matematiken som liggende mellem den rene Logik (»Dialektik«) og
Iagttagelsernes Verden (Staten p. 511 A. D; 525 A; 527 B) bekræftes derved.

I den nyeste Tid har \textsc{Hilbert} udtalt sig ikke blot mod dem, der som
Frege o. a.%
% footnote 3 on p. 83
\footnote[3]{Saaledes B. \textsc{Russell:} \textit{Introduction to mathematical
philosophy.}Chap. 2.} opfatter Tallæren som en Del af Logiken, idet de
betragter Tallene som Artsbegreber, men ogsaa mod dem, der som Poincaré og
andre bygger den paa Evnen til at gentage en Tankeakt. Denne sidste Opfattelse
betegner Hilbert som mystisk. I formentlig Modsætning hertil hævder han%
% footnote 4 on p. 83
\footnote[4]{\textit{Neubegründung der Mathematik.} (Abhandlungen des
mathematischen Seminars der hamburgischen Universität. I, 2 (1922), p.
161---163.}, at man maa begynde med »visse diskrete Objekter, der paa anskuelig
Maade er givne forud for al Tænken«. Saadanne Objekter har vi i Tegnene,
Talteoriens egentlige Genstande: »Am Anfang ist das Zeichen!« Tegnet 1 er et
Tal. Og et Tegn, der begynder med 1 og ender med 1, idet der i Mellemrummet
efter hvert 1 stedse kommer et + og efter hvert + et 1, er ogsaa et
% --- p. 84 ---
\setcounter{footnote}{0}%
\opage{84}Tal. Saaledes faar vi Rækken 1 + 1 + 1 + 1.... For lettere at kunne
meddele os til andre danner vi saa Tegnet 2 for 1 + 1, 3 for 1 + 1 + 1 o, s. v. %
% PRINTED AS IS: „o, s. v.“ (p. 84)


Jeg kan ikke se andet, end at denne Opstilling fører tilbage til den
Opfattelse, som er gjort gældende af Poincaré og hans matematiske og
filosofiske Forgængere. Et Tegn og en Række af Tegn bliver dog kun til ved
Opmærksomheds- eller Tankeakter, og en Række som den opstillede forudsætter
Muligheden af, at en saadan Akt kan gentages paa en og samme Maade (hvad der
naturligvis, psykologisk set, kun tilnærmelsesvis er muligt). Der er intet
mystisk i den Kendsgerning, at vi kan gentage vore Tanker eller, om man vil,
vore Tegn paa den Maade, at Rækken stadig fortsættes efter samme Regel, som
førte til Opstilling af de to første Led.

Deri er dog (og det er i denne Sammenhæng det vigtigste) Hilbert enig med den
anden Gruppe af Forskere, at Matematik ikke kan udledes af Logik. Et Tegnsprog
af den Art, han indfører, har Logiken ingen Brug for.

Der raader da ikke Identitet mellem Aritmetik og Logik. Derimod raader der
Analogi. Tal kan forholde sig til Tal som Grund til Følge, og Hilbert faar da
ogsaa Brug for det logiske Tegn → som Udtryk for Forholdet mellem Grund og
Følge. (Det er samme Tegn, jeg i »Den menneskelige Tanke« p. 168 brugte til at
betegne Forholdet mellem Leddene i en partiel Identitetsrække). De rent logiske
Forhold genfindes paa det matematiske Omraade, og det paa en saa enestaaende
klar og nøjagtig Maade, at de bedste og frugtbareste Exempler paa logisk
Slutten foreligger her. Men det er dog kun Exempler mellem andre Exempler. Og
indenfor Tallæren er der Analogi mellem Algebra og Aritmetik, idet man i
Algebra kun fæster Op%
% --- p. 85 ---
\setcounter{footnote}{0}%
\opage{85}mærksomheden paa den Maade, hvorpaa Størrelserne kombineres, uden at
kende selve Størrelserne. Saaledes udtrykker (\textit{a} + \textit{b})
\textit{c}\textsuperscript{2} alle Tilfælde, i hvilke en Størrelse bestemmes
ved Produktet af to Størrelsers Sum og en tredie Størrelses Kvadrat,
ligegyldigt, hvilke Værdier \textit{a, b} og \textit{c} har i de specielle Tilfælde%
% footnote 1 on p. 85
\footnote[1]{\textsc{Pierre Boutroux:} \textit{L'idéal scientifique des
mathématiciens.}1920, p. 84.}. Vi genfinder her det logiske Forhold mellem et
Begreb og de Exempler, for hvilke det gælder. Ethvert Exempel maa kunne
substitueres for et andet overalt, hvis Begrebet skal gælde. I Tallæren er
denne Substitution ubetinget; ved Begreber om kvalitative Emner gælder den i
hvert enkelt Tilfælde kun fra ét bestemt Synspunkt.

Forholdet mellem Logik og Aritmetik bliver Analogi; paa det Grundlag, der
bliver til ved Konstruktion af Talrækken gennem frembringende Gentagelse, er en
ny og fuldkomnere Anvendelse af det logiske Forhold mellem Grund og Følge
mulig; men hin Konstruktion kan ikke udledes af den rene Logik som saadan. Der
gives en ny Indsats i Erkendelsens Udvikling ved Talbegrebets Opstaaen.

\begin{center}17. Aritmetik og Geometri.\end{center}

Ikke blot Tallæren, ogsaa Geometrien beror paa Postulatet om Muligheden af
Gentagelse af en og samme Tankeakt. Hvilken saa en Tankeakts Indhold eller
Genstand er, kan jeg bestandig føje en ny Tankeakt til de tidligere. Jeg kan
tænke mig flere Heste end dem, jeg har set, og flere »Tegn« end dem, jeg
allerede har nedskrevet, flere Liniestykker end dem, jeg allerede har tegnet.
Men ved nøjere Undersøgelse er man dog bleven ført til at gøre en stedse
skarpere Forskel mellem Læren om Tal og Læren
% --- p. 86 ---
\setcounter{footnote}{0}%
\opage{86}om andre Emner, hvad enten disse er geometriske Figurer eller
konkrete Genstande. Det ejendommelige for Tallet er dets diskrete Natur. Det er
bestemt blot ved at være en Ener ved Siden af andre Enere eller ved at bestaa
af Enere, der har denne Egenskab. Men ved de Emner, der tælles, er der stedse
andre Egenskaber foruden den at kunne tælles. Det er allerede Tilfældet med
geometriske Linier. I og for sig kunde man istedetfor at udtrykke Enere ved
lodrette Streger (l + l + l + l....) benytte vandrette Streger (--- --- ---
---....), og naar saa den ene umiddelbart fort satte den anden, var en ret
Linie konstrueret. Denne opstaar da forsaavidt ved samme »Synthesis des
Gleichartigen« som Talrækken. \textsc{Kant} har netop tænkt sig Sagen paa denne
Maade, naar han siger: »Wir können uns keine Linie denken, ohne sie im Gedanken
zu ziehen«, og naar han hævder, at vi »im Ziehen der geraden Linie bloss auf
die Synthesis des Mannigfaltigen Acht haben«. (»Kritik der reinen
Vernunft«\textsuperscript{2} p. 154). Der sker jo ogsaa en Tilføjelse hvert
Øjeblik, medens vi drager Linien, ligesaa hvergang vi føjer en Ener til i
Talrækken. Men det store Spørgsmaal er, om vi har Ret til at sige, at de
stadige Tilføjelser, medens vi drager Linien, er ligestore. Ved Talrækken
opstod dette Spørgsmaal ikke, og det af den simple Grund, at hvert enkelt Led
blot betragtes under Synspunktet af en Tilføjelse, --- ene og alene som noget,
der føjes til, uden at der er Tale om Udstrækning eller om andre Egenskaber.

Denne principielle Forskel mellem Tallære og Geometri har successivt arbejdet
sig frem. Den græske Matematik havde sit Grundlag i Geometrien. Talforhold
anskueliggjordes ved Rumforhold, Produktet af to Tal f. Ex. ved Rektanglet af
to Linier, der forestillede Faktorer. Med Descartes kom et Vendepunkt. Ikke
blot blev nu omvendt
% --- p. 87 ---
\setcounter{footnote}{0}%
\opage{87}Figurer udtrykte ved Tal og Ligninger, der angav Forholdet mellem
Figurernes Elementer, men geometriske Problemer blev reducerede til algebraiske
Problemer. Der blev, som \textsc{Zeuthen} \textsc{%
% footnote 1 on p. 87
\footnote[1]{\textit{Matematikens Historie.} II, p. 285.}} udtrykker sig, vendt
helt om paa det traditionelle Forhold mellem Algebra og Geometri, og derved
betegnes en Reform i den hele Matematik, som hidtil havde haft sit anerkendte
Grundlag i Geometrien, men fra nu af fik det i Algebraen. Et Skridt videre gik
\textsc{Maxwell} (1855)%
% footnote 2 on p. 87
\footnote[2]{\textit{Über Faraday's Kraftlinien.} Übers. von \textsc{Boltzmann}
(1895). Einleitung p. 4.}, da han udtalte, at alle videnskabelige Anvendelser
af Matematiken beror paa Forhold mellem de fysiske Størrelsers Love og de hele
Tals Love. Naar \textsc{Boltzmann} i en Note til sin Oversættelse af dette Sted
hos \textsc{Maxwell} siger, at saavidt han ved, er det Standpunkt, at Maalinger
af Rums- og Tidsstørrelser ved Tal blot beror paa Analogi med de mellem Tallene
bestaaende Forhold, ikke siden taget op, saa gælder dette ikke længere. Flere
Forskere har fortsat \textsc{Maxwell's} Tankegang. Saaledes lægger
\textsc{Ernst Mach} (som idethele har henledet Opmærksomheden paa
\textsc{Maxwell's} erkendelsesteoretiske Synsmaade) stor Vægt paa Grænserne for
Maalingens Nøjagtighed. »Dækningen af en Maalestok og det Emne, der skal
maales, kan«, siger han%
% footnote 3 on p. 87
\footnote[3]{\textit{Erkenntnis und Irrtum.} (1905), p. 328 f.}, »ikke
fastslaas med ubegrænset Nøjagtighed«. Og Skridtet er gjort helt ud af
\textsc{Hjelmslev} i hans Fremstilling af den analytiske Geometri%
% footnote 4 on p. 87
\footnote[4]{\textit{Elementær Geometri.} Tredie Bog. 1921, p. 71; 91: 180 f.}.
\textsc{Hjelmslev} skelner skarpt mellem Virkelighedsgeometrien, der handler om
Ting, og den abstrakte Geometri, der handler om Tal. I et Kapitel om den
geometriske Maalings Teori viser han, at ingen Maaling af et paa en Tegning
forelagt Liniestykke kan være helt nøj%
% --- p. 88 ---
\setcounter{footnote}{0}%
\opage{88}agtig; der opnaas ikke ét, men flere Resultater, mellem hvilke der
maa vælges, for at Liniestykket kan blive »fixeret«. Tegninger er stadig
nødvendige til Vejledning og Beskrivelse, men de kan kun opfattes som Symboler
paa de abstrakte Linier, der kan udtrykkes ved bestemte Tal.

\textsc{Zeuthen} betragtede i det sidste Skrift, vi fik fra hans Haand%
% footnote 1 on p. 88
\footnote[1]{\textit{Hvorledes Matematiken i Tidsrummet mellem Platon og Euklid
blev rationel Videnskab} (1917), p. 9.}, de ideale eller abstrakte Figurer,
hvis Begreber dannes ved Løsrivelse fra, hvad Sansning viser, som Symboler,
ligesom den nyere Matematiks Bogstaver og Operationstegn kan siges at være det.
Saadanne ideale eller abstrakte Figurer kan nu ifølge \textsc{Hjelmslev} kun
udtrykkes i Tal. Paa en interessant Maade har saaledes de to danske
Matematikere grebet hver sin Side af den Korrelation, Begrebet Symbol
indeholder. Naar \textit{A er} Symbol paa \textit{B,} er \textit{B} nemlig
ogsaa Symbol paa A. \textsc{Hjelmslev,} hvis Opfattelse af Geometrien nærmest
er at kalde realistisk, opfatter de virkelige geometriske Former som Symboler
paa de abstrakte Former, der efter ham kun udtrykkes ved Tal. Han kan sige med
\textsc{Goethe:} »Alles Vergängliche ist nur ein Gleichnis!« Men
\textsc{Zeuthen,} hvis Opfattelse af Geometrien nærmest er at kalde
idealistisk, da han hævder ideale Definitioners Nødvendighed i Geometrien,
betragter omvendt disse Idealer som Symboler paa de virkelige geometriske
Forhold. ---

Der foreligger her et overordentlig interessant Exempel paa, hvorledes Analogi
kan være den eneste Maade, paa hvilken Forholdet mellem forskellige
Erkendelsesomraader kan betegnes. Og her, som paa andre Punkter, har
Tænkningens Udvikling ført til, at man finder Analogi dér, hvor man forhen
fandt Identitet.
% --- p. 89 ---
\setcounter{footnote}{0}%
\opage{89}

\begin{center}18. Formal og real Videnskab.\end{center}

Analogi vil nu vise sig ikke blot at være det rette Udtryk, naar Forholdet
mellem de forskellige formale Videnskaber (Logik, Aritmetik, Geometri)
indbyrdes skal bestemmes, men ogsaa, naar Forholdet mellem formal Videnskab og
real Videnskab skal bestemmes. --- Vi holder os foreløbig til Naturvidenskaben
som Repræsentant for Realvidenskaben og gaar senere over til at drøfte
Forholdet mellem Naturvidenskab og Aandsvidenskab.

Det er Forholdet mellem Begreberne Grund og Aarsag, det her drejer sig om. Paa
dette Punkt ligger \textsc{Kant's} Stordaad. Han har sat Analogi som Udtryk for
Forholdet mellem Grund og Aarsag istedetfor den Forudsætning om disse Begrebers
Identitet, som dogmatiske Filosofer og Naturforskere gik ud fra. Forholdet
mellem Aarsag og Virkning betyder ifølge \textsc{Kant} kun, at to Emners
(Begivenheders) Rækkefølge i Tiden er analogt med det logisk-matematiske
Forhold mellem Grund og Følge. Realvidenskaben arbejder uafbrudt paa at
godtgøre, at enhver Forandring, der sker i Verden, alt nyt, der optræder for
vor Iagttagelse, forholder sig til forudgaaende Forandringer, som en Slutning
forholder sig til sine Forudsætninger, altsaa som Følge til Grund. Al Mystik i
Aarsagsforholdet falder derved bort. Vi véd intet om dette Forhold, uden at det
paa én Gang er et Grundforhold og et Tidsforhold. Skematisk har vi her en
Analogi mellem to Rækker, en partiel Identitetsrække (\textit{A} → \textit{B} →
\textit{C} → \textit{D} \ldots), i hvilken hvert Led forholder sig til det
følgende som Grund til Følge, og en ubestemt varierende Forskelsrække
(\textit{a}\textsubscript{μ} \textit{b}\textsubscript{γ}
\textit{c}\textsubscript{π} \textit{d} \ldots), hvor Leddenes Rækkefølge blot
er kronologisk bestemt. Fra et Led i den første Række kan man slutte til et
bestemt Led i den anden Række, og derigennem indenfor den anden Række fra
% --- p. 90 ---
\setcounter{footnote}{0}%
\opage{90}et Led til et følgende, selv om man derved kommer udenfor de faktisk
foreliggende Led. En Aarsag er et Emne, der optræder for vor Iagttagelse paa et
vist Tidspunkt, og fra hvilket der kan sluttes til, at et vist andet Emne vil
optræde paa et senere Tidspunkt. Det var gennem sit Studium af Newton's Fysik
og under Indflydelse af den Brod, Hume's Kritik af Aarsagsbegrebet havde
efterladt, at \textsc{KANT} naaede til denne Opfattelse%
% footnote 1 on p. 90
\footnote[1]{De klareste Steder hos \textsc{Kant} er \textit{Kritik der reinen
Vernunft}\textit{1} p. 112 og \textit{Prolegomma,} § 30.}. Formedelst Analogien
mellem et logisk-matematisk Forhold og et Successionsforhold bliver Erfaring
andet og mere end en Opsamling af enkelte Iagttagelser. Dette nye
Erfaringsbegreb stiller \textsc{Kant} i Modsætning til Hume's Opfattelse af
Erfaring som Sum af isolerede, indbyrdes uafhængige Iagttagelser, altsaa som en
kaotisk Forskelsrække, Spinoza's experientia vaga. Men hans Fremstilling
trænger i høj Grad til nærmere Bestemmelse og Simplificering%
% footnote 2 on p. 90
\footnote[2]{I tredie Udgave af \textit{Den nyere Filosofis Historie} har jeg
allerede (III, p. 278 f.) begrundet Nødvendigheden af en Simplificering af
\textsc{Kant's} Fremstilling. Men dér holdt jeg mig mest til den rent formelle
Side.}, før den kan optages som Led i vor Tids Erkendelsesteori.

a. Selve Aarsagsforholdet drøfter \textsc{Kant} i det Afsnit af »Kritik der
reinen Vernunft«, der har til Overskrift »Erfaringens Analogier«. Men forud
gaar der to Afsnit, som benævnes »Anskuelsens Axiomer« og »Iagttagelsens
Antecipationer«, og her synes \textsc{Kant} ikke at mene at have Brug for
Begrebet Analogi. --- »Anskuelsens Axiomer« slaar fast, at alle Fænomener har
Udstrækning i Rum eller Tid eller begge Dele. Hvad \textsc{Kant} her omtaler,
er ikke andet end en Anvendelse af, hvad der ovenfor i denne Afhandling er
kaldt en identisk Kvalitetsrække, som kan konstrueres ad logisk-matematisk Vej,
og hvorved Begreberne Tal, Sted
% --- p. 91 ---
\setcounter{footnote}{0}%
\opage{91}og Tid faar en ideal Form, idet der er Forskelsidentitet mellem
Leddene i Rækken, --- mellem Tallene, Stederne og Øjeblikkene hver for sig. Men
Udtrykket Analogi passer her bedre end Udtrykket Axiom. Meningen er, at ethvert
Emne maa kunne henføres til eller indtræde som bestemt Led i en af disse
Rækker. Trods alle specielle Kvalitetsforskelligheder gælder disse Rækker for
alt, hvad vi skal have streng Erfaring om. Men mellem de specielle
Kvalitetsforskelligheders Rækker og disse Rækker er der kun Analogi, ikke
Identitet. --- Og paa tilsvarende Maade forholder det sig med »Iagttagelsens
Antecipationer«. Her er der Tale om Gradernes Række. Ethvert Emne, der skal
kunne være Genstand for streng Erfaring, maa kunne indgaa som Led i en Række,
indenfor hvilken hvert Led, hvad Intensitet angaar, stedse har samme Forskel
fra sin Forgænger, altsaa en fremskridende Forskelsrække, ligesom Tallenes,
Stedernes og Tidernes Række. Da de faktiske Emner frembyder mange andre
Egenskaber end Intensitet, kan Forholdet mellem Graderækken og de faktiske
Emners Række kun være Analogi. Men det er træffende, naar \textsc{Kant} her
taler om Iagttagelsens Antecipationer; kun burde han have brugt dette Udtryk om
alle sine erkendelsesteoretiske Principer. Analogi kan ikke føre til mere end
en begrundet Forventning (ɔ: en Hypotese), der saa muligvis bekræftes ved nye
Iagttagelser og Forsøg.

b. \textsc{Kant} opstiller tre Erfaringsanalogier. --- Den vigtigste er den
anden, som angaar Aarsagsforholdet i sin Almindelighed. Her vises, som allerede
omtalt, at et Aarsagsforhold er et Successionsforhold mellem to Emner, der
frembyder Analogi med Forholdet mellem Grund og Følge. --- Den første Analogi
hævder, at der ved et Aarsagsforhold ingen
% --- p. 92 ---
\setcounter{footnote}{0}%
\opage{92}Forøgelse eller Formindskelse sker af, hvad der bestaar. Det er
»Substansens Analogi«. Hvis den overhovedet skulde opstilles, burde den være
stillet efter den almindelige Aarsagssætning, som den tydelig nok forudsætter.
Hvad den forlanger, er, at det Forhold mellem Grund og Følge, Aarsagssætningen
forudsætter, skal kunne vendes om, saa at Følgen kan blive Grund, og Grunden
Følge, hvortil svarer, at hvad der først var Virkning, kan være Aarsag til det,
der før var Aarsag, og omvendt, saa at der er et Ækvivalensforhold mellem
Aarsag og Virkning. Men en saadan Sætning kan kun opstilles som Antecipation og
som Ideal. Videnskabens Historie viser, at man overalt, hvor det er muligt,
søger at paavise et saadant Ækvivalensforhold. Derved nærmer man sig et rent
Identitetsforhold i den Grad, det er muligt, naar man har med kvalitativt
forskellige Emner at gøre, og man ser da tillige bort fra den Tidsforskel, der
i hvert enkelt Tilfælde er tilstede og kan være af afgørende Betydning for den
Retning, i hvilken de følgende Forandringer sker. Men det er jo ikke alle
Slutninger, der kan vendes om, saa lidt som alle Prædikater kan gøres til
Subjekter for deres tidligere Subjekter som Prædikater. Hvis \textsc{KANT'S}
»Substansanalogi« gælder, er den i hvert Tilfælde et særlig lykkeligt Tilfælde
af Kausalanalogien. Og det er et uheldigt Udtryk, \textsc{KANT} bruger, naar
han her taler om Substans som særlig Kategori. Han tager ikke Ordet Substans i
den gamle, metafysiske Betydning, men i relativ Betydning, som han selv
udtrykkelig erklærer (Kritik der reinen Vernunft\textsuperscript{1}, p. 525):
»Was in der Erscheinung Substanz heisst, ist nicht absolutes Subjekt, sondern
beharrliches Bild der Sinnlichkeit und nichts als Anschauung, in der überall
nichts Unbedingtes angetroffen wird«. Substansbegrebet i den gamle, absolute
Betydning,
% --- p. 93 ---
\setcounter{footnote}{0}%
\opage{93}er en hendøende Kategori%
% footnote 1 on p. 93
\footnote[1]{Smlgn. \textit{Den menneskelige Tanke,} p. 217 f. --- Om Spinoza's
Substansbegreb se \textit{Spinoza's Ethica,} p. 6; 19---20. Smlgn. desuden
ovenfor § 7 i Slutningen.}. --- Hvad den tredie Analogi angaar, der hævder, at
alt samtidigt maa staa i Vexelvirkning, er Sagen ganske simpel, thi
Vexelvirkning forudsætter Kausalitet. Skal vi ud fra et givet Emne forstaa den
samtidige Existens af et andet Emne, maa der igennem flere eller færre
Mellemled være et Kausalitetsforhold mellem dem eller mellem deres Aarsager. Vi
maa ud fra det ene Emne ved fornøden Orientering træffe paa det andet.

c. Alle tre Analogier tilsammen udsiger, ifølge \textsc{Kant,} at alle Emner
ligger i én Natur. Ved »Natur« forstaar \textsc{Kant} netop »Emnernes
Sammenhæng, hvad deres Existens angaar, efter nødvendige Regler, d. v. s. efter
Love«. (Kritik der reinen Vernunft\textsuperscript{1}, p. 216). \textsc{Kant's}
Naturbegreb, der falder sammen med det strenge Erfaringsbegreb, hævder altsaa
Analogi mellem de successive og samtidige Emners Rækker og Rækker af Grund og
Følge og af Tal, Tider, Steder og Grader. Skal Erfaring være mulig, maa Emnerne
kunne ordnes i Overensstemmelse med saadanne Rækker.

d. Medens \textsc{Kant} omhyggelig konstruerer de »Skematismer«, d. v. s. de
fremskridende Forskelsrækker, hvis Anvendelse paa foreliggende Emner skal gøre
Erfaring i streng Forstand mulig, og medens han omhyggelig skelner mellem den
Analogi, som herved spiller en Rolle, og Identitet, har han ikke set, at kun en
fremskridende, tilnærmende Anvendelse af denne Analogi er mulig. Hans Begreb om
Natur eller streng Erfaring er rigtig dannet, men han saa ikke, at »Erfaring« i
denne Betydning er et Ideal, en stor Hypotese, som egentlig aldrig fuldt ud kan
blive bekræftet. Det gælder ogsaa her, at en Hypoteses Held ikke beviser dens
% --- p. 94 ---
\setcounter{footnote}{0}%
\opage{94}Sandhed. Det store Spørgsmaal (som Hume opkastede, og som
\textsc{Kant} mente at have besvaret ved sin Definition af Erfaring) er, om
alle Emner, der nu og i Fremtiden kan opfattes, kan tillade en Anvendelse af
hine ideale Rækker. \textsc{Kant's} »Erfaringsanalogier« formulerer paa genial
Maade den store Hypotese, Naturvidenskaben i tre hundrede Aar har arbejdet
efter, og som atter og atter har vist sig frugtbar. Og selv om han ikke har set
den hypotetiske Karakter af sit Erfaringsbegreb, saa var det dog en stor
Bedrift, at han saa klart saa og formulerede det Princip, der herligger til
Grund. Derfor betegner han stadig et Hovedpunkt og et Midtpunkt i den nyere
Tids Erkendelsesteori.

e. Det var intet Under, at en Opfattelse som \textsc{Kant's} ikke strax kunde
trænge igennem. Tværtimod opløstes den Sammenhæng mellem formal og real
Erkendelse, han havde søgt at paavise. Tendensen til at naa absolut Identitet
som paa én Gang højeste Tankeform og højeste Tankeresultat, Ønsket om én Gang
for alle at naa ud over alt kaotisk og mangfoldigt til Hvile i Tanken om noget,
der ikke igen viser ud over sig selv, var --- særlig under Indflydelse af
poetiske og religiøse Strømninger i Tidens Aandsliv --- for stærk til, at man
kunde lade sig nøje med \textsc{Kant's} tilsyneladende beskedne Analogier.
Realiteterne skulde gaa helt op i »Ideen«, i Tankens rene Form. I den tyske
romantiske Spekulation fremtræder denne Tendens. \textsc{Goethe} og
\textsc{Schelling} hævdede, at Forklaringen af Kvaliteterne skulde findes
indenfor selve Kvaliteternes Verden, idet visse Kvaliteter lagdes til Grund som
»Urfænomener«. \textsc{Hegel} forsøgte at udlede Kvaliteternes og Tidens Verden
af de rene logiske Identiteters Verden, et Forsøg, der paa én Gang, ved sin
Tillid til den rene Tankes Magt, bar Romantikens Præg, og pegede ud over
Romantiken, idet \textsc{Hegel} anerkendte, at
% --- p. 95 ---
\setcounter{footnote}{0}%
\opage{95}fuld Forstaaelse ikke kan naas ved at blive staaende ved Kvaliteterne%
% footnote 1 on p. 95
\footnote[1]{I sin aandfulde Karakteristik af Hegel's Naturfilosofi lægger
\textsc{émile meyerson} et lignende Synspunkt til Grund, som det her anvendte.
Smlgn. \textit{De l'Explication dans les Sciences,} II, p. 78: »Hegel a
réellement tenté une entreprise qui constitue un postulat éternel de I'esprit
humain«.}. Enhver Videnskab, der vil give absolute Løsninger, maa ende i et
Forsøg som \textsc{Hegel's.} Hvis de har Ret, der hævder Naturvidenskaben som
den absolute Videnskab, maa de gøre Forsøget paa at udlede de kvalitative
Forskelligheder af de logisk-matematiske Principer, ved hvis Hjælp
Kvaliteternes Verden bliver forstaaelig for os. Hvis man ikke vil anerkende, at
streng Realvidenskab naar sit Maal gennem exakte Analogier mellem
Kvalitetsrækker og matematisk-logiske Rækker, maa man gøre \textsc{Hegel's}
Forsøg om og krone den moderne Naturvidenskabs Stræben efter at forklare alle
Forandringer som Bevægelser, hvis Love kan formuleres som Ligninger, med et
Forsøg paa at udlede Kvalitetsforskellighederne af de i Ligningerne til Grund
liggende Identiteter. Da vilde der være givet et Bevis for, at Naturvidenskaben
har løst Verdensgaaden. Den »parfait retour«, som Leibniz krævede, vilde være
naaet. Et saadant Ideal staar sikkert i Baggrunden for mange Naturforskeres
Tankegang. Spørger man om, hvad der da bliver af Kvaliteterne, som dog er det
stadige og umiddelbart givne, gaar Svaret ud paa, at de kun er
»Antropomorfismer« --- det samme Svar, Teologerne giver, naar der spørges,
hvorledes de psykologiske Udtryk, der bruges om Guddommens Væsen og Virken,
egentlig er at forstaa. Men i begge Tilfælde stiller »Antropomorfismen« et nyt
Problem. Kvaliteter som Farve, Klang o. s. v. og som Tænken og Villen er nu
engang Emner, der foreligger for vor Erkendelse, og Spørgsmaalet er, om de
nogensinde kan absorberes, gaa fuldstændig op i rationel Viden (hvad enten
% --- p. 96 ---
\setcounter{footnote}{0}%
\opage{96}denne fremtræder som en Række af Ligninger eller som en Række af
dialektiske Overgange). Vi kan, som Naturvidenskabens Historie paa glimrende og
enestaaende Maade viser, opbygge en Tankernes Verden, der kan vejlede os i den
»Antropomorfismernes« Verden, som vi lever i. Og denne Vejledning giver
Naturvidenskaben os, hvis den i det foregaaende fremsatte Betragtning er
rigtig, ved Hjælp af exakt gennemførte Analogier. --- Om Teologien kan give os
en tilsvarende Vejledning, lader jeg her staa hen. Forsøg paa en streng
Gennemførelse af de anvendte Analogier fører i helt anden Retning end den
positive Teologi%
% footnote 1 on p. 96
\footnote[1]{Smlgn. min \textit{Religionsfilosofi,} § 16---25.}.

Umuligheden af at udlede Kvalitetsforskelligheder og Kvalitetsforandringer af
de Ligninger, der betegner det formale Grundlag for Naturvidenskabens exakte
Analogier, og i hvilke der er en vis Tilbøjelighed til at se den sande
Realitet, ses allerede deraf, at Dannelsen af de identiske Kvalitetsrækker
(Tallenes, Tidernes, Stedernes og Gradernes Rækker), sker ved en analyserende
Udrensning af de empiriske Kvalitetsrækker. Af de foreliggende konkrete
Helheder tager man de formale Elementer ud, som tillader streng
logisk-matematisk Behandling. I den strenge Videnskabs exakte Resultater er det
stedse kun en Del af Virkeligheden, der udtrykkes, selv om det er den Del af
Virkeligheden, der ad Analogiens Vej kan føre til at forstaa og at beherske de
øvrige Dele eller Sider af Virkeligheden. Det er stedse Naturen, der forklares
ved Naturen selv, idet de mest gennemskuelige Forhold lægges til Grund for
Tydningen af andre Forhold. Eller, for at benytte \textsc{Kant's} Definition af
»Natur«: det, der gør Naturen til »Natur«, benyttes til at tyde de Dele af
Naturen, som endnu ikke fyldestgør det strenge Naturbegreb. I Modsætning til den
% --- p. 97 ---
\setcounter{footnote}{0}%
\opage{97}primitive Mentalitet er nu Tydningsgrundlaget draget bestemt frem, og
det er ikke længer nogen fremmed Myndighed, der bevidst eller ubevidst
paakaldes til Tydning. Men ligesaa lidt som Kosmologien (Metafysiken), der
stedse vil forklare hele Tilværelsen ud fra en Del af den (ud fra »Materien«,
ud fra »Aanden« o. s. v.), nogensinde har kunnet naa til rationel Afslutning,
ligesaalidt vil nogen Videnskab formaa det.

I den nyeste Tid har Marburgerskolen opfattet Elimination af de kvalitative
Elementer som den videnskabelige Erkendelses Maal, idet den kun betragter
kvantitativ Videnskab som sand Videnskab. Dens Forsøg er et Sidestykke til
Hegel's, kun at et System af Ligninger træder istedetfor spekulativ Dialektik%
% footnote 1 on p. 97
\footnote[1]{Smlgn. \textsc{Gustav Körber:} \textit{Der Marburger Logismus und
sein Verhältnis zu Hegel.} (Archiv fur die Geschichte der Philosophie. 1917).}.
I min Afhandling »Totalitet som Kategori« (p. 46---48) har jeg haft Lejlighed
til at angive min Stilling til denne Retning, der er en af de mest energiske og
mest betydningsfulde, som har udviklet sig efter Genoptagelsen af
Erkendelsesproblemet i sidste Halvdel af det 19. Aarhundrede.

Hos andre Erkendelsesteoretikere er Kant's Tanke kommen mere til sin Ret.
Positivismens Grundlægger \textsc{Auguste Comte} lagde stor Vægt paa Forskellen
mellem de forskellige Videnskaber, som han ordnede i en Række, der begyndte med
Matematiken, og i hvis følgende Led specielle og komplexe Kvaliteter spillede
en større og større Rolle. Sidste Led i denne Række er Sociologien. Mellem de
forskellige Videnskaber kan der kun være Analogi, ikke Identitet, og
\textsc{Comte} miskendte endog, som vi har haft Lejlighed til at se, saadanne
Reduktioner, som var videnskabelig berettigede. Men hans Klassifikation af
Videnskaberne, der
% --- p. 98 ---
\setcounter{footnote}{0}%
\opage{98}skrider frem fra det simple til det komplexe, følger en lignende
Ordning, som den, der kort efter af \textsc{Renouvier} anvendtes i Kategorilæren%
% footnote 1 on p. 98
\footnote[1]{Smlgn. \textit{Relation som Kategori,} p. 13.}, og som ogsaa er
fulgt i min Fremstilling. Det gælder for Kategorier som for Videnskaber, at de
mere komplexe ikke kan afledes af mere elementære, men at der dog kan bestaa
lærerige Analogier mellem dem. Med fuld Bevidsthed og stor Klarhed blev ved
Midten at det 19. Aarhundrede Analogisynspunktet hævdet af \textsc{Clerk
Maxwell,} som vi allerede har haft Lejlighed til at se, hvad Forholdet mellem
Aritmetik og Geometri angaar. Analogi udtrykker overhovedet for
\textsc{Maxwell} Forholdet mellem formal og real Videnskab. I Tilslutning til
ham arbejdede en anden filosoferende Naturforsker, \textsc{Ernst Mach,} hvis
Erkendelsesteori navnlig er bygget paa et selvstændigt Studium af
Naturvidenskabens Historie%
% footnote 2 on p. 98
\footnote[2]{\textsc{Maxwell} og \textsc{Mach} har jeg forsøgt at karakterisere
i \textit{Moderne Filosofer,} p. 82---92.}. Af nyere Forskere maa især nævnes
\textsc{Émile Meyerson} i hans allerede flere Gange omtalte Værker. Naar
\textsc{Pierre Boutroux} \textsc{%
% footnote 3 on p. 98
\footnote[3]{\textit{L'idéal scientifique des Mathématiciens,} p. 152; 243.}}
udtaler, at Matematik er en Metode, ikke en objektiv Videnskab, synes en
lignende Opfattelse af Forholdet mellem formal og real Videnskab at ligge til Grund.

Hos fremragende Naturforskere, selvom de ikke som \textsc{Maxwell} og
\textsc{Mach} gaar mere direkte ind paa erkendelsesteoretiske Spørgsmaal, vil
man dog ofte finde Udtalelser, der vidner om fuld Klarhed angaaende Forholdet
mellem Data og deres rationelle Tydning. Ordet Analogi nævnes maaske ikke, men
de Udtryk, der bruges, peger afgjort i Retning af Begrebet Analogi. Allerede i
Titelen paa \textsc{Niels Bohr's} Afhandling »Atomernes Bygning og Stoffernes
fysiske og kemiske Egenskaber« (1922) er dette Begreb egentlig
% --- p. 99 ---
\setcounter{footnote}{0}%
\opage{99}antydet. Stoffernes fysiske og kemiske Egenskaber er det givne. De
kvalitative Forskelligheder, de frembyder, Grundstoffernes Atomvægte, deres
Spektre o. s. v. kan ordnes i Rækker, til hvilke der i Teorien om Atomernes
Bygning svarer Stillinger og Overgange af Elektronerne indenfor Atomet. Hvad
der interesserer Erkendelsesteorien, bortset fra de for objektiv Videnskab
betydningsfulde Resultater, er den Maade, paa hvilken disse Resultater
udtrykkes. Der antydes i Udtrykkene stedse Analogier, ikke Identitetsforhold,
et »betyder«, ikke et »er«, mellem den Kvalitetsrække, der danner Udgangen, og
den Kvantitetsrække, som fører til en bestemt Tesis om Atombygningen. Saaledes
tales der (p. 1) om »Belysning«, (p. 33) om »Forklaring, eller rettere
Forstaaelse«, (p. 36) om »Fortolkning«, og p. 45 hedder det: »som Spektret
viser, og som Atommodellen gør forstaaelig«. Konstruktionen af Atomets indre
Bygning og af, hvad der foregaar indenfor den, ledes Skridt for Skridt af
Iagttagelse af Forandringer i Stoffernes kvalitative Egenskaber. --- Det er
klart, at hvis Grundstofferne ikke lod sig ordne i Rækker, hvor Leddene fulgte
efter hverandre ifølge en vis Lov, --- hvis de altsaa dannede en kaotisk
Forskelsrække, ikke en fremskridende Forskelsrække, --- vilde der ingen
rationel Analogi kunne opstilles.

f. Hvis det er rigtigt, at Forholdet mellem formale og reale Videnskaber er et
Analogiforhold, saa ligger deri, at Fremskridt paa de reale Videnskabers
Omraade ikke kan vindes ved Deduktion ud fra formale Videnskabers Resultater.
Som Metode kan en Analogi kun være en ledende Tanke, en Antecipation, ikke et
Bevis. Den kan føre til at stille Spørgsmaal, og et godt stillet Spørgsmaal vil
i et videnskabeligt Sind føre ind paa nærmere Undersøgelser.

% --- p. 100 ---
\setcounter{footnote}{0}%
\opage{100}Selve vor Erkendelses Grundbegreber, Kategorierne, lægger sig
tydeligst for Dagen i de Spørgsmaal, der stilles%
% footnote 1 on p. 100
\footnote[1]{Se herom \textit{Den menneskelige Tanke,} p. 137 f.}. I de ledende
Tanker, som fører til Spørgsmaal og Undersøgelser, viser sig fremfor alt
Kontinuiteten i Videnskabens Historie. En ny Kreds af saadanne ledende Tanker
dannedes ved den moderne Videnskabs Grundlæggelse for snart 350 Aar siden, og
disse Tanker har bevaret deres Kraft den Dag idag, især efterat det 19.
Aarhundrede med saa stor Energi havde arbejdet paa at sætte Erfaringer i
Analogiernes Sted. Som allerede omtalt i en anden Sammenhæng, falder dette
særlig i Øjnene ved Studiet af de Afhandlinger, i hvilke \textsc{Robert Mayer,
Joule, Colding} og \textsc{Helmholtz} omtrent samtidig, men uafhængig af
hverandre, grundlagde den moderne Lære om Energiens Bestaaen. De erklærede
alle, at de i Grunden ansaa det som selvfølgeligt, at der maatte være
Ækvivalens mellem Aarsag og Virkning, hvad der heldigvis ikke holdt dem fra at
undersøge Forholdet gennem Iagttagelse og Forsøg. Men den Selvfølgelighed, de
antog, er ikke til alle Tider bleven anerkendt. Den er ikke ældre end det 17.
Aarhundrede, hvis Filosofer og Naturforskere hævdede, at der i Virkningen
hverken kunde komme noget nyt til eller gaa noget tabt af, hvad der var i
Aarsagen. \textsc{Huyghens} og \textsc{Leibniz} hævdede denne Sætning i
Mekaniken, efterat \textsc{Descartes} og \textsc{Spinoza} havde hævdet den fra,
hvad vi vilde kalde et erkendelsesteoretisk Synspunkt. Sætningen blev paa
Tidens Vis opstillet rent dogmatisk. \textsc{Spinoza} begrunder den ganske vist
ved, at hvis der var andet og mere i Virkningen end i Aarsagen, saa vilde dette
mere ingen Aarsag have, altsaa staa uforklaret. Men hvem siger, at alt skal
kunne forklares? Det, der skulde vises, var netop, at Forklaring her
% --- p. 101 ---
\setcounter{footnote}{0}%
\opage{101}var mulig, og det er det store hos Energisætningens Opdagere, at de
banede Vejen derfor.

Men selv om Analogien kun har Betydning som Grundlag for Spørgen og Søgen,
eller som heuretisk Princip, er det naturligvis af Betydning, at den fattes i
sin Bestemthed og sin Begrænsning. Vi har set, at det endnu skortede herpaa hos
\textsc{Kant,} naar han mente, at han kunde opstille en Sætning om »Substansens
Bestaaen«, endog før han havde opstillet Aarsagssætningen. Hans geniale
Opdagelse var den, at der mellem Begreberne Grund og Aarsag kun var Analogi,
ikke absolut Identitet. Men hvad selve Begrebet Grund (eller Forholdet mellem
Grund og Følge) angaar, skelnede han ikke mellem Slutninger, der kan vendes om
(saa der er Ækvivalens mellem Grund og Følge), og Slutninger, der ikke kan
vendes om. Al real Videnskab staar og falder med Aarsagssætningen, ifølge
hvilken et Emne forholder sig til et andet Emne som Følge til Grund; men det at
Grund og Følge (Aarsag og Virkning) skulde kunne byttes om, er en speciellere
Antagelse, der ikke kan udledes af den almindelige Aarsagsætning, selv om den
kan siges at formulere Idealet for al real Videnskab. Der kan være Videnskab,
selv om Videnskabens højeste Ideal maaske aldrig naas. --- Hos \textsc{Robert
Mayer} gør sig en lignende Uklarhed gældende. En Ungdomsven, der stod i livlig
Tankeforbindelse med ham, før hans epokegørende Afhandling var forfattet
(1841), siger om ham: »Skønt han var langt fra at være en Skolefilosof, havde
en aldeles selvstændig Eftertanke paa Logikens og Metafysikens Omraade om
Kausalitetens Væsen maaske nok saa stor Andel i hans Opdagelse som exakt
Naturforsken«%
% footnote 1 on p. 101
\footnote[1]{\textsc{Gustav Rümelin:} \textit{Erinnerungen an Robert Mayer}
(Reden und Aufsätze. Neue Folge. 1881, p. 383).}. Men \textsc{Mayer} mente,
% --- p. 102 ---
\setcounter{footnote}{0}%
\opage{102}som det ses af hans grundlæggende Afhandling »Bemerkungen über die
Kräfte der unbelebten Natur«%
% footnote 1 on p. 102
\footnote[1]{Optrykt i \textit{Die Mechanik der Wärme.} Herausg. von Weyrauch,
p. 29.}, at hans nye Sætning »med Nødvendighed« fremgik af Sætningen causa
æqvat effectum, og at der derfor var Ækvivalens mellem Bevægelse og Varme.
\textsc{Mayer} vidste ikke, at Korrektivet til hans Sætning allerede var givet
i en Afhandling af \textsc{Carnot} (1824), der viste, at Overgangen fra Varme
til Bevægelse ikke gaar for sig saa let som Overgangen fra Bevægelse til Varme.
Og han saa ikke klart, at der ved Overgangen fra Logik og Mekanik til speciel
Fysik ikke kan være Tale om Deduktion, men kun om analogisk Genspejling af de
abstrakte Idealer i de konkrete Forhold mellem Kvaliteter. --- Analogiens
metodiske Betydning falder naturligvis ikke bort, fordi den paa mange Maader
maa omformes og begrænses ved selve det Arbejde, der gøres under dens Ledelse. ---

g. Af Analogier med heuretisk Betydning indenfor samme Omraade er især saadanne
af erkendelsesteoretisk Interesse, ved hvilke der gaas ud fra en Ordning
indenfor en vis Totalitet og undersøges, om en lignende Ordning ikke genfindes
indenfor andre Totaliteter. For Galilei indeholdt saaledes Opdagelsen af
Jupiter's Maaner en Analogi til Belysning af Solens Forhold til Planeterne
(Jorden indbefattet) og saaledes til Bestyrkelse i den kopernikanske
Opfattelse. Og i den nyeste Tid har Analogien mellem Solsystemets og Atomets
indre Bygning virket anskueliggørende og fremhjælpende. Saadanne Exempler fører
os over fra Kausalitetsbegrebet til Totalitetsbegrebet. Der ligger i
Tænkningens Natur en Trang til at danne Helheder. Enhver Slutning er en logisk
Helhed, og til Grund for Kausalitetstrangen ligger ligeledes en Trang til at
ophæve det enkelte Emnes Isola%
% --- p. 103 ---
\setcounter{footnote}{0}%
\opage{103}tion ved at bringe det i saa fast Sammenhæng som muligt med andre
Emner. Men ud herfra fører ingen Deduktion til specielle Resultater.
Totalitetsbegrebets Anvendelse beror, ligesom Aarsagsbegrebets, paa en Analogi,
der kan være af stor Betydning som ledende Tanke, men ikke har bevisende Magt.
Som jeg har haft Lejlighed til at udtale i en anden Afhandling (Totalitet som
Kategori. p. 51), kan Begrebet Totalitet aldrig bruges som Forklaring, saaledes
som det sker i vitalistiske, spiritualistiske og mystisk-sociologiske Teorier.
Som alle Kategorier er dette Begreb en Form for Tankearbejde og tillige et
Ideal, der søges genfundet og gennemført i det enkelte. Hvor Iagttagelsen
umiddelbart frembyder Emner med Helhedspræg --- et Atom, en Klode, en
Organisme, en Personlighed, et Samfund --- dér stilles os ligesaa mange
Opgaver, først for Beskrivelse og Klassifikation, saa for Undersøgelse af
Opstaaen og af Betingelser for Bestaaen%
% footnote 1 on p. 103
\footnote[1]{Smlgn. nærmere \textit{Totalitet som Kategori,} p. 52---66.}. Det
er et analytisk Arbejde paa Grundlag af, hvad Iagttagelse kan frembyde, der her
kræves. Som \textsc{Leibniz} har sagt, er det Loven for Elementernes
Sammenhæng, der udgør et Væsens ejendommelige Helhed, og det er en saadan Lov,
det analytiske Arbejde skal søge at finde.

Som det forholder sig med Totalitetsbegrebet, forholder det sig ogsaa med
Udviklingsbegrebet. Ogsaa dette Begreb er en Tankeform, som baade uvilkaarlig
og med fuld Bevidsthed leder Forskningen. Dets metodiske Betydning viser sig i,
at det indeholder en Opfordring til at forsøge at konstruere en Række, i
hvilken hvert Led betinger det følgende og betinges ved det foregaaende, og det
saaledes, at i senere Led er stedse en større Mangfoldighed af Elementer
forbundne paa en inderligere Maade. Det bliver en
% --- p. 104 ---
\setcounter{footnote}{0}%
\opage{104}regelmæssig varierende Forskelsrække, der ikke kan vendes om.
(Smlgn. Den menneskelige Tanke, p. 233 f.). Hvad den enkelte Organismes eller
det enkelte Organs Udvikling angaar, hævder den af \textsc{Caspar Friedrich
Wolff} grundlagte Epigenesislære, at der kun kan være Analogi mellem de
forskellige Stadier, medens tidligere Teorier fandt Identitet her, idet de
ansaa de senere Stadier for fuldstændig præformerede i de tidligere. Samme
Tanke som hos \textsc{Wolff} optraadte omtrent samtidig hos \textsc{Rousseau} i
hans psykologisk-pædagogiske Teori, der hævdede, at Forskellen mellem Barnet og
den voksne ikke blot var en Forskel i Dimensionerne. I det 19. Aarhundrede
forsøgtes Udviklingsbegrebet ogsaa gennemført for de organiske Arters
Vedkommende, og man søgte at paavise en Kontinuitet i Udviklingsprocesserne,
hvorved Analogien mellem de forskellige Stadier fik et realt Grundlag.
\textsc{Cuvier} var endnu tilfreds med, at Analogien havde sin Grund i
Skaberens Fantasi, der efter enhver ved geologisk Katastrofe voldet Undergang
af organisk Liv frembragte nye, analoge Former. Hos \textsc{Goethe} og de
romantiske Naturforskere var Metamorfose og Udvikling kun billedlige Udtryk for
Formernes Analogi; de antog Forskydninger i Naturens systematiske Fantasi, men
ikke nogen real Nedstamning. Mange af \textsc{Darwin's} Modstandere, især
Anatomer og Morfologer, stod paa et saadant Standpunkt. Naar undertiden moderne
strenge Kritikere af Darwinismen erklærer, at de ikke forstaar det Standpunkt,
hine Morfologer indtog, er det et Tegn paa, at Analogien mellem
Udviklingsstadierne i den Grad har vist sig i sin reale Betydning, at Forskere
i Nutiden ikke kan sætte sig ind i deres eget Fags Fortid. Intet viser bedre
Vigtigheden af Studiet af Videnskabens Historie med særligt Henblik til de
Grundformer, indenfor hvilke Forsk%
% --- p. 105 ---
\setcounter{footnote}{0}%
\opage{105}ningen bevæger sig. De, der ledes af visse Tanker, opdager ikke
altid Ledebaandet, og Følgen deraf kan ikke blot blive Mangel af historisk
Forstaaelse, men ogsaa dogmatisk Tillid til deres egen Tids Former.

Indenfor Biologien skelnes der, især siden \textsc{Owen's} Tid, skarpt mellem
Analogi og Homologi%
% footnote 1 on p. 105
\footnote[1]{Smlgn. \textsc{Yves Delage:} \textit{Comparative Anatomy and the
foundation of Morphology.} (Congress of Arts and Sciences. 1904).}, saaledes at
hin beror paa en Overensstemmelse i Funktion, altsaa er af fysiologisk Art,
medens denne beror paa Overensstemmelse i Struktur, og altsaa er af anatomisk
Art. Pattedyrenes Lunger og Fiskenes Svømmeblærer er homologe, ikke analoge;
derimod er Pattedyrenes Lunger og Fiskenes Gæller analoge, ikke homologe. Der
kan staa en Strid mellem Anatomer og Fysiologer, om Analogi eller Homologi er
af størst Betydning. For en Del var den berømte Strid mellem \textsc{Cuvier} og
\textsc{St. Hilaire} af denne Art. \textsc{Cuvier} tænkte væsentlig paa Enhed i
Struktur (unité de composition), medens \textsc{St. Hilaire} ved Enhed forstod
Analogi%
% footnote 2 on p. 105
\footnote[2]{\textsc{Cuvier:} \textit{Histoire des progrès des sciences
naturelles.}Bruxelles. 1838. II, p. 455.}. Dog maa det bemærkes, at
erkendelsesteorisk er Formlighed, Lighed i Struktur ofte at kalde Analogi, da
det mere er Forholdet mellem Delene end kvalitativ Lighed paa enkelte Punkter,
det drejer sig om. Ifølge Epigenesislæren er der, som allerede bemærket, kun
Analogi mellem forskellige Udviklingsstadier af samme Organ, selv om
Udviklingen ses at være gaaet kontinuerlig fra den ene Form over til den anden,
og ifølge \textsc{Cuvier's} egen Fremstilling tænkte \textsc{St. Hilaire}
væsentlig paa Strukturanalogier hos forskellige Arter. Det var altsaa fra denne
Side set en Strid om Formlighed eller Forholdslighed. Som jeg har bemærket i
min Psyko%
% --- p. 106 ---
\setcounter{footnote}{0}%
\opage{106}logi, er Formlighed en Overgangsform mellem Kvalitetslighed og
Forholdslighed.

h. Hvad Forstaaelsen af det organiske Liv angaar, maa, som allerede berørt, den
Mulighed være udelukket, at selve Begrebet Organisme skulde indeholde en
Forklaring. Begreber som Liv og Organisme trænger selv til Forklaring. Da nu
Fysiologiens Historie viser, at alle videnskabelige Forklaringer paa disse
Omraader er vundne gennem Tilbageførelse af organiske Fænomener til fysiske og
kemiske Forhold, kunde det synes at være berettiget at opfatte Organismen helt
igennem i Analogi med Mekanisme. Men der bliver her stadig den Vanskelighed
tilbage, at ved et mekanisk Hele kan vi forud kende Elementerne, ved hvis
Sammenføjning det bliver til, ligesom vi kan se Mursten ligge ophobede, før de
forbindes til en Bygning. Et saa udvortes Forhold er der ikke mellem Organismen
og dens Elementer. Man kan genfinde forud kendte Grundstoffer ved Analyse af
Organismen og dens Processer, men ikke forklare disse som blot Produkt af
saadanne Elementer. Som en biologisk Forsker træffende har sagt: i Biologien
fremtræder Helheden som given, og vi gaar fra Helhedsformen tilbage til de
Elementer, den viser sig at bestaa af, men i Fysik og Kemi har vi omvendt
Elementerne givne, og det gælder at forstaa deres Forbindelse til Helheden;
Atomet er ikke givet, men konstrueres ud fra de materielle Processer, der kan
paavises%
% footnote 1 on p. 106
\footnote[1]{\textsc{Johan Hjort:} \textit{The Unity of Science,} p. 166.}. ---
Overfor Analogien med Mekanismen har man ofte stillet Analogien med menneskelig
Frembringen, snart paa mere udvortes Maade, idet man benyttede Sammenligningen
med en Bygmesters Forhold til Bygningen, han opfører, snart mere aandfuldt,
idet man benyttede Sammenligningen med en genial Arkitekt
% --- p. 107 ---
\setcounter{footnote}{0}%
\opage{107}og den Plan, der bliver til i hans skabende Fantasi. Men i begge
Tilfælde skelner man saa mellem Skaber og Værk, og her glipper Analogien. ---
\textsc{Kant's} Resultat i hans Undersøgelse af dette Problem staar fast endnu
den Dag idag, naar han udtaler: »Genau zu reden, hat die Organisation der Natur
nichts Analogisches mit irgend einer Causalität, die wir kemien« (Kritik der
Urtheilskraft § 65). Det udelukker nu ikke, at Analogiens ansporende Virken
ogsaa her kan gøre sig gældende i Forskningen. Begge de anførte Analogier kan
nemlig anvendes heuretisk, og anvendes faktisk af forskellige Forskere, idet
der snart stiller sig det Spørgsmaal, hvilke fysiske og kemiske Betingelser, en
organisk Proces forudsætter, snart spørges om, hvorvidt en Organisme er
saaledes udstyret, at den kan hævde sin Bestaaen under bestemte Forhold, især
naar disse er Forandring underkastede. Fremtiden maa vise, om denne Biologiens
Stilling er den endelige.

i. I Spørgsmaalet om Analogiens Betydning i videnskabelig Forstand gør der sig
en interessant Modsætning gældende mellem to udmærkede franske
Erkendelsesteoretikere. \textsc{Émile Meyerson} betragter, som vi tidligere har
set, Analogi som den sidste Form for videnskabelig Tænkning over foreliggende
Emner. Efter hans Opfattelse finder Menneskets rationelle Trang ved
Gennemførelse af videnskabelig analyserede Analogier den Tilfredsstillelse, som
overhovedet er mulig. Men i de irreduktible Kvaliteter bliver der stedse
irrationelle Elementer tilbage%
% footnote 1 on p. 107
\footnote[1]{\textit{De l}\textit{'}\textit{Explication dans les Sciences.}I,
p. 217---224; II, p. 291---295 og \textit{Bulletin de la Société française de
Philosophie.} 24. Febr. 1921, p. 59.}. Imod denne Opfattelse indvendte
\textsc{Léon Brunschwicg,} at den lægger et falsk Erkendelsesideal til Grund.
Hvis sand Videnskab skulde være Gennemførelse af absolut Identitet som Er%
% --- p. 108 ---
\setcounter{footnote}{0}%
\opage{108}kendelsens højeste Form, vilde Opnaaelsen af dette Ideal være et med
intellektuelt Selvmord. Saasnart man opgiver dette Ideal, falder den formodede
Irrationalitet bort, da den blot betegner Modsætning til absolut Identitet. Den
sande Videnskab søger Synspunkter, der gør det muligt at finde Ækvivalens
mellem komplexe Emner og deres Elementer; men det kan ikke opstilles som dens
Ideal at forfalske de foreliggende Kvaliteter ved at forvandle dem til rene
Kvantiteter%
% footnote 1 on p. 108
\footnote[1]{\textit{Bulletin.} 1921, p. 59 f. \textit{L'Expérience humaine,}
p. 581.}.

For mit Vedkommende har jeg allerede i et Foredrag, jeg holdt i Oxford 1904, og
som er trykt i »Mind« 1905 og i mine »Mindre Arbejder« (anden Række), bekendt
mig til Analogisynspunktet, og i »Den menneskelige Tanke« (p. 181---201 og
309---318) har jeg genoptaget denne Betragtning. Brydningen mellem Identitet og
Analogi gaar gennem hele Filosofiens Historie. Det har kostet og koster
Resignation at se en saavidt muligt gennemført Analogi som det højeste, der kan
naas, og som den eneste Maade, paa hvilken de forskellige Erkendelsesveje kan
mødes, og paa hvilken Forholdet mellem de forskellige Erkendelsesomraader kan
bestemmes. Men da Analogi kun er Forholdslighed, ikke Identitet, bliver der
stedse et Irrationelt tilbage, tilmed da det ikke paa Forhaand kan vides, om de
Analogier, der nu raades over, kan gennemføres for nye Emners Vedkommende,
eller om det er muligt at finde nye Analogier, naar uforudsete Emner kræver
saadanne. Ud over exakte Analogier vilde der kun kunne naas, dersom det kunde
vises, at de Slutninger, i Analogi med hvilke vi opfatter de kvalitative og
temporale Forskelligheder, kunde vendes om, saa at man kunde være sikker paa,
at ingen andre Præmisser end de, paa hvilke vore
% --- p. 109 ---
\setcounter{footnote}{0}%
\opage{109}Videnskaber bygger, kunde føre til Resultater, der stemmer med
Iagttagelserne.

\begin{center}19. Naturvidenskab og Aandsvidenskab.\end{center}

Ved Belysningen af Forholdet mellem formal og real Videnskab er i det
foregaaende Naturvidenskaben bleven brugt som Exempel paa Realvidenskab, fordi
den uimodsigelig er den mest fuldkomne Realvidenskab. Aandsvidenskaben
(Filosofi, Sprogvidenskab, historisk Forsken) arbejder stadig med Kvaliteter og
Successioner, som det ikke er lykkedes at føre tilbage til saa klare Forhold,
som det for en væsentlig Del er lykkedes Naturvidenskaben ad exakte Analogiers
Vej at paavise for sine Emners Vedkommende. Men saa gælder det om at undersøge
Forholdet mellem Naturvidenskab og Aandsvidenskab, særlig med Hensyn til, om
dette Forhold er et Identitetsforhold eller et Analogiforhold.

Man kunde spørge, om der i Grunden gives nogen anden Videnskab end
Aandsvidenskab. Hvad vi kalder Naturvidenskab, er jo dog Frugt af aandeligt
Arbejde, og dens Metoder saavel som de Resultater, der naas ved Hjælp af disse
Metoder, maa være bestemte ved det aandelige Arbejdes Væsen. Kan
Naturvidenskaben nogensinde blive andet og mere end et stort Forsøg paa at
opfatte Tilværelsen saaledes, som aandeligt Arbejde formaar det?
Naturvidenskaben er forsaavidt kun en særegen Form for Aandsvidenskab, og
Forskellen er kun den, at Aandsvidenskaben foruden at undersøge det aandelige
Arbejde, der fører til Videnskab, ogsaa undersøger andre Former for aandeligt
Arbejde, f. Ex. saadanne, som fører til Kunst, Moral og Religion. Af særlig
Interesse er det dog i denne Sammenhæng, at selve Erkendelsesteorien utvivlsomt
er en Aandsvidenskab, hvis Emne er den Evne til Arbejde, den Energi,
% --- p. 110 ---
\setcounter{footnote}{0}%
\opage{110}der lægger sig for Dagen i, at vore Forestillinger undergaar
Forandringer, der paa én Gang fører til større Sammenhæng mellem dem og til
nøjagtigere Bestemmelse af enhver af dem. Den erkendelsesteoretiske
Betragtningsmaade giver os et Enhedssynspunkt for al Videnskab, maaske det
eneste, der gives. Langt fra at dække over Forskellighederne mellem
Videnskaberne skærper et saadant Enhedssynspunkt netop Blikket for disse
Forskelligheder. De forskellige Arbejdsmaader og det forskellige Held, hvormed
der arbejdes paa de forskellige Arbejdsomraader, vinder netop i Interesse ud
fra et saadant Synspunkt. Og det bliver nu ogsaa lettere muligt at undersøge
Forholdet mellem videnskabeligt Arbejde og aandeligt Arbejde, der øves paa
andre Omraader af menneskeligt Sjæleliv end det intellektuelle. Ogsaa her vil
Enhedssynspunktet lade de ejendommelige Forskelligheder fremtræde med større
Tydelighed.

Men selv om en Betragtning som den nu antydede skulde være den fundamentale og
den definitive, stiller Sagen sig dog efter manges Mening anderledes, naar vi
retter Opmærksomheden mod de forskellige Omraader for intellektuelt Arbejde.
Naar det nu viser sig, at det naturvidenskabelige Arbejde, hvis Omraade er de
saakaldte materielle Fænomener, hvortil regnes alt, hvad der fremtræder for os
i Rummets Form, foregaar efter fuldkomnere Metoder og følges af større Held end
det aandsvidenskabelige, som ikke kommer ud over Kvalitets- og
Tidsforskelligheder, --- var det saa ikke rimeligt at gøre alt intellektuelt
Arbejde til naturvidenskabeligt Arbejde, --- at lade den fuldkomnere Viden
belyse den ufuldkomnere Viden og optage den i sig derved, at man opfattede alle
saakaldte aandelige Emner som Led i den store, kontinuerlige Sam%
% --- p. 111 ---
\setcounter{footnote}{0}%
\opage{111}menhæng mellem alle saakaldte materielle Emner, som Naturvidenskaben
mere og mere afslører for os? Alle indrømmer (hvorledes de saa ellers opfatter
Forholdet mellem »Sjæl og Legeme«), at aandelige Emner, hvilke Beskaffenheder
de end har, hænger sammen med visse fysiologiske Processer, og naar nu disse
sidste afgiort %
% PRINTED AS IS: „afgiort“ (p. 111)
falder ind under Naturvidenskabens Omraade, var det saa ikke rigtigst at tage
sit Udgangspunkt her og betragte den naturvidenskabelige Opfattelse som den
eneste gyldige, og saa opfatte Aandsvidenskaben (forsaavidt den overhovedet kan
kaldes Videnskab) som en Del af Naturvidenskaben?

Naar jeg forsøger at besvare dette Spørgsmaal, maa jeg først henvise til, hvad
jeg i forudgaaende Afsnit af denne Afhandling har forsøgt at vise, at al real
Videnskab bygger paa Analogi med logisk og matematisk Tænken, saa at
Spørgsmaalet tilsidst bliver om denne Tænkens Gyldighed. Er Slutninger og
Ligninger aandelige eller materielle? Hvis de maa siges at være Former for og
Resultater af intellektuelt Arbejde, maa der spørges, hvem der drager
Slutningerne og begrunder Ligningerne. Vi staar herved overfor Forholdet mellem
Erkendelsesteori og Psykologi og maa dvæle lidt ved det, før vi gaar over til
Forholdet mellem Psykologi og Fysiologi, der mere umiddelbart vedkommer vort Problem.

a. Hvad det første Forhold angaar, maa jeg henvise til min Afhandling
»Totalitet som Kategori« (p. 4\textit{---}18; 46---48), hvor jeg har søgt at
vise, at den Omstændighed, at Psykologiens Mulighed er et af
Erkendelsesteoriens Emner, ikke ophæver Psykologiens Selvstændighed, ligesaa
lidt som Fysikens Selvstændighed ophæves derved, at ogsaa den er et saadant
Emne. Det erkendelsesteoretiske Subjekt, som forudsættes ved alle Slutninger og
Ligninger, er nu et af
% --- p. 112 ---
\setcounter{footnote}{0}%
\opage{112}Psykologiens interessanteste Emner. Ved Hjælp af Videnskabens
Historie søger Psykologien at vise, hvorledes dette Subjekt har udviklet sig
fra at være blot iagttagende og beskrivende til et saa højt Stade. at man har
kunnet sige, at det er det, der foreskriver Naturen dens Love. Det er i stadig
Udvikling; til hver Tid er det paa et bestemt Trin af denne Udvikling. Det
existerer ikke i og for sig; det kan man kun antage ved et Spring fra
Erkendelsesteori til Metafysik. Og paa ethvert af dets Udviklingstrin stiller
det Psykologien den Opgave at paavise, hvorledes dets Udvikling overhovedet har
været mulig, og det af den gode Grund, at det paa ethvert Trin frembyder
specielt psykologiske, ikke blot rent logiske og matematiske Egenskaber. Det er
Principerne for aandelig Væxt, Psykologien her anvender paa det
Erkendelsessubjekt, al Videnskab forudsætter. Der er Erkendelsesteoretikere,
der helst vil smøge alt Psykologisk fra sig; de rammes af den Ironi, hvormed
\textsc{Kierkegaard} omtaler sin Tids spekulative Filosofer. Den i andre
Henseender saa fortjenstfulde Marburgerskole rammes af saadan Ironi.
Naturligvis kan Psykologien ikke bygge paa andre Principer end dem, der
anerkendes af Erkendelsesteorien, men det har den fælles med ---
Erkendelsesteorien selv, der naturligvis ikke kan bygge paa andre Principer end
saadanne, den kan anerkende, naar den selv engang bliver færdig. (Smlgn.
ovenfor § 7).

b. Hvad Forholdet mellem Psykologi og Fysiologi angaar, kan jeg ogsaa henvise
til tidligere Udtalelser, dels i min »Psykologi« og mine »Psykologiske
Undersøgelser«, dels og især til mit Universitetsprogram (»Filosofiske
Problemer«) (1902), hvis Indhold i forkortet Form er optaget i »Den
menneskelige Tanke« (1910) (p. 22---34; 329---338). Jeg vil i denne Sammenhæng
kun dvæle ved enkelte
% --- p. 113 ---
\setcounter{footnote}{0}%
\opage{113}Punkter. En iagttagende (henholdsvis experimenterende) og
beskrivende Psykologi er Forudsætningen for, at den fysiologiske Sammenhæng,
der efter nogles Mening%
% footnote 1 on p. 113
\footnote[1]{I denne Retning gaar indenfor vor Literatur den skarpsindige
Afhandling af \textsc{Jørgen Fr. Jørgensen:} \textit{Problemet om Forholdet
mellem Sjæl og Legeme.} (Tidsskriftet »Vor Tid«. 1913.)} skal give os den
eneste videnskabelige Forklaring af Sjælelivet, kan findes. For at »forklare«
noget, maa jeg først vide, hvad det er, der skal forklares; ellers kan jeg ikke
vide, om den rette Forklaring er funden. Den saakaldte Psykofysiologi
forudsætter derfor stedse den egentlige Psykologi. Programmet om at føre al
Psykologi tilbage til Fysiologi kan kun gennemføres, naar den
empirisk-experimentale Psykologi først har gjort sit Arbejde og derved
præciseret de Fænomener, der skal forklares. Psykologien har paa sit eget
Omraade store Vanskeligheder at gennemføre, og det nytter ikke at gaa forhastet
frem. De sjælelige Fænomeners Diskontinuitet, deres kvalitative Forskelligheder
og deres stadige Skiften, saavel som deres mangfoldige og komplicerede
Sammenvævninger stiller stadig nye Opgaver for beskrivende, genetisk og
analyserende Psykologi.

Ser vi nu paa de fysiologiske »Forklaringer«, der skalimponere Psykologen ved
deres Exakthed, viser det sig, at de er vundne ved Hjælp af Analogier ud fra
direkte psykologiske Iagttagelser. De er Forsøg paa at konstruere en
fysiologisk Mekanik, der nogenlunde svarer til, hvad psykologisk Iagttagelse
viser. Det er et psykologisk Faktum, at jo hyppigere to Forestillinger følges
ad, des lettere sker Overgangen senere fra den ene til den anden. Dette skal
finde sin Forklaring ved Nervecellers og Nervetraades »Balling«. Men i
Virkeligheden véd man ikke mere om, hvad saadan Nervebaning egentlig er, end
hvad der ligger i det
% --- p. 114 ---
\setcounter{footnote}{0}%
\opage{114}Faktum (som baade er et psykologisk og et fysiologisk Faktum), at en
Funktion gaar lettere for sig ved Gentagelse. Man har faaet et mekanisk
Billede; det er det hele. Og et saadant Billede kunde meget godt direkte være
anvendt i selve Psykologien, da alle psykologiske Udtryk etymologisk viser sig
at være hentede fra Yderverdenen.

Et af de første og mærkeligste Forsøg paa at opfatte Psykologi som fysiologisk
Mekanik er gjort af \textsc{Thomas Hobbes,} der vil hævde, at Sansning og
Tænkning er Bevægelse i Rummet og ikke andet. Men han standser saa uvilkaarlig
ved den Kendsgerning, at Bevægelse kan opfattes (blive Fænomen), og det
forklares jo ikke ved selve Bevægelsen. (Se Tillægget til denne Paragraf). I
\textsc{Richard Avenarius'} skarpsindige og lærerige Værk »Kritik der reinen
Erfahrung« forsøges den Opfattelse gennemført, at en videnskabelig Opfattelse
af Erfaring og Tænkning kun faas ved at betragte disse Funktioner som Processer
i Centralnervesystemet. Der skildres da en hel Række saadanne Processer, men
naar der spørges, hvad de betyder, saa viser det sig, at de er byggede paa
Analogi med direkte iagttagne sjælelige Fænomener. Den saakaldte »objektive
Vitalrække«, der skal betragtes som den uafhængig variable i Forhold til den
»subjektive Vitalrække«, er i Virkeligheden afhængig af denne. (Smlgn. min
Karakteristik af \textsc{Avenarius} i »Moderne Filosofer«). Medens
\textsc{Hobbes} dog glimtvis kommer til at antyde de subjektive Forudsætninger
for den fysiologiske Psykologi, og medens \textsc{Avenarius} udtrykkelig
anerkender, at der foreligger to »Vitalrækker« i indbyrdes Funktionsforhold,
har \textsc{BECHTEREW} i sin »Objektive Psychologie« (1913) forsøgt at skrive
en Psykologi med Undgaaen af alle psykologiske Betegnelser. Men Følgen er, at
Læseren stadig maa gætte paa, hvorfor netop disse
% --- p. 115 ---
\setcounter{footnote}{0}%
\opage{115}bestemte fysiologiske Processer beskrives i denne bestemte
Sammenhæng. Saaledes lægges der stor Vægt paa »Koncentrationsreflexer«, i
Modsætning til sporadisk isolerede Reflexer, og Læseren gætter da, at det er
Opmærksomhed, der tænkes paa. Udtrykket »Koncentration« kan meget vel benyttes
i rent psykologisk Betydning, om vor Samling af Sindet i en bestemt Retning,
ligesom vi før saa, at man taler om Baning som et godt Billede for indøvet
Tankegang. Den Plads, der gives »Koncentrationsreflexerne«, skyldes en Analogi
med sjælelige Koncentrationer, saaledes som beskrivende Psykologi kan
fremstille dem. \textsc{Bechterew} indrømmer selv, at til den Forskel, den
iagttagende Psykologi gør mellem det ubevidste og det bevidste, det
uvilkaarlige og det vilkaarlige, kan han intet fysiologisk tilsvarende finde,
og derfor, siger han, har denne Forskel ingen Betydning for den objektive
Psykologi. Hvis den objektive Psykologi maa give tabt overfor Forskelle af den
angivne Art, saa har den ikke stor Betydning. Jeg tror nu, at
\textsc{Bechterew} giver for tidlig op. Selv om man ikke mener, at den
fysiologiske Psykologi er den eneste videnskabelige, kan og maa det alligevel
stilles som en uafviselig Opgave (hvad enten den kan løses eller ikke) at finde
fysiologiske Korrelater til hine i psykologisk Henseende saa betydningsfulde
Forskelle.

Fysiologiske Processer er direkte maalelige i en Grad, i hvilken psykiske
Processer ikke er det. Vi staar, naar vi stiller de to Arter af Processer
sammen, igen overfor Forholdet mellem en Kvantitetsrække og en Kvalitetsrække.
Atter her maa Identitet og Analogi ikke forvexles. Leddene i en af de to Rækker
kan bruges som Symptomer paa, hvad der foregaar i den anden. Hvis man ikke kan
forbinde nogen Mening med, at en Tanke kan være Pro%
% --- p. 116 ---
\setcounter{footnote}{0}%
\opage{116}dukt af materielle Elementers Vexelvirkning, eller med, at en Tanke
frembringer en Forandring i Legemet, saa maa man standse ved et Analogiforhold%
% footnote 1 on p. 116
\footnote[1]{løvrigt vilde ogsaa den populære Vexelvirkningsteori bygge paa
Analogi, idet den forudsætter, at Forholdet mellem Sjæl og Legeme er analogt
Forholdet mellem to Legemer (to Billardkugler f. Ex., eller --- naar det skal
være mere »videnskabeligt« --- to Atomgrupper). \textsc{Descartes} lader Legeme
og Sjæl »støde« til hinanden, og \textsc{Lotze} har ikke noget imod denne
Teori.}, der maaske engang kan blive til et Proportionsforhold. Den Opfattelse,
vi kommer til ved at holde os til, hvad der kan iagttages, er, at der i hvert
Tilfælde er nogle Væsener, som har baade psykiske og materielle Egenskaber, der
svarer til hverandre, uden at den ene Gruppe af Egenskaber kan udledes af den
anden. Det er en Gaade, der her bliver tilbage, --- et Exempel paa, at hvor vi
maa standse ved en Analogi, standser vi ved et Spørgsmaal. Et saadant Forhold
kan iøvrigt forekomme i rent psykologisk Form. Der kan i et Menneskes Karakter
være Egenskaber, der ikke kan udledes af hverandre, og som vi derfor ikke
forstaar, fordi vi ikke kender selve Karakteren dybt nok til at indse,
hvorledes de to forskellige Egenskaber i deres inderste Grund hænger sammen.
Hvad Forholdet mellem det aandelige og det materielle angaar, er vi saaledes
stillede, at vi snart kan slutte fra sjælelige Tilstande til Tilstande i
Organismen, snart omvendt. \textsc{Spinoza,} Grundlæggeren af den Opfattelse,
jeg her nærmest slutter mig til, gaar da ogsaa allerede begge Veje%
% footnote 2 on p. 116
\footnote[2]{\textit{Spinoza's Ethica.} Analyse og Karakteristik, p. 48 f.; 87 f.}.

Det er tilsidst Forholdet mellem Begreberne Personlighed og Organisme, vi her
kommer til at staa overfor. Hvad \textsc{Kant} sagde om Organismen, at der i
Grunden ingen Analogi kan findes til den, gælder ogsaa for Personligheden.
% --- p. 117 ---
\setcounter{footnote}{0}%
\opage{117}En saa inderlig Sammenhæng som mellem Sjælelivets Elementer kender
vi intet Sidestykke til. Des interessantere er det, at hine to Begreber
indbyrdes er Analoga.

\textit{Tillæg til § 19.}\textsc{Hobbes} opfattede alt, ogsaa Sansning og
Tænkning, som Bevægelse i Rummet og \textit{intet andet.} Men især paa et
enkelt Sted dukker det op for ham, at Sagen dog ikke er saa simpel. Det er
tilsidst gennem Sansning, at Fænomener, altsaa Bevægelser, bliver til for os,
og at Videnskab bliver mulig. Og naar der saa spørges om selve Sansningens
Aarsager, saa kan vi ikke tage vort Udgangspunkt fra noget andet Fænomen end
netop den (Sansningen) selv (ad causarum sensionis investigationem ab alio
phænomeno \textit{præter eam ipsam} initium sumi non posse. De Corpore XXV, 1).
Det er paa Grund heraf, at \textsc{Hobbes} erklærer, at af alle Fænomener er
selve det, at noget kan blive Fænomen (id ipsum τὸ φαίνεσθαι), det mærkeligste
(admirabilissimum).

Jeg forstaar dette mærkelige Sted paa følgende Maade. Det, at noget kan blive
Fænomen, finder \textsc{Hobbes} saa mærkeligt, fordi en Bevægelse ikke kan
blive Fænomen for en anden Bevægelse (d. v. s. \textit{kendes} af denne). Den
ene Bevægelse kan fortsætte den anden, hæmme den eller standse den. Sansning
maa da være forskellig fra \textit{blot} Bevægelse, skønt den \textit{ogsaa} er
Bevægelse. Naar det nu er gennem Sansning, vi kender alt, saa maa det ogsaa
være gennem Sansningen selv, vi kender Sansningen. Men at sanse sin egen
Sansning er at erindre, og der er da ogsaa, hævder \textsc{Hobbes,} til al
Sansning knyttet Erindring og Sammenligning. (Sentire se sensisse, memoria est.
--- Sensioni necessario adhæret memoria aliqva, qva priora cum posterioribus
comparari et alterum ab altero distingui potest. De Corpore XXV, 1; 5). Det er
i Konsekvens af denne
% --- p. 118 ---
\setcounter{footnote}{0}%
\opage{118}sin Opfattelse af Sansningen, at han opstiller sin berømte Sætning,
at der i al Sansning indeholdes en Mangfoldighed (phantasmatum varietas), og at
det stedse at sanse det samme og slet ikke sanse kommer ud paa et (sentire
semper idem et non sentire ad idem recidunt) (5). Ud fra denne rent
psykologiske Iagttagelse slutter han saa, at den Bevægelse, til hvilken
Sansningen er knyttet, maa kunne fastholdes og genkaldes, og at der i
Sanseorganet maa være Betingelser derfor. Tydeligere kan Forskellen mellem
Erkendelsessubjektet og den Bevægelse, der ifølge \textsc{Hobbes} er det eneste
objektivt existerende, ikke udtales. Men \textsc{Hobbes} var for meget
Dogmatiker til, at dette Tankeglimt kunde faa ham til at opgive sin
Bevægelsesmetafysik. Det er gaaet ham, som det senere gik \textsc{Schelling,}
der (som jeg har vist i »Den nyere Filosofis Historie«\textsuperscript{3} III,
p. 158) havde et analogt Tankeglimt, idet han meget klart opstillede det
Problem: hvorledes kan der ikke blot existere Natur, men ogsaa en
Naturvidenskab? Dette Glimt hindrede ligesaalidt \textsc{Sghelling} %
% PRINTED AS IS: „Sghelling“ (p. 118)
i at gaa videre i sin romantiske Naturfilosofi, som \textsc{Hobbes'} Glimt
hindrede denne i at gaa videre i sin mekanistiske Naturfilosofi.

I \textsc{Frithiof Brandt's} udmærkede Monografi »Den mekaniske Naturopfattelse
hos Thomas Hobbes« (1921), i hvilken der kastes Lys over saa mange hidtil
oversete eller uforstaaede Punkter i \textsc{Hobbes'} Filosofi, drøftes det
ovenfor anførte Sted udførlig (p. 365---373). Forfatteren søger at vise, at der
i Virkeligheden ingen Inkonsekvens her er hos \textsc{Hobbes.} Men der tages
ikke Hensyn til den i mine Øjne afgørende Sætning, at Forstaaelsen af
Sansningen maa vindes ud fra den selv, og til den deraf følgende Omstændighed,
at man ud fra den rent psykologiske Forstaaelse af
% --- p. 119 ---
\setcounter{footnote}{0}%
\opage{119}Sansningens Natur maa \textit{slutte} til Beskaffenheden af
Sanseorganet og af, hvad der maa foregaa i det.

Min Tydning af den omtalte Ytring af \textsc{Hobbes} som et Tankeglimt, der
fører ud over hans absolute Bevægelsesmetafysik, bekræftes ved hans Udtalelse
(De Corpore VI, 7) om, at man kan komme til Etik og Statslære uden at gaa den
naturvidenskabelige Vej, men blot ved at bygge paa de psykologiske Erfaringer
om Trang og Drift, enhver direkte kan gøre hos sig selv (per uniuscujusqve
proprium animum examinantis experientiam). Ikke blot Trang og Drift, men ogsaa
Sansefornemmelse kender vi direkte ad denne Vej.

\begin{center}20. Etik og teoretisk Videnskab.\end{center}

Naar Etik stilles i en vis Modsætning til teoretisk Videnskab, beror det ikke
paa, at man indenfor Etiken skulde have Ret til at tænke og slutte paa anden
Maade end indenfor al anden Videnskab. Der gives ingen særlig etisk Tænken,
hvad Begrundelse og Bevisførelse angaar. Men Etiken har Forudsætninger, som
andre Videnskaber ikke behøver at opstille. Enhver nok saa teoretisk Tænken
hviler paa den Forudsætning, at Sandhed har Værdi, at det er værdifuldt at søge
efter den. I den Grad er denne Forudsætning nødvendig, at man end ikke anser
det for nødvendigt at nævne den. Naar i Nutiden en filosofisk Retning, hvis
betydeligste Repræsentant er \textsc{Heinrich Rickert,} alligevel har lagt Vægt
paa denne Forudsætning%
% footnote 1 on p. 119
\footnote[1]{\textit{Den nyere Filosofis Historie}\textit{8} IV, p. 219.},
ligger Berettigelsen i, at Værdiproblemet derved drages ud af den Isolation,
der let kan blive Følgen, naar det kun ses i sin Betydning for Æstetik, Etik
(Økonomi indbefattet) og Religionsfilosofi. Naar saadan Isolation afværges,
fremtræder der en fundamental Sammenhæng i det menneskelige Aands%
% --- p. 120 ---
\setcounter{footnote}{0}%
\opage{120}liv, en Sammenhæng, der ligger dybere end Forskellen mellem de
forskellige Problemer, men tillige begrunder den Forskel, der er. Naar
Videnskab foretrækkes for kunstnerisk eller religiøs Stræben, saa er det Værdi,
der staar imod Værdi, og der finder en Brydning Sted mellem forskellige
aandelige Fornødenheder. Problemerne bliver til derved, at enhver af disse
aandelige Fornødenheder sætter Tanker i Bevægelse. Det kommer da an paa, om
disse Tanker kan bringes i Overensstemmelse med hverandre. Men her er det, at
den store Betydning af Sandhedsværdien, al Videnskabs stiltiende Forudsætning,
fremtræder, netop i Forhold til alle andre Værdier. De Tanker, de forskellige
Værdier sætter i Bevægelse, maa --- for Værdiernes egen Skyld --- underkastes
en rent teoretisk Prøvelse, en Prøvelse, hvis eneste Formaal er at finde
Tankernes Begrundelse og deres indbyrdes Forhold, en Prøvelse altsaa ud fra
Sandhedsværdien som raadende. Af alle Problemerne faar derfor
Erkendelsesproblemet Forrangen. Af det Resultat, Drøftelsen af
Erkendelsesproblemet fører til, afhænger den Gyldighed, der kan tillægges
Tanker, som sættes i Bevægelse af andre Motiver end de rent intellektuelle. Ved
Behandlingen af Erkendelsesproblemet forudsættes, at vi ikke har anden
Interesse end den at finde Sandheden, hvorledes denne end skulde være
beskaffen. Derefter bliver det da Spørgsmaalet, hvorledes andre Interesser, der
rører sig hos os, kan forenes med denne Interesse. I min Bog »Den menneskelige
Tanke« (se især p. 22) er jeg derfor gaaet saaledes frem, at jeg forudsætter,
at der har udviklet sig en intellektuel Interesse, en Glæde ved selve
Tankearbejdet, der fører til, at Tanken kan udvikle sig efter sine egne Love og
efter Emnernes Krav. De andre Interesser, for hvem Tankearbejdet blot er
Middel, ikke Maal i og for sig, vil tilsidst
% --- p. 121 ---
\setcounter{footnote}{0}%
\opage{121}finde sig bedst tjente med, at den intellektuelle Interesse for en
Tid raader, og med at der fra Tid til anden foretages en teoretisk Revision af
de Tanker, de særlig har kaldt til Live.

Paa det etiske Omraade kan det nu ikke undgaas, at andre Interesser end de rent
intellektuelle gør sig gældende. Al etisk Tænken forudsætter en Grundværdi og
gaar ud paa at finde Midler og Veje til at hævde den i de Domme, der fældes, og
de Samfundsordninger, der dannes. Paa en træffende Maade har \textsc{Henri
Poincaré} i sin Afhandling »La Morale et la Science«%
% footnote 1 on p. 121
\footnote[1]{\textit{Dernières Pensées.} 1913, p. 229.} udtrykt dette derved,
at etisk Tænken har en Præmis, der er imperativisk, idet den udspringer af en
Trang, som kræver Tilfredsstillelse, medens de andre Præmisser er
indikativiske, idet de angiver de Midler og Veje, ved hvilke Tilfredsstillelsen
under de givne Forhold kan naas. Da der er en imperativisk Præmis, maa
Konklusionen nødvendigvis ogsaa blive imperativisk. Den Konsekvens, der her kan
være Tale om, bliver en Konsekvens af Forholdet mellem Midler og Formaal.
Spørges der om Begrundelsen af den Trang eller Værdi, der afføder Imperativet,
maa Svaret, da Værdi kun kan maales ved Værdi, hentes fra den Grundværdi, den
specielle Værdi forudsætter. Og spørges der igen om Begrundelse af denne
Grundværdi, maa vi holde os til den Kendsgerning, at enhver Værdibestemmelse
forudsætter en Helhed. Værdier svæver ikke i Luften, men er Udtryk for en
Helheds Trang til Bestaaen. Det kan være et enkelt Øjeblik, en Livsperiode, en
enkelt Interesse, man gør til et og alt; det kan være den egne Personlighed
eller den Stand og Stamme, den hører til, eller tilsidst hele Menneskeslægtens
Bestaaen og Udvikling, der udgør den sidste Forudsætning for den
% --- p. 122 ---
\setcounter{footnote}{0}%
\opage{122}perativiske Etik. Teoriernes Strid beror tilsidst paa, at der kan
lægges forskellige Helheder til Grund, og at derved ogsaa Grundværdierne bliver
forskellige%
% footnote 1 on p. 122
\footnote[1]{Sammenhængen mellem Begreberne Totalitet og Værdi har jeg især
paavist i \textit{Den menneskelige Tanke} 145---146; 151, og i
\textit{Totalitet som Kategori,} Kap. V.}. Et Forsøg paa videnskabelig Etik kan
her kun gaa saaledes frem, at man tager Udgangspunktet fra en bestemt Helhed,
drager alle Slutninger fra Betingelserne for dens Bestaaen og Udvikling og
derefter søger at vise, hvorledes de andre Helheder er at vurdere i Henhold til
de saaledes vundne Resultater. Ved denne sidste Undersøgelse bliver der særlig
Grund til Analogi. Selv om de andre Helheder i Kraft af Udgangspunktet maa
behandles som blotte Midler, bliver saadan Behandling dog kun formaalstjenlig,
naar der er Forstaaelse af deres Natur. Og hvis de er mere end Midler, --- hvis
et Solidaritets- eller endog et Sympatiforhold er Følgen af det givne
Udgangspunkt, kræves der endnu dybere gaaende Forstaaelse af dem; heldigvis kan
i saadanne Tilfælde Evnen til Forstaaelse vokse med Trangen. Mellem
ejendommelige Helheder indbyrdes er der et Analogiforhold, ikke et
Identitetsforhold. Etisk Tænkens Gyldighed (ud fra det givne Udgangspunkt)
beror paa den Strenghed, hvormed Analogierne gennemføres.

Man kunde mene, at den mest omfattende Helhed i og for sig maatte afgive det
rette Udgangspunkt for etisk Vurdering, saa at dens Bestaaen blev den
selvskrevne Grundværdi. Men foruden Omfanget kommer ved Værdibegrebet ogsaa
Inderligheden og Styrken i Betragtning. Indenfor en mindre omfattende Helhed
kan der være en Inderlighed og en Styrke i Sammenslutningen, som ikke saa let op%
% --- p. 123 ---
\setcounter{footnote}{0}%
\opage{123}staar indenfor en større Helhed, og den større Helheds Værdi vil for
en væsentlig Del bero paa den Selvstændighed, den formaar at anerkende og at
fremhjælpe hos de mindre Helheder, den omfatter.

Ved Analogierne mellem de forskellige Helheder bliver det forstaaeligt, at der
er visse etiske Begreber, der dannes, og maa dannes paa alle Standpunkter,
fordi de udtrykker det formelle Forhold mellem en Helhed, betragtet som
Grundværdi, og Betingelserne for dens Bestaaen og Udvikling. Saadanne formale
etiske Begreber er Ret og Pligt. Den enkelte Helhed skal hævde sig overfor
andre Helheder og overfor ydre Forhold idethele; det er dens Ret, der følger af
dens Stilling som Grundværdi. Men dens Opgave er desuden at værne om Midlerne
til sin Bestaaen og ogsaa at virke for de Helheders Bestaaen, der er Midler for
den, eller til hvilke den staar i Solidaritets- eller Sympatiforhold. Under
Pligtbegrebet hører Klogskaben, Forstaaelsen af Midlernes Natur og
Nødvendigheden af at fremskaffe dem. Især med Hensyn til disse formale
Begrebers Betydning og Anvendelse vil det ene etiske Standpunkt kunne lære af
det andet. Det indskærpes, hvad Klogskabspligten angaar, af Parablen om den
utro Forvalter (Lucas XVI), som viser, at Mørkets Sønner ofte indenfor deres
Sfære raader over en Klogskab, der savnes hos Lysets Sønner i deres Arbejde.
Det er i vore Dage udtalt af en erfaren Statsmand, at de kloge Menneskers
Opgave ofte bestaar i at raade Bod paa de Ulykker, gode Mennesker har voldet.
--- \textsc{Leibniz} har defineret Begrebet Retfærdighed saaledes, at Ret og
Pligt, Menneskelighed og Klogskab er forenede i det: caritas sapientis!

Ikke blot mellem forskellige etiske Standpunkter indbyrdes kan Analogi gøre sig
gældende paa frugtbar Maade, men ogsaa mellem etisk Bevidsthed og videnskabelig Be%
% --- p. 124 ---
\setcounter{footnote}{0}%
\opage{124}vidsthed. Den videnskabelige Bevidsthed, der søger Klarhed og
Forstaaelse for deres egen Skyld, kan blive et Forbillede (et Paradigma, for at
bruge Aristoteles' Udtryk) for etisk Bevidsthed, baade hvad Sindelag og hvad
Konsekvens angaar. Det samme kan siges om den kunstneriske Bevidsthed i dens
Stræben efter at yde det fuldkomne. Intet, siger \textsc{Michel Angelo,} gør
Sjælen saa ren og ædel som det at stræbe efter at øve et fuldkomment Værk. Paa
alle aandelige Omraader gælder det at forene Enhed og Sluttethed med fri
Udfoldelse af en Kraftfylde. Retfærdighed (i den anførte Betydning af Ordet)
kan kaldes den etiske Syntese, ligesom Sandhed er den intellektuelle Syntese.
(Den menneskelige Tanke, p. 365---367). I den ovenanførte Afhandling om Moral
og Videnskab siger \textsc{Henri Poincaré:} »De intellektuelle Vaner har en
moralsk Genklang. Naar man vænner sig til at se bort fra Enkeltheder og
Tilfældigheder, fordi de ikke er af Betydning for Forstaaelsen, vil der opstaa
en naturlig Tilbøjelighed til overhovedet at underordne de særlige Interesser
under de almene«.

\begin{center}21. Verdensanskuelse og Videnskab.\end{center}

Det er et sandt Ord af \textsc{Wundt}, at det er lettere at forkaste al
Metafysik end at handle efter denne Forkastelse. Ad mange skjulte Veje sniger
der sig Kombinationer og Konstruktioner ind i Baggrunden af Bevidstheden, og af
og til trænger de sig endog frem i Forgrunden. Atter og atter vil en kritisk
Undersøgelse af saadanne Konstruktioner være nødvendig. Og først og fremmest er
det af Betydning, at man bliver klar over Beskaffenheden af det Problem, man
her staar overfor. Som jeg et andet Sted (Den menneskelige Tanke, p. 306 f.)
har forsøgt at vise, er der en Analogi mellem Tilværelsesproblemet og
Erkendelsesproblemet. Erkendelsen bliver, hvad dens Gyldighed
% --- p. 125 ---
\setcounter{footnote}{0}%
\opage{125}angaar, et Problem derved, at menneskelig Tanke vel er en
Kendsgerning, et faktisk foreliggende Emne, men at den ikke er det eneste Emne
og dog skal rumme baade sig selv og alle andre Emner. Den skal udvikle sig
efter sine egne Love og dog bygge en Verden, i hvilken alle Emner, den selv
med, kan finde en til deres Ejendommelighed svarende Plads. Tanken skal raade
og er dog Element i en Verden, der er større end den selv. Verden er stor og
Tanken er lille, saa energisk den ogsaa rører sig. Tilværelsesproblemet, d. v.
s. Spørgsmaalet om en videnskabelig Verdensanskuelses Mulighed, skyldes
ligeledes et Forhold mellem Del og Helhed. Men her gælder Spørgsmaalet, om
noget enkelt Emne, nogen Del af Tilværelsen kan staa som Type for hele
Tilværelsen, som det Urfænomen, der giver hele Tilværelsens Karakteristik i
kort Begreb. Alle historisk foreliggende Verdensanskuelser, religiøse eller
spekulative, er uvilkaarlige eller med bevidst Eftertanke konstruerede
Krystallisationer om et Urfænomen. 1 Rækken af disse Krystallisationer kan man
aflæse, hvilke Emner der til forskellige Tider har draget Interesse til sig og
er blevne ubevidste eller bevidste Midtpunkter i Menneskers Sind. --- Det er
tydeligt, at der ved Tilblivelsen af en Verdensanskuelse tænkes efter Analogi,
ἐν παραδείγματι som Aristoteles siger. Et Emne bruges som Forbillede for alle Emner.

Vi har her uden videre opereret med Begreberne Del og Helhed. Men har vi Ret
til at gaa ud fra, at Tilværelsen udgør en Helhed? --- \textsc{Bertrand
Russell} har udtalt, at Forestillingen om et Univers som Helhed er en Rest af
prækopernikansk Astronomi, som end ikke \textsc{Kant's} Kritik har kunnet
skaffe bort%
% footnote 1 on p. 125
\footnote[1]{\textit{Mysticism and Logic,} p. 987.}. Det er vel ogsaa saa, at
det er det sanselige Verdensbillede, der bestemmer Forestillingerne
% --- p. 126 ---
\setcounter{footnote}{0}%
\opage{126}paa dette Omraade, selv om Videnskaben forlængst har opgivet dets
absolute Gyldighed. Men der er dog et dybere Grundlag for den Tendens, der her
gør sig gældende. Grundloven for alle Kategorier, saavel som for Bevidsthedsliv
overhovedet, er Syntesens Lov. Der arbejder her, for at bruge \textsc{Kant's}
Udtryk, en skjult Kunst i Menneskets Indre, som væver Sansefornemmelser til et
Sansebillede, Sansebilleder til Begreber og Love, og tilsidst ved den
lovmæssige Sammenhængs Hjælp et stort Helhedsbillede, mere eller mindre
gennemført og udformet. Kritiken kommer stedse bagefter. Og den vigtigste
kritiske Indsigelse mod et Forsøg paa at slaa saadanne Helhedsbilleder fast er
den, at det vilde være Tankens Død. Naar al Tænken tilsidst er Sammenfatten,
forudsættes der ved alt Tankearbejde en Mangfoldighed, en Diskontinuitet, som
skal overvindes. Og ved en fuldstændig og afsluttet Tilfredsstillelse af
Trangen til Helhed, vilde selve hin »skjulte Kunst« dø Straadød. Den vilde ikke
mere, hvad den i sine fineste og ædleste Former har kunnet, udforme store
Idealer, stille Maal, der peger ud over det hidtil vundne. Diskontinuitet maa
have en Plads blandt de fundamentale Kategorier for at minde om, at der stadig
er Arbejde at gøre. Enhver Helhed, som findes, kan kun betegne en foreløbig
Station. Vor Erkendelse arbejder stadig paa en Skala, hvis øverste Trin vilde
være den uopnaaelige absolute Identitetsrække, og hvis nederste Trin vilde være
den upaaviselige kaotiske Forskelsrække.

I de mest udprægede kosmologiske Systemer --- \textsc{Parmenides', Spinoza's,
Hegel's} --- har Identitet været det sidste og højeste Synspunkt. Idealet var
at opfatte alt som i sidste Instans ét. Det var igrunden et orientalsk Ideal,
der her kom til at raade i europæisk Filosoferen. Et Vende%
% --- p. 127 ---
\setcounter{footnote}{0}%
\opage{127}punkt kom med \textsc{Leibniz,} for hvem Identitet betegner et
metodisk Princip i vor begrundende og bevisende Tankegang og kun ad Analogiens
Vej kan faa kosmologisk Betydning. \textsc{Leibniz} hævder netop med stort
Eftertryk, at al Kosmologi (Metafysik), og dermed ogsaa al Teologi, bygger paa
Analogi, tilsidst Analogi med den menneskelige Aand og den Harmoni og Helhed,
som kan findes i den, eller som den i hvert Tilfælde stræber efter. Endnu mere
indgaaende søgte \textsc{Kant} at drage Grænsen mellem den med
Identitetsprincipet arbejdende Videnskab og den med Analogier arbejdende Tro. I
det 19. Aarhundredes Begyndelse gjordes nye Forsøg paa kosmologisk
Identitetslære. Men det synes, som om Analogisynspunktet vinder mere og mere
Tilslutning, selv blandt Tænkere af overvejende spekulativ Retning%
% footnote 1 on p. 127
\footnote[1]{Smlgn. f. Ex. \textsc{Bosanquet:} \textit{The meeting of extremes
in contemporary philosophy} (1921), p. 17.}. Den billedlige Karakter af de
Forestillinger, som uvilkaarlig eller med fuld Bevidsthed dannes paa Tankens
Grænse og ud over den, trænger mere og mere igennem. Denne billedlige Karakter
tyder paa, at man har gjort et Valg mellem de indbyrdes analoge Rækker, den
formale (logisk-kvantitative) og den reale (kvalitativt-successive) Række, og
indenfor den reale Række den naturvidenskabelige og den aandsvidenskabelige.
Men selve Analogien mellem disse Rækker er og bliver Videnskabens sidste
Resultat. En videnskabelig Afslutning vilde kræve Rækkernes Identitet, ikke den
blotte Analogi. Det er nu engang saaledes, at vort exakte Virkelighedsbegreb
bliver til ved, at de formale Rækker gør de reale Rækker forstaaelige, og at
selve de formale Rækker er en Side af Virkeligheden, den Side, vi bliver
opmærksomme paa, saasnart vi prøver paa at orientere os i Tilværelsen.

% --- p. 128 ---
\setcounter{footnote}{0}%
\opage{128}Hvorledes man stiller sig paa selve Grænsen for videnskabelig
Tænken, vil tildels bero paa, hvilken Erkendelsesteori man hylder. Der bliver
her en noget forskellig Situation for Pragmatismen og for Kriticismen. Begge
Opfattelser hævder, at de Tanker, der kan formes, naar vi har naaet hin Grænse,
kun kan være Analogier. Men Pragmatismen, som idethele har Tilbøjelighed til at
tillægge alle Tanker Gyldighed, der udspringer af menneskeligt Behov, og som
ikke interesserer sig for kritisk Undersøgelse af vore Grundbegreber, vil
heller ikke have nogen Interesse af at undersøge, hvorvidt man endnu i de
Grænsebegreber, der er den menneskelige Erkendelses sidste Ord, kan spore de
Grundformer, som er karakteristiske for den. Kriticismen derimod vil
interessere sig for at genfinde den menneskelige Aands Struktur ud over det
Omraade, hvor der kan dannes videnskabelige Begreber; i selve Arten af de
billedlige Forestillinger, der endnu er mulige ud over dette Omraades Grænser,
vil den søge efter Ejendommeligheder for menneskeligt Aandsliv. Den lægger Vægt
paa Tankens Homologier, paa den Beslægtethed, fælles Struktur kan betinge, ikke
blot paa de Analogier, der opstaar ved, at Tanken skal afpasses efter nye Forhold.

Kosmologi kan ikke blive Videnskab. Den staar mellem Videnskab og Kunst. Den
danner Overgangen til Poesi og Religion, de Omraader, hvor Trangen til at finde
konkrete Udtryk for Livserfaringer frit og uvilkaarlig viser sig. Selv der vil
Filosofien ikke blot søge Analogi, men ogsaa Homologi, idet den menneskelige
Aands Natur ogsaa der vil lægge sig for Dagen.
% --- p. 129 ---
\setcounter{footnote}{0}%
\opage{129}

\begin{center}\large V. Poetisk og religiøs Symbolik.\end{center}

22. Af den Opgaaen i uvilkaarlige Analogier, som er ejendommelig for
menneskeligt Aandsliv paa det mest primitive Stadium, vi kender, har gennem
mange Overgangsformer udviklet sig den videnskabelige Bevidsthed, der kræver
Identitet som den strenge Forudsætning ved enhver Tankeovergang, og som, hvor
denne Forudsætning ikke kan paavises, arbejder med rationelle Analogier.
Videnskabens Udvikling er et Led i den store aandelige Arbejdsdeling, som
kommer ved fremskridende Kultur. Men naar ét Element udløses af den primitive
Sammenhæng, i hvilken alle aandelige Kræfter usondrede er forbundne, udløses
ogsaa andre Elementer eller faar i hvert Tilfælde andre Vilkaar at virke under.
Ved Siden af, og mere eller mindre under Indflydelse af den intellektuelle
Kulturs Selvstændighed, har ogsaa Poesi og Religion i større eller mindre Grad
udløst sig af den oprindelige Bundethed. I Modsætning til de rationelle
Analogier træder da efterhaanden de for Poesi og Religion karakteristiske
emotionelle Analogier.

Den umiddelbare Optagethed, eller rettere Betagethed af et Emne kan afføde en
ligesaa umiddelbar Trang til at forme det i et Billede eller i Ord og derved
kundgøre dets Værdi for sig selv og for andre. Vi lever eller føler os ind i
Emnet, stræber at følge det i dets Skæbne og at faa Del i dets Egenskaber. Det
Individuelle og umiddelbart Anskuelige, som ingen Videnskab kan udtømme i sine
Begreber eller i sine Love, giver Poesien os i sine af Oplevelse og Indlevelse
udsprungne Billeder, og maaske endog med en saa fuldendt Individualisering, som
de i Erfaringen foreliggende Emner ikke havde. Og hvad Tanken kan udtrykke i
blege Former, Modsætningen mellem det lave og det høje,
% --- p. 130 ---
\setcounter{footnote}{0}%
\opage{130}det bitre og det liflige, det flygtige og det bestaaende, fører
Poesien, især som Drama, os umiddelbart ind i. Om det, der gives os paa denne
Maade, er Regel eller Undtagelse, kan Poesien ikke selv afgøre. Ethvert Exempel
er et Symbol, men ikke ethvert Symbol er et Exempel. (smlgn. ovenfor § 10).
Poesien kan kun forsaavidt stille Problemer, som allerede det, at et enkelt,
enestaaende, individuelt Emne hører med til Tilværelsen, kan tjene til
Karakteristik af denne, eller dog give Anledning til alvorlig Eftertanke. Men
ofte vil den ved Emnet vakte Sindsbevægelse være saa stærk, at intet enkelt
Symbol tilfredsstiller, men en hel Række af Symboler udvikles. Det er en
Rækkedannelse af anden Art end den intellektuelle. (Se § 5). Ikke det indre,
logiske Forhold mellem Leddene i Rækker, men den alle Leddene underliggende
Sindsbevægelse betinger her Kontinuiteten og knytter Leddene sammen i en Række.

Det betegner, som allerede \textsc{Schiller} har set, et Kulturfremskridt, naar
der kan være Glæde ved Billedet, bortset fra, om Emnet er i vort Eje eller
ikke, og bortset fra, om det har nogen Betydning som Tegn eller som Middel.
Saalænge »Participationen« (»Helhedsrealismen«) raader, kan denne Glæde ikke
opstaa. Men selv om Billedet er frigjort for praktisk Behov af en eller anden
Art, bliver Glæden ved det dog ikke af rent intellektuel Art. Det er en
Forskydning, der har fundet Sted, og med Forskydningen følger en Forædling;
Interessen er flyttet fra Emnets Realitet til Billedets Idealitet og er derved
selv undergaaet en Metamorfose.

Der maa her skelnes mellem det Sind, der føder Billedet, og det Sind, der
paavirkes af Billedet, naar det er kommet til Verden.

Hos Digteren er Følelsen det primære. De digteriske
% --- p. 131 ---
\setcounter{footnote}{0}%
\opage{131}Billeder eller Lignelser fremgaar af en Følelse af Emnets Værdi, som
vækker en Trang til at dvæle ved det og se det fra alle Sider. Maaske faar
Sindet først Luft ved en hel Række af Lignelser. Disse frembyder da Analogier
indbyrdes, men disse Analogier skyldes den Inderlighed, hvormed Digterens Sind
fastholder Emnet og dets Værdi. Saaledes er der Analogi mellem Jesu Parabler
indbyrdes, fordi de bæres af en og samme Grundstemning, den, der er vakt ved
Ideen om Himmeriges Rige.

Læseren eller Tilhøreren er her anderledes stillet end Digteren selv. Billedet
eller Lignelsen viser ham maaske med slaaende Klarhed og i et Helhedsudtryk det
Forhold, hvorom det gælder. Og i Forhold til denne intellektuelle Virkning
bliver Følelsen afledet. Digteren behøver ikke at have sat sig det udtrykkelige
Formaal at frembringe denne Virkning. I Begejstringens Øjeblik gør han det i
hvert Tilfælde ikke. Han synger, som Fuglen synger.

I sin interessante Afhandling »Metafor och Fiktion. Ett bidrag til Poetikens
Teori« (Lund 1922) støtter Dr. \textsc{Nyman} sig til Udsagn af
Forsøgspersoner, hvem forskellige digteriske Billeder blev forelagte. Her
følger det af sig selv, at Følelsen bliver sekundær, og at Billedets Virkning
til at fremkalde et intellektuelt Fugleperspektiv kommer til at staa i første
Række. Men til fuld æstetisk Værdsættelse af Billeders og Lignelsers Betydning
maa der baade spørges om Udspringet og om Virkningen. Og Digteren kan man ikke
gøre til Forsøgsperson, saa at man ad den Vej kunde komme efter, hvorledes
Billedet bliver til. Han maa belures. Og \textsc{Julius Lange} har vist sig at
have luret godt, naar han i sin Definition af Kunstværdi som første Led satte
»den Værdi, et Emne har for Kunstneren«. Til Gengæld lagde han for liden Vægt
paa de Veje og Midler,
% --- p. 132 ---
\setcounter{footnote}{0}%
\opage{132}gennem hvilke den Værdi kunde fremkaldes for Tilhører eller Læser
(Smlgn. ovenfor § 4 b).

Al Poesi staar i Forbindelse med Livet og dets Skæbner, og det kan derfor blive
vanskeligt at skelne mellem Poesi og Religion, at afgøre, hvor hin ender og
denne begynder, eller omvendt. Ogsaa Religionen taler i Billeder og Symboler.
Den Forskel, der kan gøres, beror væsentlig paa, at Udformningen af Billederne
som Billeder spiller en større Rolle for Poesien end for Religionen. Sagen er
vel, at de til Grund liggende Oplevelser her er dybere og stærkere betagende og
mere henvender sig til Villien end til Fantasien med de Fordringer, der
udspringer af dem. Trangen til umiddelbar Leven i Symbolerne og til derigennem
at faa Del i de Oplevelser, af hvilke Symbolerne udsprang, da de dannedes paa
første Haand, er karakteristisk religiøs, og derpaa beror den store Rolle,
Overlevering ofte spiller i Religionen. Det er ikke blot, fordi Overleveringen
i og for sig stadig øver sin Magt over Menneskenes Sind, at Dogmer og
Kultusformer holder sig saa længe. Det er ogsaa, fordi Individet dér staar
overfor Symboler, ved hvis Tilblivelse han véd, at han ikke har været
medvirkende, saa at Billede og Sag kommer til at staa som ét, og den
umiddelbare Indlevelse lettere bliver mulig. Ligesaalidt som Individet er sig
bevidst, at han har medvirket ved Symbolets Dannelse, ligesaalidt skelner han
mellem Tydning og Oplevelse, mellem Billede og Virkelighed. Digteren klager
over, at »Livet svæver mellem Syn og Eje«; men denne Forskel er ikke til paa
udpræget religiøse Stader. Der vil her i hvert Tilfælde være et Punkt, hvor Syn
og Eje falder sammen. Herpaa beror det, at Religion længere end Videnskab og
Poesi bærer Præget af det primitive Participationsstandpunkt og dets
uvilkaarlige Analogier.

% --- p. 133 ---
\setcounter{footnote}{0}%
\opage{133}Skilsmissen mellem Oplevelse og Tydning, mellem Eje og Syn gaar
næppe for sig uden Smerte. Det religiøse Problem, som det stiller sig i vore
Dage, angaar væsentlig det Spørgsmaal, om denne Skilsmisse kan gennemføres uden
Skade for Aandslivets Energi og Koncentration.

I Oplevelsen forholder vi os nemlig modtagende; Livet strømmer ind over os med
sine vældige Bølger. I Tydningen er vi mere selvvirksomme; den er inddæmmende
og begrænsende, et Selvforsvar overfor de stærke Paavirkninger. De gamle
Frisere sagde, at Gud skabte Havet, men Friserne Digerne. (Deus mare, Frisones
littora fecerunt). Inddæmningen kan være nødvendig for at bevare den Jordbunds
Frugtbarhed, hvorpaa Livet skal føres. Men her griber individuelle Kræfter ind,
og der kan endog komme en Spænding mellem Oplevelse og Tydning. Saalænge
Overlevering paa én Gang giver os begge Dele, eller rettere, saalænge det ikke
ses, at Overleveringen selv indeholder en Tydning, saalænge opstaar der intet
egentlig religiøst Problem. Der kan idethøjeste opstaa Spørgsmaal om, hvorledes
det overleverede skal tilegnes og forstaas.

Men det egentlige Problem har Religionen fælles med Kosmologien. I ethvert
Symbol tjener en Del til at belyse Helheden. Ogsaa i de religiøse Symboler skal
en Del eller Side af Tilværelsen belyse, ja aabenbare hele Tilværelsens
inderste Væsen. Mennesker klamrer sig, i alle Former for Religion eller
Religiøsitet, til et enkelt Emne i Tillid til, at det betegner et Gennembrud af
det højeste Værdifulde. Valget staar mellem denne bestemte Tydning og
Muligheden af et Kaos af Værdiløshed og Meningsløshed. Hos \textsc{Pascal} kan
man særlig studere dette Enten---Eller, der danner en Analogi til den logiske
Modsætning mellem kaotisk Forskelsrække og partiel eller absolut Identitetsrække.
% --- p. 134 ---
\setcounter{footnote}{0}%
\opage{134}Der var for \textsc{Pascal} kun ét eneste Lysglimt i Tilværelsens
Mørke. Ogsaa den dybsindige, af de store Gaader dybt grebne \textsc{Jacob
Böhme} standsede for en Stund ved dette Enten---Eller. Men han greb den Udvej
at se Modsætningerne under Øjne og at opfatte Tilværelsen som en stor
Kampplads, hvor Striden stod mellem Lys og Mørke, Renhed og Brynde, Kærlighed
og Had, og han fandt her Grundlaget for en djærv etisk Livsførelse, medens
\textsc{Pascal's} Dilemma førte ham til streng og ensom Askese som eneste Udvej.

Naar en Del skal staa som Udtryk for Helheden, vil der uvilkaarlig gøre sig en
Stræben gældende for at bevise, at vedkommende Del virkelig er typisk, maaske
endog, at alle andre Emner, Erfaringen frembyder, kan afledes af den. Derved
bliver Religion til Teologi og Spekulation, og Sofisterierne ligger nær. Som vi
har set, er et Exempel altid Symbol, men et Symbol ikke altid et Exempel. Det,
der bliver Symbol, kan være et Unicum. Den religiøse Digter siger om sin Gud,
at han »bjerger den bævende Due«. Men er der ikke mange stakkels Duer, som ikke
bjerges? Og hvis de alle bjergedes, hvorledes vilde det saa gaa Duehøgene, der
jo dog ogsaa er Emner? Det mangler jo heller ikke paa Forsøg paa at tage Høgene
til Symboler. --- Den religiøse Symbolik vil paa sit Højdepunkt betragte det
Emne, den oprindelig gik ud fra, som afledet i Forhold til den Generalisation
eller Absolutificering af Emnet, den uvilkaarlig fører til. Fadersymbolet er
oprindelig hentet fra den Faderkærlighed, der kan erfares i menneskelige
Forhold. Men saa vendes Forholdet om, og det, at der i menneskelige Forhold kan
øves Faderkærlighed, udledes af, at Gud er Fader (Eph. III, 15). (Smlgn.
ovenfor p.32). --- Eller det er en enkelt historisk Skikkelse, som ved Sindets Renhed
% --- p. 135 ---
\setcounter{footnote}{0}%
\opage{135}og ved sin Kærlighed til Mennesker gør det Ønske levende, i den at
se en Aabenbaring af Verdens inderste Grund. Men ikke blot staar der overfor en
saadan Skikkelse Typer af helt anden Art. Der opstaar ogsaa det Spørgsmaal, om
overhovedet nogen enkelt Skikkelse formaar i sig at forene alle de værdifulde
Egenskaber, der gennem de skiftende Tider og Kulturer maa forudsættes hos den,
der skulde være det stadige Forbillede for stræbende Mennesker?

\textsc{Stanley Hall} har i en mærkelig og indholdsrig Bog (Jesus Christ in the
light of Psychology. 1921) gjort et Forsøg paa at vise, at Jesusskikkelsen,
naar den opfattes med historisk Kritik og med grundig psykologisk Sans, er og
bliver det store Forbillede for Menneskelivet, under alle Livets Skæbner og i
al Livets Stræben. \textsc{Stanley Hall} mener ved Psykologiens Hjælp at være
kommen ud over Modsætningen mellem Ortodoxi og Kritik, og navnlig at have
overvundet den Vanskelighed, som fremkommer, naar Forbilledet skal vejlede
overfor Interesser, Opgaver og Erfaringer, der ikke forelaa, og ikke kunde
foreligge for den historiske Skikkelse, der er gjort til Forbillede. Den »nye
Tids« Jesus, som \textsc{Stanley Hall} vil fremstille, og til hvis Fremstilling
han benytter alle den nyere Psykologis Hjælpemidler, skal i sig forbinde alle
store Idealer, ligegyldigt hvorfra de saa stammer. Han er klar over, at der
mellem disse, fra forskellige Egne af Livet og Historien fremgaaede Idealer vil
være saa store Modsætninger, at en fuldstændig Harmoni vil være umulig; men han
hævder, at den jo overhovedet ikke er mulig i nogen Personlighed%
% footnote 1 on p. 135
\footnote[1]{He will come to seem a baffling paradox, until it is understood
that personality means a synthesis of elements too manifold and diverse ever to
be completely harmonized (p. 238).}. Men Spørgsmaalet er, om der virkelig
frembringes et nyt Per%
% --- p. 136 ---
\setcounter{footnote}{0}%
\opage{136}sonlighedsbillede gennem Kombination af de mange forskellige Tiders
og Kulturers Idealer. Og selv om det var muligt, vilde det nye Forbillede dog
blive saa væsentlig forskelligt fra det, der gennem Aarhundreder har virket paa
Menneskers Sind, at Kontinuiteten med den historiske Skikkelse, der stadig
skulde lægge Navn til, vil blive brudt%
% footnote 1 on p. 136
\footnote[1]{\textsc{Stanley Hall's} Forgængere i religionspsykologisk Forsken
indenfor amerikansk Literatur, især \textsc{Starbuck, Coe} og \textsc{Leuba}
fremhævede stærkt Ensidigheden i den traditionelle Udformning af Jesusbilledet.
Og navnlig \textsc{Coe} lægger ogsaa megen Vægt paa det betænkelige ved at anse
en enkelt Type for ene berettiget. Se min Afhandling \textit{Om nogle
religionsfilosofiske Arbejder fra den nyeste Tid.} 1903. (Mindre Arbejder.
Anden Række).}.

Hverken i Forbilleder eller i Dogmer kan Kontinuiteten i det menneskelige
Aandsliv søges, men i den dybe Trang, der i Vexelforhold med Livets Gang, dets
Gaver og dets Opgaver, dets Højder og dets Lavninger, afføder de skiftende
Symboler, mellem hvilke der derfor kan finde Analogi, men ikke Identitet Sted.
Analogien gør det muligt at lære af andre Forbilleder og andre Dogmer end vore
egne. Men Identificering er umulig, da den strider mod Livets Natur og Kaar.

Jo mere Religion bliver en Livspoesi uden dogmatisk Tendens, des mindre vil den
paastaa, at de Værdikilder, den haaber vil vedblive at rinde, selv om de møder
nok saa megen Modstand, tillige skulde være de Kilder, af hvilke alle i vor
Erfaring foreliggende Emner udspringer. Det vil være den nok, at der stadig
gaar en Værdistrøm gennem Verden, og den vil vække en Stræben, ikke blot efter
at faa Del i, hvad denne Strøm fører med sig, men ogsaa efter at bortrydde
Hindringerne for dens Løb, saavidt det staar i Menneskers Magt. ---

% --- p. 137 ---
\setcounter{footnote}{0}%
\opage{137}I min Religionsfilosofi findes disse Tanker nærmere udviklede.
Hensigten med dette Slutningsparti af Afhandlingen har kun været at antyde, at
de emotionelle Analogier, der gør sig gældende i praktisk og religiøs Symbolik,
er ejendommelige Sidestykker, om end tildels Modstykker, til de rationelle
Analogier, der er af saa stor Betydning i Videnskaben. Begge Arter af Analogi
staar mere og mere i Modsætning til den Analogi, der ligger skjult i den
primitive »Participation«.

\begin{center}\rule{2cm}{0.4pt}\end{center}

% --- p. 138 --- INDHOLD: not in the e-book; the scan, read at the image.
% Chapter heads letterspaced in the print (\emph); A.--F. not.
\clearpage
\opage{138}
\begin{center}{\large INDHOLD}\end{center}
\medskip
\hfill Side\par
\tocline{}{\emph{Indledning}}{3}
\tocline{I.}{\emph{Uvilkaarlige Analogier}}{7}
\tocline{II.}{\emph{Analogi og Logik}}{37}
\tocline{III.}{\emph{Analogi mellem Erkendelsesfunktioner}}{61}
\tocline{IV.}{\emph{Analogi mellem Erkendelsesomraader}}{79}
\tocline{}{\quad A. Logik og Aritmetik}{80}
\tocline{}{\quad B. Aritmetik og Geometri}{85}
\tocline{}{\quad C. Formal og real Videnskab}{89}
\tocline{}{\quad D. Naturvidenskab og Aandsvidenskab}{109}
\tocline{}{\quad E. Etik og teoretisk Videnskab}{119}
\tocline{}{\quad F. Verdensanskuelse og Videnskab}{124}
\tocline{V.}{\emph{Poetisk og religiøs Symbolik}}{129}
\vfill
\begin{center}\footnotesize
Forelagt paa Mødet den 23. Februar 1923.\\
Færdig fra Trykkeriet den 29. Juni 1923.
\end{center}

\end{document}
